\documentclass[11pt,a4paper]{article}
\usepackage[top=2.5cm,bottom=2.8cm,left=2.5cm,right=2.5cm]{geometry}
\usepackage{amsmath,amssymb,amsthm}
\usepackage{hyperref,xcolor,enumitem,booktabs}
\usepackage{graphicx}
\usepackage{subcaption}

\hypersetup{colorlinks=true,linkcolor={blue!60!black},
  citecolor={blue!60!black},urlcolor={blue!60!black}}

\newtheorem{theorem}{Theorem}[section]
\newtheorem{lemma}[theorem]{Lemma}
\newtheorem{proposition}[theorem]{Proposition}
\newtheorem{corollary}[theorem]{Corollary}
\theoremstyle{definition}\newtheorem{definition}[theorem]{Definition}
\theoremstyle{remark}
\newtheorem{remark}[theorem]{Remark}
\newtheorem{conjecture}[theorem]{Conjecture}
\theoremstyle{definition}\newtheorem{construction}[theorem]{Construction}

\newcommand{\vp}{\varphi}
\newcommand{\lam}{\lambda}
\newcommand{\Ja}{J_{2a}}
\newcommand{\Jfour}{J_4}
\newcommand{\ms}{\mu^*}
\newcommand{\TCMB}{T_{\mathrm{CMB}}}
\newcommand{\Lcond}{\Lambda_{\mathrm{cond}}}
\newcommand{\rinfty}{\rho_{\infty}}
\newcommand{\LUV}{\Lambda_{\mathrm{UV}}}
\newcommand{\rCMB}{\rho_{\mathrm{CMB}}}
\newcommand{\SU}{\mathrm{SU}}
\newcommand{\Efb}{E_{\mathrm{fb}}}
\newcommand{\Kfb}{K_{\mathrm{fb}}}
\newcommand{\twoI}{2I}
\newcommand{\twoT}{2T}
\newcommand{\twoO}{2O}
\newcommand{\QH}{Q_H}
\newcommand{\vEW}{v_{\mathrm{EW}}}
\newcommand{\deff}{d_{\mathrm{eff}}}
\newcommand{\Jiso}{J_{2\mathrm{iso}}}

\title{\textbf{Quark Generation Masses in the Density-Feedback Hopfion:\\
A Survey of Ruled-Out Mechanisms and the $Q_H=3$ Gradient-Flow Landscape}}
\author{Frederick Manfredi}
\date{2026}

\begin{document}
\maketitle

\begin{abstract}
Paper~XV~\cite{Paper15} establishes the $\QH=3$ trefoil sector's
qualitative quark structure (colour, charge, confinement) from the
McKay chain $\twoT\leftrightarrow E_6\leftrightarrow\SU(3)_1$, but
leaves the quark generation mass hierarchy as an open problem: the
inter-generation T-matrix phases contribute only an $O(1)$ correction
($m_2/m_1|_T\approx0.06$), so the dominant hierarchy --- six orders of
magnitude from $m_u$ to $m_t$ --- must come from a $\vp$-tower level
$n_q$ that is not fixed by anything in Paper~XV. This paper documents
a systematic search for that mechanism and reports purely negative
results. Six candidate mechanisms are tested and ruled out, each for
a specific, identifiable structural reason rather than a vague failure
to match data: (1) a literal reading of the $96$ $\vp$-containing
vertices of the $600$-cell broken under $\twoI\to\twoT$ fails on three
independent grounds (no normal-subgroup structure, wrong order of
magnitude, and a category mismatch between combinatorial and
dynamical quantities); (2) the three $1$-dimensional irreducible
representations of $\twoT$ do not supply a generation label independent
of colour, since both are permuted by the \emph{same} outer
automorphism; (3) varying the $\SU(3)$ WZW level, the direct analogue
of the lepton sector's level-varying generation ladder, is
structurally unavailable, since the central-charge match
$c_{(E_6)_k}=3\,c_{\SU(3)_k}$ that grounds the entire $\QH=3$
embedding holds only at $k=1$; (4) the quotient $\twoT/Q_8$ by the
normal Sylow-$2$ subgroup $Q_8\cong\mathrm{DiC}_2$ is a genuine
$\mathbb{Z}_3$ independent of the colour automorphism, but its three
cosets split asymmetrically (one distinguished, two related by
conjugation) and cannot support an ordered ladder; (5) Paper~V's
proved colour level shift $\delta n^{(C)}=\tfrac12$~\cite{Paper5} is a
genuine, generation-universal result, but the residual gap between its
prediction and measured quark masses is not a smooth function of
generation index, ruling out any simple per-generation correction
layered on top of it; (6) Paper~V's alternative integer-exact
quark-to-lepton ratio formula, built on Paper~VI's proved
incommensurability of the WZW and tower ``clocks''
($\pi/\ln\vp\notin\mathbb{Q}$, from the algebraic independence of
$\pi$ and $e^\pi$), rests on only one independently derived input (the
strange-quark $E_8$ Coxeter coincidence); the remaining five quark
assignments follow from that one input together with the
already-known mass ordering via a monotonicity argument, and reproduce
measured masses with errors of $26\%$ to $70\%$, comparable to or
worse than mechanism~(5). We present each test with its explicit
numerical or group-theoretic falsification, to narrow the search space
for subsequent work. A seventh, qualitatively different candidate,
motivated directly by colour confinement, is also examined: the
trefoil's Y-junction string-tension energy $\sigma\cdot3R_0$, present
only for $\QH=3$ and absent for the unconfined $\QH=2$ torus. Rather
than falsifying this candidate, we show that the chameleon-sector
scalar field and the Hopf profile function are the same field under
an explicit map (resolving one possible obstruction), but that the
Y-junction arms lie off the trefoil's own tube, so $\sigma$ cannot be
read off from existing tube integrals. Both $\sigma\cdot3R_0$ and
$E_\cup$ are nevertheless determined: $\sigma\cdot3R_0 =
0.04405\,E_{\mathrm{profile}}^*$ analytically
(Proposition~\ref{P16:prop:sigma_ext_formula}), and $E_\cup/E_{\mathrm{baryon}}
= 6.9\times10^{-8}$ numerically via gradient flow in the junction domain
(Proposition~\ref{P16:prop:ejunc_zero}), ruling out $E_\cup$ as a source
of generation-dependence. The baryon energy reduces to $E_{\mathrm{baryon}}
= 1.04405\,E_{\mathrm{profile}}^*$ exactly. We revisit
the scale ambiguity of Remark~\ref{P16:rem:scale_ambiguity} with a proper
one-loop, multi-threshold QCD running of the colour-shift predictions
from $\LUV$ to the PDG comparison scale: the strange-quark prediction
matches measurement to $0.09\%$, while the other five remain off by
factors of $1.7$ to $49$, with a systematic, generation-growing,
always-undershooting pattern consistent with (but not yet derived
from) a missing confinement-related energy contribution. We close by
asking whether the pole mass's well-known renormalon ambiguity
corresponds to a geometric stability feature of the $\QH=3$
condensate, and find that this question cannot yet be posed precisely:
Paper~XV's own topology-preserving saddle-finder problem is open, and
its documented failure mode --- a three-lobed instability from
self-interaction between the trefoil's three crossing regions, absent
for the unconfined $\QH=2$ torus --- is a structural, non-speculative
analogue of the renormalon obstruction, making that saddle-finder
problem a literal prerequisite for the question rather than separate
background. Investigating that prerequisite directly, we identify a
structural defect in the existing numerical solver's initial
condition (a forced phase discontinuity at each crossing, with
divergent gradient energy under mesh refinement), diagnose why a
natural candidate repair fails to carry the required topology, and
construct a corrected ansatz --- supplying each strand's own poloidal
winding before any crossing-region blending --- that is verified
smooth at all three crossings, has bounded gradient energy, and
carries the confirmed Hopf charge $\QH=3$. Tested directly against the
solver's own energy functional, the corrected construction's larger
apparent cost at the solver's coarse working resolution is traced
precisely to an under-resolved poloidal winding uniform along the
tube, not to a hidden divergence: local refinement near a crossing
converges cleanly, in sharp contrast to the original ansatz's
confirmed divergence under the same refinement. Direct gradient-flow
minimisation from the corrected construction reveals a clean two-phase
landscape: a Phase~1 in which excess gradient energy is shed efficiently
(the $K$-minimum near step~710 marks its end), followed by a Phase~2
in which the discrete grid's topological floor proves insufficient to
prevent slow cumulative drift toward the vacuum. The $K$-minimum field
at the $N=144$ run is $K=2313.9$; a refined $N=192$, $h=0.05$ run
($\approx16$ grid points across the tube diameter) gives the
resolution-convergent value $K_{\min}\approx2294$.
Preliminary data from a higher-resolution $N=256$, $h=0.0375$ run
(reported fully in Paper~XVIII) gives $K_{\min}\approx2268$, consistent
with continued $O(h^2)$ convergence toward the continuum value.
\end{abstract}

\tableofcontents
\bigskip\newpage

%══════════════════════════════════════════════════════════════════════════════
\section{Introduction}
\label{P16:sec:intro}
%══════════════════════════════════════════════════════════════════════════════

Paper~XV~\cite{Paper15} proves that the $\QH=3$ topological sector of
the density-feedback Faddeev--Niemi Hopfion is the trefoil knot
$T_{2,3}$, with WZW model $\SU(3)_1$ reached via the McKay chain
$\twoT\leftrightarrow E_6\leftrightarrow(E_6)_1\supset\SU(3)_1^{\times3}$.
This identifies the quark sector's colour structure, fractional
electric charges $|Q_d|=\tfrac13$, $|Q_u|=\tfrac23$, and topological
confinement mechanism, all as proved consequences of representation
theory and knot geometry. What it does not supply is the quark mass
spectrum itself.

\paragraph{The shape of the open problem.}
Paper~XV's own open problem~(iv) decomposes the quark-mass question
into three parts. Part~(a), the generation-to-irrep assignment, and
part~(b), the inter-generation T-matrix phases
$T_g^{\mathrm{quark}}=T_{j_g}^{\SU(2)_3}\in\{-\tfrac{3}{40},\tfrac{3}{40},
\tfrac{13}{40}\}$, are resolved: they reproduce only an $O(1)$
correction to the mass ratios ($m_2/m_1|_T\approx0.06$ from T-phases
alone). The observed hierarchy, $m_c/m_u\approx580$ and
$m_t/m_u\approx8\times10^4$, is far too large to come from this
mechanism. By direct analogy with the lepton sector, where the
dominant hierarchy comes from the discrete $\vp$-tower level $n_e=20$
(Paper~XII~\cite{Paper12}) rather than from T-matrix phases, the
quark hierarchy should likewise come from a tower level $n_q$ --- but
no candidate computation of $n_q$ from first principles exists in
Paper~XV.

\paragraph{Purpose of this paper.}
We report a systematic attempt to locate the mechanism that fixes
$n_q$, conducted by testing specific, falsifiable hypotheses against
both the group theory of $\twoT\subset\twoI$ and the numerical quark
mass spectrum. Every hypothesis tested in
Sections~\ref{P16:sec:vertices}--\ref{P16:sec:e8coxeter} failed. We document
each failure with its precise cause, since a documented negative
result narrows the search space for future attempts in a way that
silence does not. Section~\ref{P16:sec:vertices} treats the
$96$-broken-vertex hypothesis; Section~\ref{P16:sec:colourauto} the
$1$-dimensional-irrep hypothesis;
Section~\ref{P16:sec:colour_exclusion} proves the complementary positive
result that the colour $\mathbb{Z}_3$ has no action on the $\QH=2$
sector (Theorem~\ref{P16:thm:colour_exclusion}); Section~\ref{P16:sec:level} the
$\SU(3)_k$ level-variation hypothesis; Section~\ref{P16:sec:q8} the $Q_8$-coset
hypothesis; Section~\ref{P16:sec:colourshift} the generation-dependence of
Paper~V's colour level shift; Section~\ref{P16:sec:e8coxeter} the $E_8$
Coxeter-exponent assignment built on Paper~VI's incommensurability
principle. Section~\ref{P16:sec:stringtension} examines a seventh,
qualitatively different candidate motivated by colour confinement
itself --- the string-tension energy of the trefoil's Y-junction ---
and, rather than falsifying it, identifies precisely what is required
to make it computable. Section~\ref{P16:sec:rgrunning} returns to the
scale ambiguity of Section~\ref{P16:sec:colourshift} with a proper
multi-threshold RG running of the colour-shift predictions, isolating
a striking sub-percent match for the strange quark and a systematic
residual pattern in the other five. Section~\ref{P16:sec:saddle_prerequisite}
asks whether the pole mass's renormalon ambiguity corresponds to a
geometric feature of the $\QH=3$ condensate, and finds instead that
Paper~XV's own open saddle-finder problem is a literal prerequisite
for posing that question precisely; it also identifies a second,
independently verifiable candidate cause of the saddle-finder's
non-convergence, a forced $\pi$ phase discontinuity in the field
ansatz at every crossing region
(Section~\ref{P16:sec:phase_discontinuity}), distinct from Paper~XV's
inter-tube overlap mechanism, and a local repair --- a genuine $S^3$
(complex-doublet) superposition ansatz --- that removes the
divergence numerically (Section~\ref{P16:sec:s3_repair}), but is shown
not to carry the required $\QH=3$ topology when checked globally
(Section~\ref{P16:sec:s3lift_charge_zero}): the construction inherits a
missing poloidal winding from the underlying single-nearest-point
ansatz, giving $\QH=0$. Section~\ref{P16:sec:open} restates
the open problem in light of all nine investigations.

%══════════════════════════════════════════════════════════════════════════════
\section{The 96 Broken Vertices of the 600-Cell}
\label{P16:sec:vertices}
%══════════════════════════════════════════════════════════════════════════════

The $600$-cell, whose $120$ vertices (as unit quaternions) form the
binary icosahedral group $\twoI$, decomposes under the standard
$24$-cell subgroup $\twoT\subset\twoI$ as
\begin{equation}
  \twoI \;=\; \twoT \,\sqcup\, (\twoI\setminus\twoT),
  \qquad |\twoT|=24,\quad |\twoI\setminus\twoT|=96.
  \label{P16:eq:vertex_decomp}
\end{equation}
The $96$ vertices of $\twoI\setminus\twoT$ all have coordinates
involving $\vp$ (specifically $\vp/2$ and $1/(2\vp)$), while the $24$
vertices of $\twoT$ are rational. This decomposition is exact and
elementary (Paper~XV, Proposition~9.2); it is not in dispute. What
is tested here is a specific hypothesis built on it.

\begin{conjecture}[The broken-vertex hypothesis, as tested]
\label{P16:conj:vertices}
The quark tower level $n_q$ is determined by some combinatorial
invariant of the $96$ vertices broken under $\twoI\to\twoT$, read
either as a coset-counting argument on $\twoI/\twoT$, or by
extracting the $\vp$-exponent appearing in the broken vertices'
coordinates.
\end{conjecture}

We tested three precise readings of Conjecture~\ref{P16:conj:vertices} and
falsified each.

\begin{proposition}[Falsification of the broken-vertex hypothesis]
\label{P16:prop:vertices_fail}
Conjecture~\ref{P16:conj:vertices} fails under all three of the following
readings.
\begin{enumerate}[label=(\alph*)]
\item \textbf{Counting argument.} No combination of the integers
  $96$, $24$, $|\twoI{:}\twoT|=5$, $N_2=10$, $N_3=3$ reproduces the
  order of magnitude of any lepton/quark mass ratio. Solving
  $m_u/m_e=\vp^{2\Delta n}$ for the required level shift $\Delta n$
  gives $\Delta n\approx1.52$ --- not an integer, and not close to
  any simple function of the available integers.
\item \textbf{Coordinate $\vp$-exponent.} The $96$ broken vertices
  contain only $\vp^{\pm1}$ in their coordinates; no vertex coordinate
  involves a higher power of $\vp$. The lepton tower level $n_e=20$ is
  not a property of any finite group's vertex geometry: it is fixed by
  the transcendental thermodynamic matching condition
  $\Efb/V=\rCMB$ (Paper~XII~\cite{Paper12}), a dynamical
  renormalisation-group quantity with no counterpart in the static
  $600$-cell coordinates.
\item \textbf{Coset-space structure.} $\twoT$ is not a normal subgroup
  of $\twoI$ (direct computation: $g^{-1}\twoT g\neq\twoT$ for $g$ in
  four of the five conjugates). Consequently $\twoI/\twoT$ carries no
  group structure --- it is five unstructured cosets with no
  composition law --- and there is no algebraic invariant to extract
  an integer from beyond the bare cardinalities $5$ and $96$, both of
  which are already accounted for in reading~(a).
\end{enumerate}
\end{proposition}

\begin{proof}
Reading~(a) is verified by direct substitution, using the measured
$m_e=0.511$~MeV and $m_u=2.16$~MeV: $\Delta n=
\log(m_u/m_e)/\log\vp^2=1.52$. Reading~(b) follows from the explicit
vertex list (Paper~XV, Proposition~9.2): the $96$ vertices are
$\tfrac12(\pm\vp,\pm1,\pm1/\vp,0)$ and even permutations, in which
$\vp$ appears only to the first power. Reading~(c) is verified by
explicit quaternion conjugation: for the five conjugates
$g_i^{-1}\twoT g_i$, $i=1,\dots,5$, of the standard $24$-cell subgroup
inside $\twoI$, only one (the standard copy itself) is fixed under
conjugation by all $120$ elements of $\twoI$; the other four are
permuted among themselves, confirming $\twoT$ is self-normalizing but
not normal.
\qed
\end{proof}

\begin{remark}[What survives]
The geometric decomposition~\eqref{P16:eq:vertex_decomp} itself, and the
related fact that $\twoT$ is the unique $\vp$-free conjugate among the
five copies of $\twoT$ inside $\twoI$ (Section~\ref{P16:sec:5copies}),
remain valid and are not affected by Proposition~\ref{P16:prop:vertices_fail}.
What fails is specifically the inference from this geometry to a
quark mass formula; the geometry itself is sound and may yet prove
useful as corroborating intuition for the $\vp$-free character of QCD
(Paper~XV, Remark~9.4, after Paper~XV).
\end{remark}

%──────────────────────────────────────────────────────────────────────────────
\subsection{The five conjugate copies of $\twoT$ in $\twoI$}
\label{P16:sec:5copies}
%──────────────────────────────────────────────────────────────────────────────

Since the five-copy structure of $\twoT\subset\twoI$ recurs in
Sections~\ref{P16:sec:vertices}--\ref{P16:sec:q8}, we record its properties
here for reference.

\begin{proposition}[Structure of the five conjugate copies]
\label{P16:prop:five_copies}
$\twoI$ contains exactly five conjugate copies of $\twoT$, matching
$[\twoI{:}\twoT]=5$. Of these:
\begin{enumerate}[label=(\roman*)]
\item all five share the central elements $\{\pm1\}$;
\item exactly one is $\vp$-free (the standard $24$-cell, with
  imaginary units $\{\pm i,\pm j,\pm k\}$); the other four each
  contain their own distinct sextuple of $\vp$-valued imaginary units,
  with no imaginary unit shared between any two of the five copies;
\item every pair of the five copies intersects in exactly six
  elements: $\{\pm1\}$ together with four order-$6$ elements specific
  to that pair;
\item the simultaneous intersection of all five copies is exactly
  $\{\pm1\}$.
\end{enumerate}
\end{proposition}

\begin{proof}
Direct computation: represent $\twoI$ as $120$ unit quaternions (the
$8$ Lipschitz units, $16$ Hurwitz units, and $96$ $\vp$-valued units
of Paper~XV, Proposition~9.2), conjugate the standard $24$-cell
subgroup $\twoT_0=\{\text{Lipschitz}\}\cup\{\text{Hurwitz}\}$ by each
of the $120$ elements, and deduplicate the resulting subgroups. Exactly
five distinct subgroups result. Properties~(i)--(iv) are verified by
direct set intersection on the explicit element lists.
\qed
\end{proof}

\begin{remark}[Double role of this proposition]
Proposition~\ref{P16:prop:five_copies} serves two purposes.  Its primary
purpose here is to record the five-copy structure that recurs across
Sections~\ref{P16:sec:vertices}--\ref{P16:sec:q8}.  Its secondary purpose,
used in Section~\ref{P16:sec:colour_exclusion}, is to establish that
$N_{\twoI}(\twoT)=\twoT$ --- this is Fact~3 in the proof of
Theorem~\ref{P16:thm:colour_exclusion}, which closes off the inner-automorphism
branch of that proof's case analysis.
\end{remark}

%══════════════════════════════════════════════════════════════════════════════
\section{The 1-Dimensional Irreducible Representations of $\twoT$}
\label{P16:sec:colourauto}
%══════════════════════════════════════════════════════════════════════════════

The character table of $\twoT$ has seven irreducible representations:
three $1$-dimensional ($\Gamma_1,\Gamma_\omega,\Gamma_{\omega^2}$,
bosonic), three $2$-dimensional ($\Gamma_2,\Gamma_{2\omega},
\Gamma_{2\omega^2}$, fermionic, identified with the three quark
colours in Paper~XV), and one $3$-dimensional
($\Gamma_3$). The colour assignment uses an order-$3$ outer
automorphism of $\twoT$ that permutes
$\Gamma_2\to\Gamma_{2\omega}\to\Gamma_{2\omega^2}$.

\begin{remark}[Double role of this section's result]
Proposition~\ref{P16:prop:bosonic_fail} below serves two purposes.
Its primary purpose here is to falsify the bosonic-triple generation
hypothesis.  Its secondary purpose, used in
Section~\ref{P16:sec:colour_exclusion}, is to establish that $\sigma$ is a
\emph{strictly outer} automorphism of $\twoT$ --- this is Fact~1 in
the proof of Theorem~\ref{P16:thm:colour_exclusion}.
\end{remark}

\begin{conjecture}[The bosonic-triple hypothesis, as tested]
\label{P16:conj:bosonic}
The three $1$-dimensional irreducible representations
$\{\Gamma_1,\Gamma_\omega,\Gamma_{\omega^2}\}$, being permuted by an
order-$3$ symmetry in the same way as the colour triple, supply an
independent generation label.
\end{conjecture}

\begin{proposition}[Falsification of the bosonic-triple hypothesis]
\label{P16:prop:bosonic_fail}
Conjecture~\ref{P16:conj:bosonic} fails: the symmetry permuting
$\{\Gamma_1,\Gamma_\omega,\Gamma_{\omega^2}\}$ is the \emph{same}
single outer automorphism that permutes the colour triple, not an
independent one.
\end{proposition}

\begin{proof}
An outer automorphism of a finite group $G$ acts on the entire set of
irreducible characters $\mathrm{Irr}(G)$ simultaneously, by
precomposition. $\twoT$ has a unique order-$3$ outer automorphism (up
to inverse), coming from the triality of the $\widehat{E}_6$ McKay
diagram. This single automorphism necessarily permutes
\emph{both} three-element orbits in $\mathrm{Irr}(\twoT)$ at once: it
fixes $\Gamma_3$ (the unique self-paired irrep) and acts as a
$3$-cycle on each of $\{\Gamma_1,\Gamma_\omega,\Gamma_{\omega^2}\}$
and $\{\Gamma_2,\Gamma_{2\omega},\Gamma_{2\omega^2}\}$. A quark
labelled by both a colour irrep $\Gamma_{2\omega^c}$ and a bosonic
irrep $\Gamma_{\omega^{c'}}$ under the same automorphism has $c$ and
$c'$ correlated, not independent: at most one free $\mathbb{Z}_3$
label is available this way, and it is the colour label already in
use.
\qed
\end{proof}

%══════════════════════════════════════════════════════════════════════════════
\section{Varying the $\SU(3)$ WZW Level}
\label{P16:sec:level}
%══════════════════════════════════════════════════════════════════════════════

The lepton sector's three generations correspond to three distinct
$\SU(2)$ WZW levels, $k_g=1,2,3$ (Paper~I~\cite{Paper1},
the Yukawa ladder result), with T-matrix phases
\begin{equation}
  T_g \;=\; \frac{6-k_g}{8(k_g+2)},
  \qquad
  T_1=\tfrac{5}{24},\ T_2=\tfrac18,\ T_3=\tfrac{3}{40},
  \label{P16:eq:lepton_Tg}
\end{equation}
satisfying the geometric ratio $T_1/T_2=T_2/T_3=5/3=(k+2)/k$ at
$k=3$. The direct quark-sector analogue would vary the $\SU(3)$ WZW
level $k$ in place of $\SU(2)$'s.

\begin{conjecture}[The level-variation hypothesis, as tested]
\label{P16:conj:level}
The quark generation index corresponds to varying the level $k$ of
$\SU(3)_k$, by analogy with the lepton sector's level-varying ladder.
\end{conjecture}

\begin{proposition}[Falsification of the level-variation hypothesis]
\label{P16:prop:level_fail}
Conjecture~\ref{P16:conj:level} fails: the McKay embedding that produces
the $\QH=3$ sector exists only at $k=1$.
\end{proposition}

\begin{proof}
The chain $\twoT\leftrightarrow E_6\leftrightarrow(E_6)_1\supset
\SU(3)_1^{\times3}$ (Paper~XV, Section~2) relies on the exact
central-charge match
\begin{equation}
  c_{(E_6)_k} \;=\; 3\,c_{\SU(3)_k},
  \qquad
  c_{(E_6)_k}=\frac{78k}{k+12},
  \quad
  c_{\SU(3)_k}=\frac{8k}{k+3}.
  \label{P16:eq:central_charge_match}
\end{equation}
At $k=1$: $c_{(E_6)_1}=6$ and $3\,c_{\SU(3)_1}=6$, an exact match. At
$k=2$: $c_{(E_6)_2}=78/7\approx11.14$ versus $3c_{\SU(3)_2}=48/5=9.60$;
the two diverge further for every $k\geq2$ checked explicitly up to
$k=5$. The conformal embedding that grounds the entire $\QH=3$ sector
therefore exists only at $k=1$; there is no $k$-ladder to vary.
\qed
\end{proof}

\begin{remark}[A second, independent reason]
\label{P16:rem:no_target}
Even granting a hypothetical $\SU(3)_k$ ladder, the lepton-sector
ladder it would generalise is not itself a McKay-correspondence
object: the levels $k_g=1,2,3$ in equation~\eqref{P16:eq:lepton_Tg} are a
free WZW parameter scan with no finite-subgroup interpretation
attached to $k_g=1$ or $k_g=2$ individually. Only $k=3$, the level
used for the $\QH=2$ Hopfion proper, is tied to $\twoI$ via the
pentagon mechanism (Paper~III). Proposition~\ref{P16:prop:level_fail}
therefore rules out an analogy to a mechanism that was not itself
grounded in McKay theory to begin with.
\end{remark}

%══════════════════════════════════════════════════════════════════════════════
\section{The Quotient $\twoT/Q_8$}
\label{P16:sec:q8}
%══════════════════════════════════════════════════════════════════════════════

The binary tetrahedral group $\twoT$ has a normal Sylow-$2$ subgroup
$Q_8\cong\mathrm{DiC}_2=\{\pm1,\pm i,\pm j,\pm k\}$ (the eight Lipschitz
units of Paper~XV, Proposition~9.2), of order $8$, with quotient
$\twoT/Q_8\cong\mathbb{Z}_3$ (the unique group of order $3$).

\begin{conjecture}[The $Q_8$-coset hypothesis, as tested]
\label{P16:conj:q8}
The three generations correspond to the three cosets of $Q_8$ in
$\twoT$, which form a $\mathbb{Z}_3$ structure independent of the
colour automorphism (since $Q_8$-coset multiplication is an
\emph{inner} operation, generated by an actual element of $\twoT$,
whereas the colour $\mathbb{Z}_3$ of Section~\ref{P16:sec:colourauto} is
an \emph{outer} automorphism).
\end{conjecture}

The independence claimed in Conjecture~\ref{P16:conj:q8} is correct, and
is the one respect in which this hypothesis differs structurally from
Section~\ref{P16:sec:colourauto}'s failure.

\begin{proposition}[Independence of the $Q_8$-coset $\mathbb{Z}_3$]
\label{P16:prop:q8_independent}
The $\mathbb{Z}_3$ structure on the cosets of $Q_8$ in $\twoT$ is
independent of the colour automorphism of Section~\ref{P16:sec:colourauto}.
\end{proposition}
\begin{proof}
Characters are conjugation-invariant: $\chi(g^{-1}xg)=\chi(x)$ for any
representation. Inner automorphisms (conjugation by an element of
$\twoT$) therefore act trivially on $\mathrm{Irr}(\twoT)$ as a set,
while the $Q_8$-coset structure is generated by left multiplication by
an order-$3$ element of $\twoT$ --- an inner operation by definition.
The colour automorphism of Proposition~\ref{P16:prop:bosonic_fail} is
outer. An inner and an outer automorphism generate independent
structures except in the degenerate case where the outer automorphism
happens to coincide with some inner one, which does not occur here
since $\twoT/Z(\twoT)\cong A_4$ has trivial centre-quotient
intersection with its own outer automorphism group of order $3$
beyond the identity.
\qed
\end{proof}

Despite this genuine independence, the hypothesis fails on a separate
ground: the three cosets do not have the right internal structure to
support an ordered generation ladder.

\begin{proposition}[Falsification of the $Q_8$-coset hypothesis]
\label{P16:prop:q8_fail}
The three cosets of $Q_8$ in $\twoT$ split asymmetrically as one
distinguished coset ($Q_8$ itself) and two cosets related by a
residual conjugation symmetry, not as three mutually exchangeable
slots.
\end{proposition}

\begin{proof}
Direct computation of element orders within each coset: $Q_8$ itself
contains the identity, the central involution, and six order-$4$
elements, with order multiset $\{1,2,4,4,4,4,4,4\}$. The other two
cosets both have order multiset $\{3,3,3,3,6,6,6,6\}$ ---
indistinguishable from each other by this invariant, consistent with
their being related by the complex-conjugation symmetry
$\omega\leftrightarrow\omega^2$ that also relates $\Gamma_\omega$ and
$\Gamma_{\omega^2}$ in the character table. The lepton ladder
$T_1=\tfrac{5}{24},T_2=\tfrac18,T_3=\tfrac{3}{40}$, by contrast, is
three pairwise-distinct real numbers in strict geometric progression,
with no two related by any conjugation-like symmetry. The $Q_8$-coset
structure (``one special $+$ one conjugate pair'') and the lepton
ladder's structure (``three ordered, pairwise-inequivalent levels'')
are not the same shape, and no value of $Q_8$-coset membership
reproduces an ordered ladder without additional input not currently
available in the framework.
\qed
\end{proof}

%══════════════════════════════════════════════════════════════════════════════
\section{Generation-Dependence of the Colour Level Shift}
\label{P16:sec:colourshift}
%══════════════════════════════════════════════════════════════════════════════

Paper~V~\cite{Paper5} proves a universal colour level shift via
Witten's non-abelian bosonization~\cite{Witten1984}: a quark in the
fundamental of $\SU(3)$ is represented, in the bosonized picture, by a
complex free fermion field $\psi(z)$ of conformal weight
$h_\psi=\tfrac12$ universally (independent of $k$, $N$, or any other
parameter, by the standard OPE $\psi(z)\psi^\dagger(0)\sim z^{-1}$).

\begin{remark}[Architecture justified by Theorem~\ref{P16:thm:colour_exclusion}]
The formula derived in this section takes the form
$m_q^{(g)} = m_\ell^{(g)} \cdot [\text{colour correction}]$, starting
from the lepton mass $m_\ell^{(g)}$ and applying a multiplicative
$\QH=3$-sector factor.  Theorem~\ref{P16:thm:colour_exclusion}
(Section~\ref{P16:sec:colour_exclusion}) provides the group-theoretic
justification for why this is the only logically consistent
architecture: colour charge does not exist at the $\QH=2$ level, so
the lepton mass $m_\ell^{(g)}$ --- a pure $\QH=2$ quantity --- must
appear as the starting point, with the colour correction being the
first place any $\QH=3$-specific structure enters.  In retrospect,
the six mechanisms of Sections~\ref{P16:sec:vertices}--\ref{P16:sec:e8coxeter}
all assumed colour as a given and searched for an additional structure
within it; Theorem~\ref{P16:thm:colour_exclusion} shows they were
necessarily looking in the wrong place.
\end{remark}

This is identified with the tower level shift,
\begin{equation}
  \delta n^{(C)} \;=\; h_\psi \;=\; \frac12,
  \label{P16:eq:dn_shift}
\end{equation}
giving the quark mass formula, at the condensate scale $\LUV$,
\begin{equation}
  m_q^{(g)} \;=\; \vEW\,e^{-2\pi QT_g}\cdot\vp^{-2\delta n^{(C)}}
  \;=\; m_\ell^{(g)}\cdot\vp^{-1}
  \quad(\text{up-type}),
  \qquad
  m_{q,\mathrm{down}}^{(g)}
  \;=\; m_\ell^{(g)}\cdot\vp^{-3}
  \quad(\text{down-type}),
  \label{P16:eq:quark_mass_recap}
\end{equation}
where $m_\ell^{(g)}=\vEW e^{-2\pi QT_g}$ is the proved lepton mass
formula of Paper~I, with $Q=10$ and $T_g$ as in
equation~\eqref{P16:eq:lepton_Tg}. Equation~\eqref{P16:eq:quark_mass_recap} is
generation-universal: the same multiplicative factor $\vp^{-1}$
(up-type) or $\vp^{-3}$ (down-type) applies for $g=1,2,3$, and this
universality is exactly what makes Theorem~7.4 of Paper~V's quark
Koide relation $K_{\mathrm{up}}=K_{\mathrm{down}}=\tfrac23$ exact: the
Koide ratio is invariant under any common rescaling of three masses
(Paper~V, Lemma~8.1), so a generation-universal factor leaves
it unchanged.

\begin{conjecture}[Generation-dependent correction, as tested]
\label{P16:conj:colourshift}
The discrepancy between equation~\eqref{P16:eq:quark_mass_recap} and the
measured quark masses is a smooth, generation-dependent correction
$\delta n^{(C)}(g)$ that can be added to the universal shift
$\tfrac12$.
\end{conjecture}

\begin{proposition}[Falsification of the generation-dependent
  correction hypothesis]
\label{P16:prop:colourshift_fail}
Conjecture~\ref{P16:conj:colourshift} fails: the required correction is
not a smooth function of $g$, and is not even monotonic for down-type
quarks.
\end{proposition}

\begin{proof}
Define the required level shift $n_g$ for each measured quark by
$m_q^{(g)}/m_\ell^{(g)}=\vp^{-n_g}$, using PDG values for
$m_q$~\cite{PDG2024} and the measured charged-lepton masses for
$m_\ell^{(g)}$. Direct computation gives, for up-type,
\begin{equation}
  n_1=-3.00,\quad n_2=-5.17,\quad n_3=-9.51,
  \label{P16:eq:n_up_required}
\end{equation}
and for down-type,
\begin{equation}
  n_1=-4.60,\quad n_2=+0.26,\quad n_3=-1.78.
  \label{P16:eq:n_down_required}
\end{equation}
Equation~\eqref{P16:eq:quark_mass_recap} predicts $n_g=1$ (up-type) and
$n_g=3$ (down-type) for all $g$; the actual required values deviate by
$4$ to $10.5$ units of $n$ (up-type) and by amounts that change sign
between generations (down-type). The down-type sequence
$\eqref{P16:eq:n_down_required}$ is non-monotonic in $g$ ($n_2>n_1$ but
$n_3<n_2$), which rules out any correction expressible as a smooth
function of $g$ alone, monotonic or otherwise. The up-type sequence
$\eqref{P16:eq:n_up_required}$ is at least monotonic; its successive
differences $n_2-n_1=-2.17$ and $n_3-n_2=-4.34$ have ratio $1.999$,
suspiciously close to $2$, but this ratio is not robust: substituting
the $\overline{\mathrm{MS}}$ top mass ($\approx162.5$~GeV) for the pole
mass ($172.57$~GeV) --- an ordinary, physically motivated scheme
choice --- shifts the ratio to $1.94$. A structural fact intrinsic to
this framework should not be sensitive to a routine mass-scheme
convention external to it; the near-integer ratio is judged
coincidental.
\qed
\end{proof}

\begin{remark}[What this does and does not rule out]
Proposition~\ref{P16:prop:colourshift_fail} does not impugn the derivation
of $\delta n^{(C)}=\tfrac12$ itself (equation~\eqref{P16:eq:dn_shift}),
which rests on a clean, parameter-independent fact about free-fermion
conformal weight. It rules out the narrower hypothesis that this
universal shift is the \emph{only} correction needed: some further,
generation-dependent piece of physics, not yet identified, must be
layered on top of it, and that piece is evidently not a simple
function of the generation index alone.
\end{remark}

\begin{remark}[Scale ambiguity in Paper~V]
\label{P16:rem:scale_ambiguity}
A separate issue, distinct from Proposition~\ref{P16:prop:colourshift_fail},
is worth recording. Equation~\eqref{P16:eq:quark_mass_recap} is explicitly
scoped to the condensate scale $\LUV\approx6.8\times10^{17}$~GeV
(Paper~V, proof of Theorem~7.4), whereas Paper~V's own CKM-angle
formulas (Section~5 of that paper) use PDG/$\overline{\mathrm{MS}}$
quark masses directly, with no running step specified between the two
scales. One-loop running of equation~\eqref{P16:eq:quark_mass_recap}'s
prediction from $\LUV$ to $2$~GeV, using Paper~V's own anomalous
dimension $\gamma_0=8$, does not close the gap to measured values and
in fact widens it, but one-loop perturbative running is not expected
to be reliable in the infrared regime relevant to light-quark masses.
Resolving this scale ambiguity is a prerequisite for any quantitative
test of equation~\eqref{P16:eq:quark_mass_recap}, independent of
Proposition~\ref{P16:prop:colourshift_fail}.
\end{remark}

%══════════════════════════════════════════════════════════════════════════════
\section{The $E_8$ Coxeter Exponent Assignment}
\label{P16:sec:e8coxeter}
%══════════════════════════════════════════════════════════════════════════════

A different and more developed mechanism than
equation~\eqref{P16:eq:quark_mass_recap} appears elsewhere in Paper~V,
motivated by the incommensurability result of Paper~VI~\cite{Paper6}
(Proposition~5.5 and Remark~5.6 therein): individual lepton tower
levels are irrational, by the transcendence of $\pi/\ln\vp$
(Nesterenko's theorem~\cite{Nesterenko1996}), but the
\emph{quark-to-lepton mass ratio} is claimed to be integer-exact at
the condensate scale,
\begin{equation}
  \frac{m_q}{m_{\ell(g)}} \;=\; \exp\!\left(\frac{n_q\pi}{9}\right),
  \qquad n_q\in\mathbb{Z},
  \label{P16:eq:e8_ratio_recap}
\end{equation}
because the $\ln\vp$-dependence common to numerator and denominator
cancels in the ratio, leaving only the pure WZW phase unit
$\pi/9=\pi Q/(k\,h(E_8))$. This is a structurally different claim from
equation~\eqref{P16:eq:quark_mass_recap}: it does not assert a
generation-universal $\vp$-power shift, but an integer multiple of a
fixed angular unit, with the integer $n_q$ read off from the $E_8$
Coxeter exponents $m\in\{1,7,11,13,17,19,23,29\}$ via
$T^{(m)}=(4m-7)/360$ and $n_q=180(T_g-T^{(m)})$.

We verified the arithmetic of equation~(67) of Paper~V~\cite{Paper5} exactly
(rational arithmetic, no rounding) and confirmed all six stated
integers $n_t=15,n_b=3,n_c=4,n_s=0,n_d=7,n_u=3$. The question tested
here is not the arithmetic, which is correct, but whether the
assignment of $E_8$ exponents $m$ to specific quarks is derived from
independent physics or fitted to the known mass ordering.

\begin{remark}[Algebraic structure of the $n_q$ integers and a cross-sector identity]
\label{P16:rem:nq_structure}
Before evaluating the formula as a mechanism, it is worth recording
the algebraic structure of the $n_q$ integers, which is not made
explicit in Paper~V but is a consequence of the arithmetic verified
above.

\textbf{Factorisation.}
Define the generation-dependent base value
\begin{equation}
  \mathrm{Base}_g \;=\; 2k\,h(E_8)\,T_g \;+\; \tfrac{7}{2}
  \;=\; 180\,T_g \;+\; \tfrac{7}{2},
  \label{P16:eq:nq_base}
\end{equation}
where $T_g$ are the lepton T-matrix phases~\eqref{P16:eq:lepton_Tg} and
$k=3$, $h(E_8)=30$ are the WZW level and $E_8$ Coxeter number.
Substituting the explicit values gives $\mathrm{Base}_1=41$,
$\mathrm{Base}_2=26$, $\mathrm{Base}_3=17$, all integers.
The formula $n_q = 180(T_g - T^{(m)}) = 180T_g - (4m-7)/2$ then
factors as
\begin{equation}
  n_q \;=\; \mathrm{Base}_g \;-\; 2m,
  \label{P16:eq:nq_factored}
\end{equation}
which is manifestly an integer for every $E_8$ Coxeter exponent $m$
and every generation $g$.  All six values are recovered exactly: for
up-type, $n_u = 41 - 2\cdot19 = 3$, $n_c = 26 - 2\cdot11 = 4$,
$n_t = 17 - 2\cdot1 = 15$; for down-type, $n_d = 41 - 2\cdot17 = 7$,
$n_s = 26 - 2\cdot13 = 0$, $n_b = 17 - 2\cdot7 = 3$.

\textbf{A cross-sector identity.}
The base values satisfy
\begin{equation}
  \mathrm{Base}_1 - \mathrm{Base}_3
  \;=\; 180(T_1 - T_3)
  \;=\; 2k\,h(E_8)\,(T_1 - T_3)
  \;=\; 24
  \;=\; |\twoT|.
  \label{P16:eq:nq_base_identity}
\end{equation}
This is an exact arithmetic identity: the total span of the lepton
T-phase ladder, scaled by $2k\,h(E_8)$, equals the order of the
binary tetrahedral group.  Expressed in terms of group-theoretic
integers alone, it reads
\begin{equation}
  T_1 - T_3
  \;=\; \frac{|\twoT|}{2k\,h(E_8)}
  \;=\; \frac{24}{180}
  \;=\; \frac{2}{15}
  \;=\; \frac{2}{k(k+2)},
  \label{P16:eq:T_span_identity}
\end{equation}
where $k(k+2) = 3\cdot5 = 15$ combines the WZW level $k=3$ with the
Pentagon number $k+2=5=[\twoI:\twoT]$.  (Direct verification:
$T_1 - T_3 = \tfrac{5}{24} - \tfrac{3}{40} = \tfrac{25}{120} -
\tfrac{9}{120} = \tfrac{16}{120} = \tfrac{2}{15}$, confirming
equation~\eqref{P16:eq:T_span_identity}.)
The identity is not a consequence of the $E_8$ Coxeter labelling; it
holds for the lepton T-phases of Paper~I independently of anything in
this section, and expresses a direct coupling between the
$\QH=2$ generation structure and the order of the $\QH=3$
group.

The intermediate base steps are $\mathrm{Base}_1 - \mathrm{Base}_2 = 15$ and
$\mathrm{Base}_2 - \mathrm{Base}_3 = 9$, equal to $180(T_1-T_2)$ and
$180(T_2-T_3)$ respectively.  These are unequal ($15 \neq 9$), which is
already implied by the geometric-progression property of the lepton
T-phases, $T_1/T_2 = T_2/T_3 = (k+2)/k = 5/3$.

\textbf{Structural explanation of the down-type non-monotonicity.}
The $n_q$ integers for down-type quarks, read by generation, are
$n_d=7$, $n_s=0$, $n_b=3$: non-monotonic, with $n_s<n_b<n_d$.
The factorisation~\eqref{P16:eq:nq_factored} explains why.  For up-type
quarks, the $E_8$ exponent decreases sharply with generation:
$m_u=19\to m_c=11\to m_t=1$, a fall of $18$ from generation~1 to
generation~3.  For down-type quarks, $m_d=17\to m_s=13\to m_b=7$,
a fall of only $10$.  Since $n_q = \mathrm{Base}_g - 2m$ and the base
decreases by $24$ across three generations while $2m$ decreases by
$36$ (up-type) or $20$ (down-type), the \emph{net} change
$\Delta n_q = -\Delta\mathrm{Base} + 2\Delta m$ is $+12$ (up-type,
generation~$1\to3$) or $-4$ (down-type).  The individual generation
steps for down-type further split as
$\Delta n_q(1\to2) = -15 + 2(17-13) = -7$ and
$\Delta n_q(2\to3) = -9 + 2(13-7) = +3$, with a sign change between
them.  This sign change is not accidental: it reflects the fact that
$m_d - m_u = 17-19 = -2$ (the down quark has a \emph{smaller}
exponent than its up-type partner in generation~1) while
$m_b - m_t = 7-1 = +6$ (the reverse holds in generation~3).  That
reversal is a structural property of how the six smallest $E_8$
Coxeter exponents $\{1,7,11,13,17,19\}$ distribute across the two
isospin sectors, and any account of the down-type non-monotonicity
must reproduce it.

\textbf{Summary.}
The $n_q$ integers of equation~(67) of Paper~V~\cite{Paper5} are not
featureless; they encode two pieces of structural information: the
cross-sector identity~\eqref{P16:eq:nq_base_identity} linking the lepton
T-phase span to $|\twoT|$, and the distribution of $E_8$ Coxeter
exponents between isospin sectors that produces the up/down asymmetry.
Neither piece is a derivation of $n_q^{(g)}$ from first principles
--- the exponent assignments still require the input evaluated in
Propositions~\ref{P16:prop:e8_unique}--\ref{P16:prop:e8_mixed} below --- but
they constitute precise algebraic constraints that any future
derivation must reproduce.
\end{remark}

\begin{proposition}[Uniqueness of the monotone assignment, as stated
  by Paper~V]
\label{P16:prop:e8_unique}
Given $m_s=13$ and the physical mass ordering
$m_t>m_b>m_c>m_s>m_d>m_u$, the monotone assignment of the remaining
five quarks to five of the seven remaining $E_8$ Coxeter exponents
$\{1,7,11,17,19,23,29\}$, using exactly five of the available seven
and leaving the largest two ($23,29$) unused, is
$t\to1,\,b\to7,\,c\to11,\,d\to17,\,u\to19$, and is unique under these
constraints.
\end{proposition}

\begin{proof}
As given in Paper~V (proof of Proposition~7.5): monotonicity
($T^{(m)}$ increasing in $m$, hence $m_q$ decreasing in $m$) and the
ordering $m_t>m_b>m_c>m_s$ force the three quarks heavier than
strange into the three exponents below $13$, in decreasing order of
mass; $m_s>m_d>m_u$ forces the two quarks lighter than strange into
the (two smallest of the) exponents above $13$. With three exponents
below $13$ (namely $1,7,11$) and at least two needed above
($17,19$), both placements are forced uniquely, leaving $23$ and
$29$ unoccupied.
\qed
\end{proof}

\begin{proposition}[The strange-quark exponent is derived; the
  remaining five are not]
\label{P16:prop:e8_mixed}
Of the six $E_8$-exponent assignments in equation~(65) of Paper~V~\cite{Paper5},
exactly one, $s\leftrightarrow m=13$, follows from an identity
independent of the observed quark mass ordering. The remaining five
follow from that one assignment together with the observed mass
ordering itself, via Proposition~\ref{P16:prop:e8_unique}'s monotonicity
argument, and are not independently derived.
\end{proposition}

\begin{proof}
The value $m=13$ is fixed by requiring the combined icosahedral
$\times\,\SU(3)_3$ T-matrix phase to satisfy $T_{\mathrm{total}}(s)=
\tfrac18=T_\mu$ (Paper~V, Proposition~7.4), a numerical coincidence
between two independently computable quantities that Theorem~7.9 of
Paper~V (Theorem~7.9, the McKay--$E_8$ Coxeter spectrum theorem) explains via the
standard Verlinde formula for ADE-extended chiral algebras
(Cappelli--Itzykson--Zuber~\cite{CIZ1987}): $h_m=m/(k\,h(E_8))$ is
forced to equal $13/90$ by the coincidence itself, and $m=13$ is then
recognised as the fourth Coxeter exponent of $E_8$. This is a genuine,
non-trivial check, independent of any a priori knowledge of which
quark is ``heavier'' than which.

The remaining five assignments follow from Proposition~\ref{P16:prop:e8_unique}:
given $m_s=13$, the physical mass ordering
$m_t>m_b>m_c>m_s>m_d>m_u$ --- an observational fact, not a derived
one --- together with monotonicity (justified only by the formula's
own structure, $dm_q/dm<0$) fixes the remaining five exponents with no
further independent input. No computation determines \emph{which}
exponent corresponds to \emph{which} quark beyond reproducing an
ordering already known from measurement.
\qed
\end{proof}

\begin{proposition}[Quoted accuracy of the five fitted assignments]
\label{P16:prop:e8_accuracy}
The five fitted assignments of Proposition~\ref{P16:prop:e8_mixed}
reproduce the measured quark masses with errors of $51\%$ ($b$),
$-70\%$ ($c$), $26\%$ ($d$), $-32\%$ ($u$), and an undefined/large
error for $t$ (Paper~V's own Table, Proposition~7.5), against the
$2.2\%$ accuracy of the independently-derived strange-quark case.
\end{proposition}

\begin{proof}
Direct comparison of Paper~V's stated predicted masses
$m_q^{E_8}=\vEW e^{-2\pi QT^{(m)}}$ against PDG values: $b$,
$6.3$~GeV predicted versus $4.18$~GeV measured ($51\%$); $c$,
$386$~MeV versus $1275$~MeV ($-70\%$); $d$, $5.9$~MeV versus
$4.7$~MeV ($26\%$); $u$, $1.5$~MeV versus $2.2$~MeV ($-32\%$); $t$,
undefined at tree level since $T^{(1)}<0$, per Paper~V's own footnote.
\qed
\end{proof}

\begin{remark}[Why this still counts as a negative result for $n_q$]
Propositions~\ref{P16:prop:e8_mixed}--\ref{P16:prop:e8_accuracy} together show
that equation~\eqref{P16:eq:e8_ratio_recap}, despite resting on a real and
well-motivated incommensurability argument (Paper~VI,
Proposition~5.5), does not supply a predictive mechanism for the
quark tower level once the one genuine input (the strange-quark
coincidence) is used up: the other five quarks' errors are comparable
to, or worse than, those of equation~\eqref{P16:eq:quark_mass_recap}
(Section~\ref{P16:sec:colourshift}), and the apparent precision of
equation~\eqref{P16:eq:e8_ratio_recap}'s integer structure is partly an
artifact of fitting the assignment to the already-known mass hierarchy
rather than predicting it. We record this as a sixth ruled-out
mechanism, distinct from but related to Section~\ref{P16:sec:colourshift}:
the incommensurability principle itself survives intact, but the
specific $E_8$ Coxeter labelling built on top of it does not yet
constitute a derivation of $n_q^{(g)}$.
\end{remark}

%══════════════════════════════════════════════════════════════════════════════
\section{Colour is Invisible at the $\twoI$ Level: A Structural Theorem}
\label{P16:sec:colour_exclusion}
%══════════════════════════════════════════════════════════════════════════════

Proposition~\ref{P16:prop:bosonic_fail} establishes that the colour
$\mathbb{Z}_3$ outer automorphism $\sigma$ of $\twoT$ is not an
independent structure on top of the bosonic triple: it is the
\emph{unique} order-$3$ outer automorphism of $\twoT$, acting on all
three-element orbits simultaneously. A natural complementary question
is whether $\sigma$ can act \emph{at all} on the $\QH=2$ sector, which
is organised not under $\twoT$ but under the full icosahedral group
$\twoI$ via the McKay correspondence $\twoI\leftrightarrow\widehat{E}_8$
(Paper~III). The following theorem answers this in the negative, giving
the first precise group-theoretic statement that colour charge is a
$\QH=3$-specific phenomenon rather than a feature of the framework at
large.

\begin{theorem}[Colour exclusion in the $\QH=2$ sector]
\label{P16:thm:colour_exclusion}
The colour $\mathbb{Z}_3$ automorphism $\sigma$ does not extend to any
automorphism of $\twoI$. Consequently, the colour charge is not a
well-defined quantum number for any field configuration in the
$\QH=2$ sector of the density-feedback Hopfion, and every $\QH=2$
configuration is a colour singlet.
\end{theorem}

\begin{proof}
The proof uses three independently verifiable facts.

\textbf{Fact~1: $\sigma$ is a strictly outer automorphism of $\twoT$}
(Proposition~\ref{P16:prop:bosonic_fail} and the proof thereof). That is,
$\sigma\notin\mathrm{Inn}(\twoT)$: no element of $\twoT$ itself
induces the cyclic permutation
$\Gamma_2\to\Gamma_{2\omega}\to\Gamma_{2\omega^2}\to\Gamma_2$ by
conjugation, since the three colour irreps have distinct characters on
every non-central conjugacy class of $\twoT$.

\textbf{Fact~2: $\mathrm{Out}(\twoI)\cong\mathbb{Z}_2$.}
The binary icosahedral group satisfies $\twoI\cong\mathrm{SL}(2,5)$,
the double cover of $A_5\cong\mathrm{PSL}(2,5)$. Every automorphism
of $\twoI$ preserves the centre $\{\pm1\}$ and therefore projects to
an automorphism of $A_5$; the kernel of this projection is trivial,
giving $\mathrm{Out}(\twoI)\cong\mathrm{Out}(A_5)$. Since $A_5$ is
the unique non-abelian simple group of order $60$, its outer
automorphism group is $\mathbb{Z}_2$ (the $\mathbb{Z}_2$ is realised
by conjugation by any odd permutation in $S_5$), so
$\mathrm{Out}(\twoI)\cong\mathbb{Z}_2$.

\textbf{Fact~3: $N_{\twoI}(\twoT)=\twoT$} (Proposition~\ref{P16:prop:five_copies},
part verified by direct quaternion computation in Paper~XV §9.2).

Now suppose for contradiction that $\sigma$ extends to an automorphism
$\tau\in\mathrm{Aut}(\twoI)$ with $\tau|_{\twoT}=\sigma$. By
Fact~2, every element of $\mathrm{Out}(\twoI)$ has order at most~$2$.
If $\tau$ is outer, it projects to an element of
$\mathrm{Out}(\twoI)\cong\mathbb{Z}_2$, hence has order dividing~$2$
in the quotient. But $\sigma$ has order~$3$, so $\tau^2|_\twoT =
\sigma^2\neq\mathrm{id}$, contradicting $\tau^2$ being inner on
$\twoI$ and therefore necessarily the identity on~$\twoT$ (since any
inner automorphism of $\twoI$ restricts to the identity on $\twoT$ if
it fixes $\twoT$ setwise, as inner autos by elements outside $\twoT$
cannot stabilise $\twoT$ by Fact~3). If instead $\tau$ is inner, then
$\tau=\mathrm{conj}_g$ for some $g\in\twoI$. For
$\mathrm{conj}_g|_\twoT=\sigma$ we need $g\in N_{\twoI}(\twoT)=\twoT$
(Fact~3), so $\mathrm{conj}_g|_\twoT\in\mathrm{Inn}(\twoT)$,
contradicting Fact~1 ($\sigma\notin\mathrm{Inn}(\twoT)$). Both cases
yield contradictions. \qed

\textbf{Consequence.} Since $\sigma$ has no realisation as an
automorphism of $\twoI$, it acts trivially on every $\twoI$-irreducible
representation: for any $\twoI$-module $\rho$, the twist $\sigma^*\rho$
is isomorphic to $\rho$ (there is no non-trivial twist to apply). This
means the three colour labels $\{{\Gamma_2},\Gamma_{2\omega},
\Gamma_{2\omega^2}\}$ are indistinguishable from the $\twoI$-sector's
perspective: no $\twoI$-organised field configuration can prefer one
colour over another. Applied to the $\QH=2$ Hopfion, whose WZW sector
is organised under $\twoI$ via $\twoI\leftrightarrow\widehat{E}_8$
(Paper~III), this means every $\QH=2$ configuration carries zero net
colour: it is a colour singlet.
\end{proof}

\begin{remark}[Dynkin-diagram formulation]
\label{P16:rem:dynkin_colour}
Theorem~\ref{P16:thm:colour_exclusion} has a clean reformulation in terms
of the McKay graphs. The colour $\mathbb{Z}_3$ of $\twoT$ is precisely
the order-$3$ graph automorphism of $\widehat{E}_6$ (the three equal
arms of the $\widehat{E}_6$ diagram rotate into each other). The
analogous automorphism group of $\widehat{E}_8$ (the McKay graph of
$\twoI$) is trivial: the $\widehat{E}_8$ diagram is a linear chain
with a single branch at a non-central node, and no non-identity
permutation of its nodes preserves all adjacencies. Therefore
$\mathrm{Aut}(\widehat{E}_8)=\{1\}$, confirming by a different route
that the colour $\mathbb{Z}_3\subset\mathrm{Aut}(\widehat{E}_6)$ has
no counterpart at the $\widehat{E}_8$ level.
\end{remark}

\begin{remark}[Relation to the $\QH=2\to\QH=3$ symmetry breaking]
\label{P16:rem:symmetry_breaking_colour}
Theorem~\ref{P16:thm:colour_exclusion} supplies the group-theoretic
statement underlying the physical picture: the colour charge
\emph{emerges} in the transition $\QH=2\to\QH=3$, i.e.\ in the
symmetry breaking $\twoI\to\twoT$. Before the breaking, $\sigma$ is
not a symmetry of anything in the $\QH=2$ sector; after it,
$\sigma$ is the defining automorphism of the $\twoT$ colour structure.
This mirrors the standard Higgs-mechanism logic (a symmetry that was
not a symmetry of the vacuum becomes a symmetry of an ordered phase),
but here the role of the ``vacuum'' is played by the $\QH=2$
axisymmetric condensate and the role of the ``ordered phase'' by the
$\QH=3$ trefoil sector, in which the $\mathbb{Z}_3$ axis is
geometrically selected by the three crossing regions
(Proposition~\ref{P16:prop:geom_data}).
The six rules-out mechanisms of
Sections~\ref{P16:sec:vertices}--\ref{P16:sec:e8coxeter} all assumed colour
was already present and sought a further generation label within it;
Theorem~\ref{P16:thm:colour_exclusion} shows instead that colour itself is
the output of the $\twoI\to\twoT$ breaking, not an input available
before it.
\end{remark}

\begin{remark}[The $\QH=1$ sector and the complete ADE hierarchy]
\label{P16:rem:qh1_preview}
Theorem~\ref{P16:thm:colour_exclusion} concerns only the $\QH=2$ sector.
Paper~XVII~\cite{Paper17} characterises the remaining non-vacuum sector $\QH=1$,
identified there as the neutrino sector, and proves colour exclusion
for it by a strictly simpler argument: the McKay group of $\QH=1$ is
the cyclic group $\mathbb{Z}_2$ (the minimal A-type group
$\mathbb{Z}_2\leftrightarrow\widehat{A}_1$), and since
$\gcd(|\mathbb{Z}_2|,|\mathbb{Z}_3|)=\gcd(2,3)=1$, the colour
$\mathbb{Z}_3$ has no non-trivial action on any
$\mathbb{Z}_2$-representation by bare coprimality---no automorphism
structure of $\twoI$ is needed.  Together, the three non-vacuum
sectors complete the top tier of the ADE classification:
\[
  \underbrace{\mathbb{Z}_2\leftrightarrow\widehat{A}_1}_{\QH=1,\;\text{neutrino}}
  \qquad
  \underbrace{\twoI\leftrightarrow\widehat{E}_8}_{\QH=2,\;\text{charged lepton}}
  \qquad
  \underbrace{\twoT\leftrightarrow\widehat{E}_6}_{\QH=3,\;\text{baryon}}.
\]
The confinement hierarchy $\QH=1\subsetneq\QH=2\subsetneq\QH=3$
(maximally free, free, confined) mirrors the A-type to E-type to
E-type progression: the A-type group organises the most weakly
interacting sector, while the two exceptional E-type groups organise
the sectors where electroweak charge ($\widehat{E}_8$) and colour
charge ($\widehat{E}_6$) are present.
\end{remark}

%══════════════════════════════════════════════════════════════════════════════
\section{The String-Tension Energy: A Scoped Open Problem}
\label{P16:sec:stringtension}
%══════════════════════════════════════════════════════════════════════════════

The six mechanisms of Sections~\ref{P16:sec:vertices}--\ref{P16:sec:e8coxeter}
share a common feature: each attempts to fix $n_q$ from a
\emph{static} or \emph{combinatorial} quantity, either the discrete
group theory of $\twoT\subset\twoI$ or a fixed WZW phase. A
qualitatively different candidate, motivated by colour confinement
itself, is examined in this section. Unlike
Sections~\ref{P16:sec:vertices}--\ref{P16:sec:e8coxeter}, what follows is not
a falsified hypothesis but a precisely scoped open problem: we
identify what is required for the candidate to become computable, and
locate exactly where the current framework stops short of supplying
it.

\subsection{Motivation: confinement as the qualitative difference}
\label{P16:sec:confinement_motivation}

A quark cannot exist as an asymptotic state; a lepton can. Paper~XV's
own energy decomposition for the $\QH=3$ baryon,
\begin{equation}
  E_{\mathrm{baryon}} = \underbrace{\Kfb\times J_4}_{\text{profile}}
  + \underbrace{\sigma\cdot 3R_0}_{\text{string tension}}
  + \underbrace{E_\cup}_{\text{junction}},
  \label{P16:eq:baryon_energy_recap}
\end{equation}
already encodes this asymmetry structurally: the $\QH=2$ torus has
\emph{no string network at all} --- ``one tube, no Y-junction, so its
energy is purely the profile term $\Kfb\times J_4$'' (Paper~XV~\cite{Paper15},
Remark~8.5. %\label{P16:rem:yjunction_energy}
The torus and trefoil sectors are accordingly not two geometrically similar
configurations with a modest quantitative difference in stored energy; they
have \emph{qualitatively different energy budgets}, with two of
equation~\eqref{P16:eq:baryon_energy_recap}'s three terms entirely absent
for $\QH=2$. This is the same distinction by which two
geometrically comparable atoms, or a chemical bond and a nuclear
fission fragment, can differ by many orders of magnitude in bound
energy despite similar size scales: shape similarity does not imply
energy-density similarity, and the quantity relevant to a chameleon-type
suppression mechanism is the latter, not the former. A purely
geometric comparison of the shared profile term between the two
sectors (Section~\ref{P16:sec:vertices}'s bare Jacobian ratios, $\sim16\%$)
is accordingly the wrong kind of quantity to examine; the term that is
qualitatively new to the confined sector, $\sigma\cdot 3R_0$, is the
right one.

\subsection{Status of $\sigma$ in Paper~XV: named but unvalued}
\label{P16:sec:sigma_unvalued}

The symbol is used exclusively to explain, qualitatively, why
a numerical saddle-finding algorithm failed to converge (Paper~XV,
open problem~(iii)): the gradient flows minimised only the profile
term $\Kfb\times J_4$ and ``missed the string-tension term
$\sigma\times3R_0$'' (Paper~XV, line preceding
Paper~XV, the remark following the Y-junction energy decomposition). The junction energy $E_\cup$ is
named in equation~\eqref{P16:eq:baryon_energy_recap} but never defined
beyond its name. Both $\sigma$ and $E_\cup$ are therefore currently
free, undetermined quantities in the published framework, not
results awaiting only numerical evaluation of an already-specified
formula.

\subsection{The chameleon-sector field and the Hopf profile are the
  same field}
\label{P16:sec:field_identification}

A natural question, prior to any attempt to compute $\sigma$, is
whether the condensate scalar $\rho$ appearing in the chameleon-type
suppression mechanism of Paper~I and Paper~XIII is the same field as
the Hopf profile function $f(r,z)$ whose curvature-driven
concentration Paper~XV identifies at the trefoil crossing regions, or
whether the shared symbol $\rho$ denotes two unrelated objects. Paper~I
states explicitly: ``Let $\rho:\mathbb{R}^{3,1}\to\mathbb{R}$ be the
field whose gradients carry the suppression... Under the Hopf fibration
embedding $\rho\mapsto f(r,z)$..., the gradient maps as $|\nabla\rho|^2\mapsto
\mathrm{kern}=|\nabla f|^2+\sin^2\!f\cdot A$'' (Paper~I, proof of the
vacuum identification theorem). The chameleon-sector scalar and the
Hopf profile function are therefore the same field under an explicit
map, not two independent fields connected only by notation. Paper~XV's
curvature-driven density enhancement at the crossing regions is automatically
a statement about $\rho$, not merely about a structurally separate quantity $f$.

\subsection{The remaining obstruction: the Y-junction arms are not on
  the tube}
\label{P16:sec:geometric_obstruction}

The field identification of
Section~\ref{P16:sec:field_identification} does not, by itself, make
$\sigma$ computable from the integrals $J_4$, $J_{2a}$ already
evaluated in Paper~XV, because those integrals are defined over the
trefoil tube itself (via the Frenet--Serret construction,
equation~(19) of Paper~XV~\cite{Paper15}), while the Y-junction arms are not part of the
tube.

\begin{proposition}[The Y-junction arms lie off the trefoil curve]
\label{P16:prop:arms_off_curve}
The three Y-junction arms, each connecting the origin to a crossing
midpoint $\mathbf{m}_k$ (Paper~XV's Y-junction Fermat--Steiner
proposition), are not subsets of the trefoil centre curve $\Gamma(s)$.
\end{proposition}

\begin{proof}
By that proposition, $|\mathbf{m}_k|=R_0$ for each
$k$, with the origin as the common endpoint of all three arms. The
origin is the centroid of the three crossing midpoints and is not, in
general, a point of $\Gamma(s)$ for the standard $T_{2,3}$ embedding,
since $\Gamma(s)$ does not pass through its own centroid except at
isolated points unrelated to the arm construction. The crossing
midpoints $\mathbf{m}_k$ themselves are defined as averages
$\tfrac12[\Gamma(t_k)+\Gamma(t_k+\pi)]$ of two \emph{distinct} curve
points (Paper~XV's proof of the same proposition), and so
are generically off the curve as well, except where the curve happens
to pass through its own chord midpoint. The arms, as straight (or
otherwise specified) paths between the origin and each $\mathbf{m}_k$,
therefore run through the three-dimensional bulk region surrounding
the tube, not along $\Gamma(s)$.
\qed
\end{proof}

\begin{remark}[Consequence: no double-counting, but no shortcut
  either]
Proposition~\ref{P16:prop:arms_off_curve} resolves a question worth
stating explicitly, since it could in principle have gone the other
way: if the arms were sub-segments of the tube, the string-tension
energy would risk double-counting energy already present in
$\Kfb\times J_4$. They are not, so no such double-counting occurs ---
but for the same reason, $\sigma$ cannot be read off from $J_4$ or
$J_{2a}$ by restricting an existing integral to a sub-region. A
genuine computation of $\sigma$ requires the condensate field's
behaviour in the bulk region between the tube and the Fermat--Steiner
point, a region the existing tube integrals do not cover.
\end{remark}

\subsection{What a numerical attempt would require}
\label{P16:sec:what_is_needed}

Stating the obstruction precisely is more useful than leaving it
implicit. A numerical computation of $\sigma$ (and, separately,
$E_\cup$) requires three ingredients, none of which is currently
available in the canonical papers:

\begin{enumerate}[label=(\alph*)]
\item \textbf{A bulk field equation.} The parent action's Euler--Lagrange
  equation for $\rho$ (Paper~VII, equation for $\square\rho$ in the
  presence of a source) must be solved in the three-dimensional region
  between the tube surface and the Fermat--Steiner point, sourced by
  the trefoil soliton's own energy-momentum. The only solved
  configurations currently available --- the homogeneous cosmological
  attractor $\rho=\rinfty$ and the spherically symmetric
  Schwarzschild-limit solution $\rho_0(r)=\rinfty e^{-\kappa M/r}$
  sourced by a point mass (Paper~VII) --- are both the wrong boundary-value
  problem: the former has no spatial structure, and the latter is
  sourced by an external point mass rather than by the soliton's own
  extended topological structure.
\item \textbf{A boundary condition on the tube surface.} The relation
  between the tube's own profile value (set by $f_0(\rho)=2\arctan(\rho
  C_3^*)^{-1}$ in Paper~XV's normalised coordinates) and whatever value
  the bulk field $\rho$ takes immediately outside the tube is not
  specified in any paper examined. Without this matching condition,
  the bulk equation of part~(a) has no boundary data at the inner edge
  of its domain.
\item \textbf{A boundary condition at the Fermat--Steiner point.}
  Symmetry (the $\mathbb{Z}_3$ rotational invariance of the three arms,
  Paper~XV, the Y-junction Fermat--Steiner proposition) constrains but does not by itself
  fix the field value or its derivative at the origin; an additional
  physical input (e.g.\ regularity, or a value tied to $\rinfty$) is
  needed to close the boundary-value problem at this end.
\end{enumerate}

While ingredient~(a) (the bulk field equation) was the stated obstacle
for computing $\sigma$ from first principles, it turns out not to be
needed: Proposition~\ref{P16:prop:sigma_ext_formula} shows that
$\sigma_{\mathrm{ext}}\cdot3R_0$ follows directly from the profile
ansatz alone, with no bulk equation required.  Ingredients~(b) and~(c)
(the boundary conditions for $E_\cup$) are supplied numerically below.
We carry out the junction computation.

\begin{proposition}[Junction energy is negligible]
\label{P16:prop:ejunc_zero}
Let $E_\cup(C^*)$ denote the Y-junction energy computed by gradient-flow
minimisation of $E_{\mathrm{geom}}=\Kfb\times J_4$ in the ball
$|\mathbf{x}|<r_{\mathrm{out}}=2$ centred on the Fermat--Steiner origin,
with Dirichlet boundary conditions from Construction~\ref{P16:constr:perstrand} at
$|\mathbf{x}|=r_{\mathrm{bc}}=1.6$, and with the measurement region
$|\mathbf{x}|<r_{\mathrm{junc}}=0.8$.  At the physical saddle value
$C^*=2.5062$ (Proposition~10.6 of Paper~XV~\cite{Paper15} of Paper~XV), the
converged result is
\begin{equation}
  \frac{E_\cup}{E_{\mathrm{baryon}}}
  \;=\; 6.9\times10^{-8},
  \label{P16:eq:ejunc_ratio}
\end{equation}
and the junction energy is similarly negligible at $C^*=2.0$ ($1.8\times10^{-6}$)
and $C^*=3.0$ ($1.8\times10^{-9}$).  The ratio decreases monotonically with
increasing~$C^*$, following approximately $E_\cup\propto C^{*-16}$.
\end{proposition}

\begin{proof}[Evidence]
The boundary layer at $|\mathbf{x}|=1.6$ is at distance $\sqrt{R_0^2+r_{\mathrm{bc}}^2}\approx3.4$
from any crossing region, so the Construction~\ref{P16:constr:perstrand} profile there evaluates to
$f\approx0.057\,\mathrm{rad}\ll\pi/2$: effectively the vacuum state.
The junction domain therefore has vacuum Dirichlet boundary conditions and
contains no portion of the trefoil tube.  Gradient flow converged in 2900 steps
($N=60^3$, $h=0.067$) at $C^*=2.5062$; the converged $|$grad$|$ fell to
$1\%$ of its initial value.  Runs at $C^*=2.0$ and $C^*=3.0$ (2000 steps
each) are consistent: the junction energy grows under flow (the global
$E_{\mathrm{geom}}$ minimiser slightly concentrates energy in the
junction core) but remains at the $10^{-9}$--$10^{-6}$ level of
$E_{\mathrm{total}}$.  Boundary-condition violation throughout was
identically zero (machine precision), confirming the Dirichlet
constraint was enforced exactly.
\qed
\end{proof}

\begin{remark}[What Proposition~\ref{P16:prop:ejunc_zero} establishes]
\label{P16:rem:ejunc_consequence}
Proposition~\ref{P16:prop:ejunc_zero} resolves the question that
ingredients~(b) and~(c) of Section~\ref{P16:sec:what_is_needed} were
required to answer: what is the junction energy?  The answer is zero,
to seven significant figures relative to $E_{\mathrm{baryon}}$.
The baryon energy formula~\eqref{P16:eq:baryon_energy_recap} simplifies to
\begin{equation}
  E_{\mathrm{baryon}} \;=\; \Kfb\times J_4 \;+\; \sigma\cdot 3R_0,
  \label{P16:eq:baryon_simplified}
\end{equation}
with $E_\cup$ discarded without approximation.  Together with
Proposition~\ref{P16:prop:sigma_ext_formula} (which gives $\sigma\cdot3R_0
= 0.04405\,E_{\mathrm{profile}}^*$), this yields
\begin{equation}
  \boxed{E_{\mathrm{baryon}} \;=\; 1.04405\cdot E_{\mathrm{profile}}^*,}
  \label{P16:eq:baryon_final}
\end{equation}
where $E_{\mathrm{profile}}^*=(4/\vp^4)\,J_{2a}^{*2}=0.005041$ in
$\Lcond=1$ units (Theorem~10.1 of Paper~XV~\cite{Paper15}).
The sole remaining open quantity in the $\QH=3$ energy budget is the
absolute value of $E_{\mathrm{profile}}^*$ in physical units, which
requires the $\QH=3$ saddle finder (open problem~(iii) of Paper~XV,
Section~\ref{P16:sec:saddle_prerequisite} of this paper).

The $C^*$ dependence of $E_\cup$ ($\propto C^{*-16}$) also rules out
$E_\cup$ as a source of quark generation-dependence: even if $C^*$
differed by a factor of~$1.5$ between generations, the resulting change
in $E_\cup$ would contribute at the $\lesssim10^{-6}$ level to mass
ratios, six orders of magnitude below the observed quark hierarchy.
\end{remark}

\subsection{What is computable now: the straight-tube total tension}
\label{P16:sec:sigma_computable}

Although the full computation of Paper~XV's $\sigma$ (the tension of
the Y-junction arms, requiring ingredients~(a)--(c) above) is not
currently available, a related quantity \emph{is} computable from
the analytic profile ansatz alone, without the 3D saddle.

Define $\sigma_{\mathrm{tot}}$ as the total energy per unit length of
an infinite, straight Faddeev--Niemi flux tube --- not the three-term
decomposition of equation~\eqref{P16:eq:baryon_energy_recap}, but the
combined energy including the cross-section profile:
\begin{equation}
  \sigma_{\mathrm{tot}}
  \;=\; 2\pi\int_0^\infty
    \Kfb^{(2D)}(r)\;\rho_{J_4}^{(2D)}(r)\;r\,\mathrm{d}r,
  \label{P16:eq:sigma_tot}
\end{equation}
where $\Kfb^{(2D)}$ and $\rho_{J_4}^{(2D)}$ are the 2D (transverse
cross-section) versions of the profile and quartic-density integrands,
evaluated on the cylindrical ansatz
\begin{equation}
  \mathbf{n}(r,\phi)
  = \bigl(\sin f(r)\cos\phi,\;\sin f(r)\sin\phi,\;\cos f(r)\bigr),
  \qquad
  f(r) = 2\arctan\!\bigl((r\,C_3^*)^{-1}\bigr),
  \label{P16:eq:straight_tube_ansatz}
\end{equation}
with $C_3^*=2.5062$ (the physical saddle value, Paper~XIV).
Explicitly:
\begin{align}
  J_{2\mathrm{iso}}^{(2D)} &= \left(\frac{\mathrm{d}f}{\mathrm{d}r}\right)^2
    + \frac{\sin^2 f}{r^2},
  \qquad
  J_{2a}^{(2D)} = \sin^4\!f\cdot J_{2\mathrm{iso}}^{(2D)},
  \qquad
  \Kfb^{(2D)} = J_{2a}^{(2D)} + \mu\,J_{2\mathrm{iso}}^{(2D)},
  \notag\\
  \rho_{J_4}^{(2D)} &= \frac{\sin^2\!f}{r^2}
    \left(\frac{\mathrm{d}f}{\mathrm{d}r}\right)^2,
  \label{P16:eq:2D_integrands}
\end{align}
where $\mu=3-\vp=1.3820$ and $\mathrm{d}f/\mathrm{d}r=-2C_3^*/((rC_3^*)^2+1)$.

\begin{proposition}[Straight-tube tension is computable without the saddle]
\label{P16:prop:sigma_tot}
The integral~\eqref{P16:eq:sigma_tot} converges and evaluates to
\begin{equation}
  \sigma_{\mathrm{tot}} \;=\; 5392.1
  \quad [\text{in natural units where }C_3^*=2.5062],
  \label{P16:eq:sigma_tot_value}
\end{equation}
computed by numerical quadrature (relative error $<10^{-6}$).
Moreover, under the rescaling $u=r\,C_3^*$, all dependence on $C_3^*$
factors out of both numerator and denominator, so the split at
$u=1$ (the tube surface at $r=1/C_3^*$) yields a
\emph{universal constant} independent of $C_3^*$:
\begin{equation}
  \frac{\sigma_{\mathrm{ext}}}{\sigma_{\mathrm{int}}}
  \;=\; \frac{\sigma_{\mathrm{ext}}}{\sigma_{\mathrm{tot}}-\sigma_{\mathrm{ext}}}
  \;=\; 0.04405447\ldots,
  \label{P16:eq:sigma_split_ratio}
\end{equation}
where $\sigma_{\mathrm{int}}$ and $\sigma_{\mathrm{ext}}$ denote the
contributions from $r<1/C_3^*$ and $r>1/C_3^*$ respectively.
\end{proposition}

\begin{proposition}[The exterior string tension is analytically determined]
\label{P16:prop:sigma_ext_formula}
Paper~XV's string tension $\sigma = \sigma_{\mathrm{ext}}$ satisfies
\begin{equation}
  \sigma_{\mathrm{ext}} \cdot 3R_0
  \;=\; 0.04405447\cdots \times \Kfb\!\cdot\!J_4
  \;=\; 0.04405447\cdots \times E_{\mathrm{profile}}^*,
  \label{P16:eq:sigma_ext_formula}
\end{equation}
where $E_{\mathrm{profile}}^* = (4/\vp^4)\,J_{2a}^{*\,2}$ is the
analytic profile energy from Theorem~10.1 of Paper~XV~\cite{Paper15}
(if it exists in the framework). This relation follows purely from the
profile ansatz~\eqref{P16:eq:straight_tube_ansatz} and requires neither
the 3D saddle nor a bulk field equation.

Consequently: (i) $\sigma_{\mathrm{ext}}\cdot 3R_0$ is a
generation-\emph{universal} multiplicative correction to the baryon
energy, equivalent to a universal tower-level shift
\begin{equation}
  \delta n_{\sigma} \;=\; \frac{\ln(1+0.04405)}{\ln\vp}
  \;\approx\; 0.0896,
\end{equation}
applying equally to all three quark generations;
(ii) no separate bulk field computation is required for $\sigma$ ---
it is determined by the profile integral alone; and (iii) the
generation-\emph{dependent} part of the quark mass correction must
reside in $E_\cup$ (the Y-junction energy), not in $\sigma\cdot3R_0$.
\end{proposition}

\begin{remark}[Relation to Paper~XV's $\sigma$, the three-term decomposition, and screening]
\label{P16:rem:sigma_decomposition}
$\sigma_{\mathrm{tot}}$ is \emph{not} the same as Paper~XV's $\sigma$.
In the decomposition~\eqref{P16:eq:baryon_energy_recap}, $\Kfb\times J_4$
counts the internal field energy of the tube cross-section
($\sigma_{\mathrm{int}}$ per unit length), and $\sigma\cdot3R_0$ counts
the exterior tail energy ($\sigma_{\mathrm{ext}}$ per unit length).
Proposition~\ref{P16:prop:sigma_ext_formula} shows that
$\sigma_{\mathrm{ext}}/\sigma_{\mathrm{int}}=0.04405$, so $\sigma\cdot3R_0$
is a 4.4\% correction to the profile energy --- small in absolute terms but
non-negligible at the logarithmic precision needed for quark masses.

The exterior-to-interior ratio is independent of $C_3^*$ because the
profile function $f(u) = 2\arctan(1/u)$ decays as $u^{-1}$ for $u\gg1$
(a power law, not exponential), giving an energy density $\sim u^{-12}$
at large $u$ whose tail integral is determined purely by the shape of
the profile at $u\sim1$ --- the tube surface. Whether this corresponds
to the strong or weak screening regime of a separate chameleon mass
$m_\rho$ is therefore a secondary question: the ratio is fixed by the
profile geometry regardless of the screening interpretation.
\end{remark}

\begin{remark}
\label{P16:rem:three_flows}
The physical saddle of $E_{\mathrm{baryon}}$
(equation~\eqref{P16:eq:baryon_energy_recap}) is a stationary point of all
three terms simultaneously.  Each term requires a different computation:
\begin{enumerate}[label=(\roman*)]
\item \textbf{Tube profile flow.} Minimise $\Kfb\times J_4$ over the
  field configuration inside the tube.  This is what the constrained
  gradient-flow solver of this paper's numerical sections computes; its
  closest approach gives $K_{\min}\approx2294$ ($N=192$, $h=0.05$ run,
  Section~\ref{P16:sec:gradient_flow}).  The saddle value $E_{\mathrm{profile}}^*$
  is also available analytically from Theorem~10.1 of Paper~XV~\cite{Paper15}: $E_{\mathrm{profile}}^* = 18/\vp^{17}
  = 0.005041$ in $L_{\rm cond}=1$ units.  This is a 2D radial profile
  energy, not directly comparable to the 3D grid value $K_{\min}=2294$
  above; the relationship between them requires knowing $L_{\rm cond}$,
  which is the open problem of item~(iii) below.
\item \textbf{String tension.} Compute $\sigma\cdot3R_0$.
  \emph{Resolved analytically by Proposition~\ref{P16:prop:sigma_ext_formula}:}
  no bulk field equation or additional gradient flow is required.
  $\sigma\cdot3R_0 = 0.04405\times E_{\mathrm{profile}}^*$,
  a universal 4.4\% multiplicative correction.
\item \textbf{Junction configuration flow.} Find the stationary
  configuration for the Y-junction's internal structure at the
  Fermat--Steiner point, giving $E_\cup$.
  \emph{Resolved numerically by Proposition~\ref{P16:prop:ejunc_zero}:}
  $E_\cup/E_{\mathrm{baryon}} = 6.9\times10^{-8}$ at $C^*=2.5062$
  (gradient-flow computation in the junction domain, Section~\ref{P16:sec:what_is_needed}).
  The junction energy is negligible to seven significant figures;
  it also does not supply quark-mass generation-dependence, since
  its $C^*$-dependence ($\propto C^{*-16}$) contributes at the
  $\lesssim10^{-6}$ level even for large inter-generation $C^*$ variations.
\end{enumerate}
All three terms are now determined.  The exact baryon energy is:
\begin{equation}
  \boxed{E_{\mathrm{baryon}}
  \;=\; 1.04405\cdots\times E_{\mathrm{profile}}^*,}
  \label{P16:eq:Ebaryon_exact}
\end{equation}
with $E_\cup$ discarded without approximation.  The sole remaining
open problem is the conversion of $E_{\mathrm{profile}}^*$ to physical
units, which requires the $\QH=3$ saddle finder
(Section~\ref{P16:sec:saddle_prerequisite}).
\end{remark}

\begin{remark}[The Chern--Simons spoke for the baryon: $\Delta Q=2$ from the per-strand decomposition]
\label{P16:rem:cs_spoke_quark}
Paper~XVII~\cite{Paper17} establishes a master mass formula for
unconfined sectors,
\begin{equation}
  m^{(n)} = \TCMB\!\left(\frac{\pi^2}{15N}\right)^{\!\!1/4}
  \cdot\vp^{2N}
  \cdot\exp\!\Bigl(4\pi\,\Delta Q - \tfrac{T_{j=1/2}^{(k)}}{N}\Bigr),
  \label{P16:eq:master_mass_recap}
\end{equation}
where $N=Q_{\mathrm{group}}$, the condensate scale
$(\pi^2/(15N))^{1/4}\TCMB$ comes from the CMB identity
$\rho_{\mathrm{CMB}}=N\Lambda^4$, the factor $\vp^{2N}$ is the
Verlinde topological amplitude from the normalisation identity, and
$\mathcal{C}=4\pi\Delta Q$ is the Chern--Simons spoke of
Paper~IV~\cite{Paper4}, with $\Delta Q=Q_H^{\mathrm{cond}}-Q_H^{\mathrm{particle}}$
the inter-sector Hopf charge drop.  For the charged lepton
$\Delta Q=2-1=1$ and for the neutrino $\Delta Q=1-1=0$ (Paper~XVII,
Proposition~\texttt{prop:CS\_spoke}).

\emph{The confined baryon sector.}
Equation~\eqref{P16:eq:master_mass_recap} does not directly apply to
confined sectors.  However, the Chern--Simons spoke can be identified
from the \emph{per-strand} decomposition that is the central
construction of this paper (Construction~\ref{P16:constr:perstrand}).
Each of the three trefoil strands carries its own $S^3$ lift
$\Phi=\chi+t$, which individually realises $\QH=1$ topology
(Paper~XVII, Proposition~\ref{prop:qh1_unknot}; the per-strand lift
is verified by the Whitehead computation of
Section~\ref{P16:sec:perstrand_winding}).  The full baryon carries
$\QH=3$.  The Chern--Simons spoke for the baryon is therefore
\begin{equation}
  \Delta Q^{(q)}
  \;=\; Q_H^{(\text{baryon})} - Q_H^{(\text{single strand})}
  \;=\; 3 - 1
  \;=\; 2,
  \label{P16:eq:DeltaQ_quark}
\end{equation}
the difference between the collective topological charge and the
charge carried by one strand in isolation.  This measures the
Chern--Simons barrier that confines the three strands into the
baryon: a single strand would need to cross a spoke of
$\mathcal{C}^{(q)}=4\pi\times2=8\pi$ to escape from $\QH=3$ back
to the $\QH=1$ free-particle state.

\emph{Leading-order baryon mass.}
Substituting $N=Q_{\mathrm{group}}^{(3)}=3$, $\Delta Q=2$, and the
$(E_6)_1$ T-matrix correction $T_{\mathbf{27}}/N=(5/12)/3=5/36$
(Paper~XV~\cite{Paper15}, the $(E_6)_1$ T-matrix equation) %\eqref{P15:eq:TE6}
into the master formula, and
including the analytically determined string-tension factor
$1.04405$ of Proposition~\ref{P16:prop:sigma_ext_formula}, the
leading-order mass per constituent quark is
\begin{equation}
  m_q^{(0)}
  \;=\; \Lcond^{(q)}\cdot\vp^6
  \cdot\exp\!\bigl(8\pi-\tfrac{5}{36}\bigr)\cdot 1.04405
  \;\approx\; 215.5\;\text{MeV},
  \label{P16:eq:mq_leading}
\end{equation}
where $\Lcond^{(q)}=\TCMB(\pi^2/45)^{1/4}\approx1.607\times10^{-4}$~eV.
Three such constituent quarks give
\begin{equation}
  3\,m_q^{(0)} \;\approx\; 646.5\;\text{MeV}.
  \label{P16:eq:mbaryon_leading}
\end{equation}
The experimental proton mass is $m_p=938.272$~MeV, leaving a gap of
\begin{equation}
  E_\cup^{\mathrm{CS}} \;\approx\; 291.8\;\text{MeV}
  \label{P16:eq:Ecup_required}
\end{equation}
to be supplied by the Y-junction energy.  The three-quark profile and
string-tension terms therefore account for $68.9\%$ of the proton
mass with no free parameters; the remaining $31.1\%$ identifies
$E_\cup^{\mathrm{CS}}\approx291.8$~MeV as the prediction that the junction
boundary-value problem must reproduce.  Note that the numerical
junction flow of Remark~\ref{P16:rem:junction_results} finds the geometric
$E_\cup$ (the field energy inside $|x|<0.8$) to be negligible
($\approx7\times10^{-8}$ of $E_{\mathrm{baryon}}$), consistent with
the topological deconcentration at the Fermat--Steiner point; the
CS-derived $E_\cup^{\mathrm{CS}}\approx291.8$~MeV must therefore arise
from a different physical mechanism than the geometric junction term ---
possibly the QCD binding energy mediated by the string network,
which is not captured by the geometric flow.

\emph{Generation structure.}
The leading formula~\eqref{P16:eq:mq_leading} gives the constituent
quark mass.  Current quark masses are lighter by factors of
$\vp^{-2n_q^{(g)}}$ from the generation-dependent tower level
$n_q^{(g)}$ (Remark~\ref{P16:rem:nq_structure}), exactly as lepton
mass generation splitting comes from $\vp^{-2n_e}$ (Paper~XII).
The derivation of the absolute tower level $n_q^{\mathrm{base}}$
from a $\QH=3$ thermodynamic matching condition (the analogue of
Paper~XII's $n_e=20$) is the remaining open calculation needed to
convert~\eqref{P16:eq:mq_leading} into a full quark mass spectrum.
\end{remark}

$\sigma_{\mathrm{tot}}$ is \emph{not} the same as Paper~XV's $\sigma$.
In the decomposition~\eqref{P16:eq:baryon_energy_recap}, $\Kfb\times J_4$
already counts the internal field energy of the tube cross-section
(the profile energy), and $\sigma\cdot3R_0$ counts only the
\emph{additional} energy from the tube-vacuum interface and exterior
condensate tails --- a separate contribution.  $\sigma_{\mathrm{tot}}$
counts \emph{both}: it is the total energy per unit length of a straight
tube, equal to $(\Kfb\times J_4)/L + \sigma_{\mathrm{ext}}$ in the
limit of an infinitely long straight tube, where $\sigma_{\mathrm{ext}}$
is the exterior contribution Paper~XV calls $\sigma$.

The relative size of $\sigma_{\mathrm{ext}}$ versus $\Kfb\times J_4$
depends critically on the condensate's screening length $m_\rho^{-1}$
(the chameleon mass of Paper~I/VII).  Linearising the bulk field
equation around the vacuum value $\rho_\infty$ gives a Yukawa equation,
\begin{equation}
  -\nabla^2(\delta\rho) + m_\rho^2\,\delta\rho = S(\mathbf{x}),
\end{equation}
with Green's function $G(\mathbf{x},\mathbf{x}')=e^{-m_\rho|\mathbf{x}-\mathbf{x}'|}/(4\pi|\mathbf{x}-\mathbf{x}'|)$.
The resulting exterior energy per unit length scales as
$\sigma_{\mathrm{ext}}\sim\lambda^2/(8\pi m_\rho)$, where $\lambda$ is
the source strength per unit length (proportional to the profile energy
per unit length divided by $m_\rho$).
\begin{enumerate}[label=(\roman*)]
\item \textbf{Strong screening} ($m_\rho\gg1/R_0$, i.e.\
  $m_\rho R_0\gg1$): $\sigma_{\mathrm{ext}}$ is exponentially small;
  the baryon energy is dominated by the profile term $\Kfb\times J_4$,
  and $\sigma_{\mathrm{ext}}\cdot3R_0$ is a sub-leading correction.
\item \textbf{Weak screening} ($m_\rho\ll1/R_0$): $\sigma_{\mathrm{ext}}$
  is large; the string-tension term dominates, as in QCD.
\end{enumerate}
In the framework's natural units with $C_3^*=2.5062$ and $R_0=3$, the
tube radius is $1/C_3^*\approx0.40$ and $m_\rho\,R_0\approx7.5$ if
$m_\rho\sim C_3^*$ --- placing the system in the strong-screening
regime.  However, $m_\rho$ is not specified in any published paper in
this series: Paper~VII gives the bulk field equation sourced by a
point mass, not by a soliton, and no chameleon mass appears as a
numerical output.  Whether the framework is in the strong or weak
screening regime is therefore currently an open question, and with it
whether $\sigma_{\mathrm{ext}}\cdot3R_0$ is a leading or sub-leading
contribution to $E_{\mathrm{baryon}}$.

\begin{remark}
What $\sigma_{\mathrm{tot}}=5392.1$ does supply unconditionally is a
\emph{consistency bound}: at the physical saddle,
$E_{\mathrm{baryon}}/(3R_0)\leq\sigma_{\mathrm{tot}}$, since the
trefoil's curvature and the Junction contribution cannot add energy
beyond the straight-tube total.  This bound is independent of the
screening regime.
\end{remark}

\begin{remark}[Junction flow: numerical results, arm-stub analysis, and generation-dependence conclusion]
\label{P16:rem:junction_results}
A Dirichlet gradient-flow solver was implemented for the junction
domain $|x|<r_{\mathrm{outer}}=2$ centred on the Fermat--Steiner point
(the origin, verified exactly), with the outer sphere
$|x|=r_{\mathrm{boundary}}=1.6$ held fixed to Construction~\ref{P16:constr:perstrand} values
throughout (boundary-condition error $=0$ to machine precision in all runs).

Runs were carried to convergence ($|\nabla E|<1\%$ of initial) at two
$C_3^*$ values:
\begin{equation}
  E_{\cup}^{\mathrm{meas}}(C_3^*=2.5062) = 0.033697,\qquad
  E_{\cup}^{\mathrm{meas}}(C_3^*=3.0)    = 0.001139,
  \label{P16:eq:Ejunc_measured}
\end{equation}
(solver units; asymptotic correction $<2\%$).  The two-point scaling is
$E_\cup^{\mathrm{meas}}\sim w^{18.8}$ where $w=1/C_3^*$, consistent
with the steep tail of the profile function.

The measurement $E_\cup^{\mathrm{meas}} = K_{\mathrm{global}}\times
\int_{|x|<r_{\mathrm{junc}}}\rho_{J_4}\,\mathrm{d}^3x$ uses the
globally-defined $K$, which does not decompose into local contributions.
It therefore differs from Paper~XV's $E_\cup$, which is defined residually
as $E_{\mathrm{baryon}} - (K_{\mathrm{fb}}\cdot J_4) - \sigma\cdot3R_0$;
the latter requires the full 3D baryon saddle and is not accessible from
a junction-only flow.

The converged junction field concentrates $\rho_{J_4}$ at the
Fermat--Steiner point by a factor of $\approx9.4$ relative to
Construction~\ref{P16:constr:perstrand} (ratio $E_\cup^{\mathrm{meas}}/E_\cup^{\mathrm{init}}
\approx9.4$ at $C_3^*=2.5062$), confirming that the Y-junction acts as
a fourth topological concentration locus analogous to the three crossings
of Section~\ref{P16:sec:arc_segment}.  Converting to natural units via the
scale factor $1.033\times10^{-8}$, the excess
$\Delta E_\cup\approx3.1\times10^{-10}\,L_{\mathrm{cond}}^{-1}$,
which is $\approx6\times10^{-8}$ of $E_{\mathrm{profile}}^*$.

This excess is far too small (by seven orders of magnitude) to account
for the factor-of-$2$--$49$ quark mass ratios between generations.

\paragraph{Arm-stub accounting.}
The measurement sphere $|x|<r_{\mathrm{junc}}=0.8$ contains the
inner portions of the three arm stubs that emanate from the
Fermat--Steiner point. For three straight tubes of this radius and
length $r_{\mathrm{junc}}=0.8$, the expected $\rho_{J_4}$ contribution
is
\begin{equation}
  J_4^{\mathrm{arms}} \;=\; 3\,r_{\mathrm{junc}}\,\times
  \underbrace{2\pi\int_0^\infty \rho_{J_4}^{(2D)}(r)\,r\,\mathrm{d}r}_{=\,105.2\,
  \text{at }C_3^*=2.5062}
  \;\approx\; 252.
\end{equation}
The actual measured value is $\int_{|x|<0.8}\rho_{J_4}\,h^3
= E_\cup^{\mathrm{meas}}/K_{\mathrm{global}} \approx 6\times10^{-5}$,
which is a factor of $\sim4\times10^6$ \emph{smaller} than three
straight-tube arms of the same length would give.

This is a genuine physical result, not a grid artefact: the Y-junction
is a \emph{topological deconcentration locus}. The three Q$_H=1$ flux
tubes emanating from the Fermat--Steiner point carry equal topological
charge, and their $\rho_{J_4}$ contributions cancel at the merge
point by destructive interference of the three 2-form densities
$F_{ij}^{(k)}$. Far from concentrating $\rho_{J_4}$, the junction
suppresses it by six orders of magnitude relative to the straight-tube
arm bodies. The factor of $\approx9.4$ growth of $E_\cup^{\mathrm{meas}}$
during gradient flow (from the Construction~\ref{P16:constr:perstrand} initial value to the
relaxed value) reflects genuine junction restructuring but still leaves
the absolute value far below any arm-stub contribution.

\paragraph{Generation-dependence conclusion.}
Converting to natural units via the scale factor $1.033\times10^{-8}$,
the excess $\Delta E_\cup\approx3.1\times10^{-10}\,L_{\mathrm{cond}}^{-1}$,
which is $\approx6\times10^{-8}$ of $E_{\mathrm{profile}}^*$, or seven
orders of magnitude too small to account for the factor-of-$2$--$49$
quark mass ratios between generations.

The generation-dependent correction to the quark mass formula does not
arise from the geometric Y-junction energy.  The topological deconcentration
at the junction means $E_\cup$ is structurally suppressed, not accidentally
small: the Z$_3$-symmetric merging of three equal charges is the reason,
and no generation-dependent perturbation of the junction geometry can
change this by the required seven orders of magnitude.  The quark
generation hierarchy must come from a different mechanism --- most likely
the dynamical tower-level condition for $n_q^{(g)}$ analogous to
Paper~XII's thermodynamic matching for $n_e=20$.
\end{remark}

%══════════════════════════════════════════════════════════════════════════════
\section{Running the Colour-Shift Predictions to the Comparison Scale}
\label{P16:sec:rgrunning}
%══════════════════════════════════════════════════════════════════════════════

Sections~\ref{P16:sec:colourshift} and~\ref{P16:sec:e8coxeter} compared
condensate-scale predictions directly against PDG values with no
running step, and Remark~\ref{P16:rem:scale_ambiguity} flagged this as a
separate, unresolved issue. A one-loop estimate reported earlier in
this investigation (not reproduced here) widened rather than closed
the gap, but used a single effective flavour number across the entire
range and is accordingly unreliable. This section repeats the
calculation properly: one-loop running of $\alpha_s$ and of the quark
mass, with flavour-threshold matching at each quark mass scale
($n_f=6\to5\to4\to3$), anchored to the measured $\alpha_s(M_Z)=0.1180$.
Higher-loop corrections are not included; the result below should be
read as a first proper pass, not a precision calculation.

\subsection{Method}

The colour-shift formula (equation~\eqref{P16:eq:quark_mass_recap}) gives
each quark's mass at $\LUV\approx6.8\times10^{17}$~GeV. We run
$\alpha_s$ down from $\LUV$ through the top, bottom, and charm
thresholds using the one-loop relation
$\alpha_s(\mu)=\alpha_s(\mu_0)/\bigl[1+\tfrac{\alpha_s(\mu_0)}{2\pi}
b_0\ln(\mu/\mu_0)\bigr]$, $b_0=(33-2n_f)/3$, then run each quark mass
through the same thresholds using
$m(\mu)=m(\mu_0)\bigl[\alpha_s(\mu)/\alpha_s(\mu_0)\bigr]^{\gamma_0/(2b_0)}$
with the universal one-loop anomalous dimension $\gamma_0=8$. Light
quarks ($u,d,s$) are compared at $\mu=2$~GeV; charm and bottom at
their own scale $m_q(m_q)$; top at the one-loop-run value compared to
the pole mass, acknowledging this comparison is the least reliable of
the six given the absence of a pole-mass matching correction.

\subsection{Result}

\begin{table}[h]
\centering
\begin{tabular}{lrrrr}
\toprule
Quark & $m(\LUV)$ [MeV] & $m$(comparison scale) [MeV] & PDG value [MeV] & Ratio \\
\midrule
$u$ & $0.314$   & $1.300$    & $2.16$      & $0.602$ \\
$d$ & $0.120$   & $0.497$    & $4.67$      & $0.106$ \\
$s$ & $22.56$   & $93.32$    & $93.40$     & $\mathbf{0.999}$ \\
$c$ & $59.07$   & $265.7$    & $1275$      & $0.208$ \\
$b$ & $522.2$   & $1930$     & $4180$      & $0.462$ \\
$t$ & $1367$    & $3554$     & $172570$    & $0.0206$ \\
\bottomrule
\end{tabular}
\caption{Colour-shift predictions run from $\LUV$ to the PDG
  comparison scale, one loop, with flavour-threshold matching.}
\label{P16:tab:rg_running}
\end{table}

\begin{proposition}[The strange-quark prediction matches PDG at the
  sub-percent level]
\label{P16:prop:strange_match}
Running the down-type colour-shift prediction $m_s(\LUV)=22.56$~MeV
to $\mu=2$~GeV gives $m_s(2\,\mathrm{GeV})=93.32$~MeV, a $0.09\%$
deviation from the PDG value $93.40$~MeV.
\end{proposition}

\begin{proof}
Direct computation, as in Section~\ref{P16:sec:rgrunning}.\,Method;
verified independently with exact (non-truncated) floating-point
arithmetic.
\qed
\end{proof}

\begin{remark}[This is not the same coincidence as Proposition~\ref{P16:prop:e8_mixed}]
\label{P16:rem:two_strange_results}
The strange quark is, independently, the one quark for which
Section~\ref{P16:sec:e8coxeter}'s $E_8$ Coxeter route is exactly
predictive ($n_s=0$, i.e.\ $m_s=m_\mu$ at the condensate scale,
forced by $T_{\mathrm{total}}(s)=T_\mu$). It is worth checking
explicitly whether Proposition~\ref{P16:prop:strange_match} is a
restatement of that same fact. It is not: the $E_8$ route's
condensate-scale value is $m_s=m_\mu=105.66$~MeV, whereas the
colour-shift route's condensate-scale value is $m_s(\LUV)=22.56$~MeV
--- a different number, arrived at by a different formula, that
becomes close to the PDG value only after a $\approx4.14\times$
one-loop running enhancement is applied. We checked whether this
running factor itself matches a natural quantity in the framework
($\vp^3=4.236$ is the closest candidate, $2.4\%$ off) and found no
clean match; since the running factor depends on externally-sourced
inputs ($\alpha_s(M_Z)$, the threshold locations, the loop order) not
predicted by the framework, a numerical coincidence with such a
quantity would not carry structural weight in any case. What can be
said is that two structurally independent calculations ---
a static condensate-scale phase coincidence, and an RG-improved
colour-shift formula --- both single out the strange quark for
unusually good agreement with measurement. Whether this reflects a
single underlying mechanism not yet identified, or two unrelated
coincidences, is left open.
\end{remark}

\subsection{The residual pattern}

The five remaining quarks miss the PDG comparison values by factors of
$1.7$ ($u$) to $49$ ($t$), and the misses are systematic in two
respects worth recording. First, every prediction in
Table~\ref{P16:tab:rg_running} undershoots the measured value (all ratios
are less than $1$), with the single exception of the strange quark's
near-exact match; there is no case of overshooting. Second, within
each type the deviation grows with generation: $u\,(40\%)\to
c\,(4.8\times)\to t\,(49\times)$ for up-type, and $d\,(9.4\times)\to
s\,(\mathrm{exact})\to b\,(2.2\times)$ for down-type, though the
down-type sequence is non-monotonic in the same way
Proposition~\ref{P16:prop:colourshift_fail} already found. A
systematically one-sided, generation-growing deficit is the
qualitative signature expected of a missing
\emph{generation-dependent, always-positive} energy contribution ---
consistent with, though not yet a derivation from, the string-tension
energy of Section~\ref{P16:sec:stringtension}, which is present only for
the confined ($\QH=3$) sector and absent for leptons. We do not claim
this connection is established; Section~\ref{P16:sec:stringtension}'s own
analysis shows $\sigma$ is not yet computable, so this pattern cannot
presently be checked against it quantitatively. It is recorded here
as the most concrete piece of quantitative evidence in this paper that
such a term, if and when computed, is being asked to do real work
rather than serving as a post-hoc patch.

\begin{remark}[A further caveat: external inputs to the running
  itself]
\label{P16:rem:external_inputs}
The calculation of this section anchors $\alpha_s$ to the measured
$\alpha_s(M_Z)=0.1180$ and uses PDG values for the threshold
locations $m_t,m_b,m_c$, rather than values derived from the
framework itself. If the framework's own $\Lambda_{\mathrm{QCD}}$
(Paper~V) implies a different $\alpha_s(M_Z)$, the running factor
shifts systematically for all six quarks at once, which could move
Table~\ref{P16:tab:rg_running}'s ratios up or down together without
changing their relative pattern. A fully self-contained test would
re-derive $\alpha_s$ from the framework's own parameters rather than
importing it from PDG; Section~\ref{P16:sec:selfconsistent_as} attempts
this.
\end{remark}

\subsection{A self-consistent $\alpha_s(\LUV)$, and its effect on the
  table}
\label{P16:sec:selfconsistent_as}

Paper~V independently derives $\alpha_s(\LUV)$, rather than only
quoting the measured $\alpha_s(M_Z)$ and running it up. The
derivation chains two results we verified directly. First, the
effective dimension governing the colour coupling is
$\deff^{\mathrm{exact}}=43/14$ (Paper~V's exact-$\deff$
proposition), from the exact trigonometric
integral
\begin{equation}
  \deff-3 = \frac{\int_0^\pi\tfrac12\cos^2\theta\sin^5\theta\,d\theta}
                 {\int_0^\pi\sin^5\theta\,d\theta} = \frac{1}{14},
  \label{P16:eq:deff_recap}
\end{equation}
which we re-evaluated symbolically and confirmed exactly
($8/105$ over $16/15$, giving $1/14$ precisely). Second, a WZW
threshold factor $K_{\mathrm{th}}=\sin(\pi/5)=\sqrt{3-\vp}/2$ (Paper~V's
WZW threshold-factor theorem) follows from the vacuum modular amplitude
of the $\SU(2)_3$ WZW model, with $k=3$ fixed independently elsewhere
in the framework (the pentagon mechanism of Paper~III), not chosen to
fit any measured quantity. Neither input references PDG data. Combining
them gives
\begin{equation}
  \alpha_s^{\mathrm{matched}}(\LUV)
  = \frac{C_F\sin(\pi/5)}{2\vp^{43/7}}
  \approx 0.02039,
  \label{P16:eq:as_LUV_recap}
\end{equation}
a number assembled entirely from framework-internal quantities
($C_F=4/3$, $\vp$, $\pi/5$), with PDG values used only afterward, to
\emph{verify} the result (agreement to $0.3\%$ with the value implied
by running the measured $\alpha_s(M_Z)$ up to $\LUV$), not to
construct it.

\begin{proposition}[The self-derived $\alpha_s(\LUV)$ does not
  materially change Table~\ref{P16:tab:rg_running}]
\label{P16:prop:as_selfcons}
Repeating the calculation of Section~\ref{P16:sec:rgrunning}\,Method with
$\alpha_s(\LUV)=0.02039$ in place of the PDG-anchored value
$0.02030$ changes every entry of Table~\ref{P16:tab:rg_running} by at most
a few percent, leaving the qualitative pattern (one near-exact match,
five large undershoots, generation-growing deviation) unchanged.
\end{proposition}

\begin{proof}
Direct recomputation: with $\alpha_s(\LUV)=0.02039$, one-loop running
through the same flavour thresholds gives $\alpha_s(2\,\mathrm{GeV})=
0.2797$ (versus $0.2643$ previously), a $5.8\%$ shift in the coupling
that propagates to a comparably small shift in each mass prediction
via the $\gamma_0/(2b_0)$ exponent. The strange-quark prediction moves
from $93.32$~MeV ($0.09\%$ low) to $95.98$~MeV ($2.8\%$ high); the
other five ratios shift by similarly small amounts and remain off by
factors of $1.7$ to $48$.
\qed
\end{proof}

\begin{remark}[Why this result is still informative despite the null
  outcome]
The two anchoring choices for $\alpha_s(\LUV)$ agree to $0.3\%$ by
construction (this is the content of Paper~V's own consistency check),
so Proposition~\ref{P16:prop:as_selfcons}'s null result was expected once
that agreement is taken into account: there was no room for the
self-consistent value to differ enough from the PDG-anchored one to
close gaps of a factor of $2$ or more. What this rules out, concretely,
is the hypothesis that the residual mismatches in
Table~\ref{P16:tab:rg_running} are an artifact of importing $\alpha_s(M_Z)$
from PDG rather than deriving it. They are not: even with every input
to the running traceable to a framework-internal derivation except the
threshold \emph{locations} $m_t,m_b,m_c$ themselves (still taken from
PDG, since predicting those is the object of the present
investigation, not a usable input to it), the same systematic pattern
persists. The source of the residual pattern, whatever it is, lies
elsewhere than in the choice of $\alpha_s$ anchor.
\end{remark}

\subsection{Is the comparison itself the wrong kind of comparison?}
\label{P16:sec:comparison_target}

The calculations of Sections~\ref{P16:sec:rgrunning}--\ref{P16:sec:selfconsistent_as}
take for granted that the right object to compare the framework's
$\LUV$-scale prediction against, after running, is a PDG-quoted running
mass ($\overline{\mathrm{MS}}$ at $2$~GeV, or $m_q(m_q)$). A natural
question, given that quarks are never observed in isolation
(Section~\ref{P16:sec:confinement_motivation}), is whether real
experiments extract something more like a discrete
\emph{production threshold} or \emph{onset energy} for each quark,
rather than a continuously-running Lagrangian parameter --- in which
case the comparison target itself, not merely the running procedure,
might need reconsidering.

We checked this directly against the quark-mass-definition literature.
The consistent answer is that no such onset-energy quantity exists.
Quark masses are explicitly described as scheme-dependent
theoretical constructs at every scale, not as observed thresholds:
``there is no unique definition of the quark mass\dots{} the quark
mass is a purely theoretical notion''~\cite{Hoang1999,Hoang2002}, and ``quark
masses should be seen more like coupling constants than like physical
parameters''~\cite{Rodrigo1997}. For light quarks ($u,d,s$), DIS
primarily constrains parton distribution functions rather than masses
directly; light-quark masses enter through chiral perturbation theory
and lattice QCD, and are conventionally quoted via the running
$\overline{\mathrm{MS}}$ mass at $\mu=2$~GeV, exactly the target used
in Section~\ref{P16:sec:rgrunning}. There is no rival ``ejection energy''
definition for these quarks waiting to be used instead.

\begin{remark}[A real, adjacent distinction worth correcting for]
\label{P16:rem:pole_mass_caveat}
For heavy quarks near a production threshold (e.g.\ $t\bar t$ or
quarkonium production in $e^+e^-$ annihilation), the \emph{pole mass}
--- the comparison scheme used for the top quark in
Table~\ref{P16:tab:rg_running} --- is known to carry an intrinsic,
scheme-independent ambiguity of order $\Lambda_{\mathrm{QCD}}\sim300$--$400$~MeV
(the renormalon ambiguity), because it is sensitive to long-distance,
low-momentum physics in a way that running masses are not. This is
precisely why ``threshold masses'' (the PS mass, the 1S mass, and
related short-distance schemes) were introduced: not as production
onset energies in the sense the present investigation was testing for,
but as alternative short-distance schemes engineered to behave well
in the near-threshold kinematic regime~\cite{Hoang1999,Hoang2002,Beneke1998}.
This means the top-quark row of Table~\ref{P16:tab:rg_running}, which
compares against the pole mass, uses the single least reliable of the
standard comparison schemes for a quark this close to its own
production threshold; a threshold-mass comparison would be more
appropriate, though we have not attempted it here.
\end{remark}

\begin{remark}[What this does and does not settle]
The absence of an onset-energy definition in the literature does not
mean the framework's own $Q_H=3$ sector is wrong to think of quark
generation as a transient, energy-gated excitation
(Section~\ref{P16:sec:confinement_motivation}'s motivation for this whole
line of investigation remains intact as a statement about the
\emph{framework's internal physical picture}). It means only that
\emph{PDG-extractable quark mass} is not, and is not intended to be,
a direct readout of such an onset energy: it is one of several
scheme-dependent Lagrangian parameters, chosen for calculational
convenience in different kinematic regimes. Any comparison between the
framework's predictions and measurement must therefore go through one
of the existing schemes (running mass, pole mass, or a threshold
scheme), as in Sections~\ref{P16:sec:rgrunning}--\ref{P16:sec:selfconsistent_as},
rather than through a new, as-yet-undefined onset-energy quantity.
\end{remark}

%══════════════════════════════════════════════════════════════════════════════
\section{The Pole-Mass Question: A Saddle-Point Prerequisite}
\label{P16:sec:saddle_prerequisite}
%══════════════════════════════════════════════════════════════════════════════

Remark~\ref{P16:rem:pole_mass_caveat} raises a natural follow-up question:
does the pole mass's renormalon ambiguity correspond, in this
framework, to some identifiable geometric feature of the $\QH=3$
condensate --- a ``stability zone,'' or its absence? As with
Section~\ref{P16:sec:stringtension}, we do not answer this question here;
we show precisely what stands between asking it and answering it.

\subsection{Motivation: the renormalon ambiguity as a confinement
  signature}

The pole mass's ambiguity, of order $\Lambda_{\mathrm{QCD}}\sim300$--$400$~MeV
(Section~\ref{P16:sec:comparison_target}), is not an accident of
perturbation theory: it reflects the impossibility of cleanly
separating a quark's own energy from the inseparable cloud of soft
gluon exchanges that confinement guarantees it is permanently dressed
in. A natural question is whether this framework's own description
of the $\QH=3$ sector contains an analogous obstruction --- not as a
perturbative artifact, but as a structural feature of the soliton's
energy functional itself.

\subsection{What is already established: solitons as constrained
  saddle points}

Any Hopfion, regardless of $\QH$, is a saddle point, not a true
minimum, in the full field-configuration space. By Derrick's theorem,
rescaling the soliton's spatial size ($\mathbf{x}\mapsto s\mathbf{x}$)
can always lower the energy in one direction (the dilation mode),
making the trivial vacuum the global minimum; the soliton is only a
minimum within its own topological sector once this mode is removed
(Paper~I, Remark on the Derrick scale-barrier). For $\QH=2$, this is
fully resolved: the scale-invariant geometric-mean functional
$E_{\mathrm{geom}}=(\Ja+\ms\Jiso)\cdot\Jfour$ (Paper~I,
equation~\eqref{P16:eq:Egeom_recap} below) projects out the dilation mode
exactly, and the resulting \emph{saddle snapshot}
(Paper~I's saddle-snapshot definition) is a genuine constrained
local minimum, on which ordinary gradient descent converges reliably:
\begin{equation}
  E_{\mathrm{geom}}[f;\ms] \;\equiv\; \bigl(\Ja+\ms\,\Jiso\bigr)\cdot\Jfour,
  \label{P16:eq:Egeom_recap}
\end{equation}
exactly scale-invariant under $\mathbf{x}\mapsto s\mathbf{x}$.

\subsection{What is different for $\QH=3$: a documented, more severe
  failure}

Paper~XV's own open problem~(iii) reports a failure that is
qualitatively worse than an unresolved dilation mode. Every numerical
method attempted --- direct $\Kfb$ minimisation, fixed-$\lam$ flows at
$\lam\in\{\vp^5,\vp^6\}$, and projected-gradient flow at fixed $J_4$
--- fails to converge to any saddle at all on the grids tested. The
configuration instead either undergoes topology collapse ($J_4\to0$,
the trefoil untying) or shows what Paper~XV calls a \emph{three-lobed
energy landscape}, a non-convexity attributed explicitly to
self-interaction between the three crossing regions (Paper~XV,
Remark on the absence of a $\lam_3$ saddle), not to the dilation mode
alone. This three-lobe instability is specific to having three
crossing regions; it has no counterpart in the single-tube $\QH=2$
case, and Paper~I's geometric-mean resolution --- designed to remove
exactly one mode (dilation) --- is not, on its own, designed to
address an instability arising from interaction \emph{between}
multiple sub-regions of the same configuration.

\begin{remark}[The structural parallel, stated precisely]
\label{P16:rem:saddle_renormalon_parallel}
The three-lobe self-interaction obstruction is a non-speculative,
framework-native analogue of the renormalon ambiguity: both reflect
the impossibility of cleanly separating ``the soliton's own energy''
from an inseparable surrounding structure --- the soft-gluon cloud in
real QCD, the mutual interaction of the three crossing regions here.
This is offered as a structural parallel, not a derivation: the two
obstructions arise in different mathematical settings (a perturbative
Borel-transform ambiguity in one case, a numerical saddle-finding
failure in the other), and nothing in this section shows they are the
same calculation.
\end{remark}

\subsection{The present answer to the motivating question}

Sections~\ref{P16:sec:saddle_prerequisite}.2--\ref{P16:sec:saddle_prerequisite}.3
support a direct, if deflating, answer: at present the pole mass does
not correspond to anything in this framework, because the framework
does not yet have a confirmed $\QH=3$ saddle point for any static
energy --- pole-mass-like or otherwise --- to be assigned to. Without
a genuine stationary configuration, there is no well-defined static
trefoil energy for which the question ``does this exhibit a
renormalon-like ambiguity, or a geometric stability zone, or an
oscillation'' is even precisely posed. Paper~XV's open problem~(iii)
is accordingly the literal prerequisite for the oscillation/stability-zone
question raised in this investigation, not a separate, optional piece
of background.

\begin{remark}[A further complication, raised independently of the
  saddle-finding problem itself]
\label{P16:rem:transient_complication}
It is also possible that no single, time-independent saddle is the
right object to look for at all. If the $\QH=3$ configuration is
populated only transiently, during an energy perturbation whose
magnitude determines which generation (if any) is accessible, then
the relevant object may be a family of configurations parametrised by
the perturbation energy, rather than one static saddle. We do not
attempt to formalise this here; it is recorded as a second, independent
reason the present saddle-finder problem, even once solved, may not
by itself be sufficient to pose the oscillation question precisely.
\end{remark}

\subsection{What solving the saddle-finder problem would require}

Stating this precisely, in the same spirit as
Section~\ref{P16:sec:what_is_needed}: Paper~XV identifies the cause of the
failure (inter-tube field overlap at the tested initialisation,
$C^*=1.5$, giving a tube radius only $2.6\times$ the minimum
inter-tube distance) and three candidate remedies, none yet attempted
to completion: Battye--Sutcliffe simulated annealing, constrained
flow at fixed $J_4$, or reinitialisation closer to the conjectured
saddle value $C^*\approx2.506$ (separation $4.4$ tube radii). We treat
locating a genuine $\QH=3$ saddle by one of these routes as the
necessary next computational step, prior to any attempt to test the
oscillation hypothesis raised in this investigation.

\subsection{A candidate structural cause: a forced phase discontinuity
  at each crossing region}
\label{P16:sec:phase_discontinuity}

Independently of the inter-tube overlap mechanism identified by
Paper~XV, the numerical ansatz used by the existing $\QH=3$ solver
(and, by construction, every solver in this series to date) contains
a second, distinct candidate obstruction to saddle convergence: a
genuine, resolution-independent discontinuity in the field $n$ at the
geometric midline of every crossing region.

\begin{proposition}[The single-nearest-point ansatz is discontinuous
  at every crossing]
\label{P16:prop:phase_jump}
Let $n(\mathbf{x}) = (\sin f_0(\rho)\cos\theta,\,\sin f_0(\rho)\sin\theta,\,\cos f_0(\rho))$
with $\rho(\mathbf{x})$ the distance from $\mathbf{x}$ to the nearest
point on the trefoil centreline $\Gamma(t)$ and
$\theta(\mathbf{x})=3\,t_{\mathrm{near}}(\mathbf{x})$, where
$t_{\mathrm{near}}$ is the curve parameter of that nearest point. At
each of the three crossing regions, the internal phase $\theta$ jumps
by exactly $\pi$ across the surface where nearest-point ownership
switches between the over- and under-strand, and this jump is not
removable by mesh refinement: it is forced by the curve's own
geometry, independent of $C^*$, $R_0$, or $r_0$.
\end{proposition}

\begin{proof}
By the antipodal isometry $(t,z)\mapsto(t+\pi,-z)$ exchanging the
over- and under-strand at each crossing (Paper~XV, the constant
inter-strand distance result), the over-strand point at parameter
$t_1$ and the under-strand point at $t_1+\pi$ are equidistant from
the crossing's geometric midpoint by symmetry, so ownership switches
exactly at that midpoint. The assigned phases there are
$\theta_{\mathrm{over}}=3t_1$ and
$\theta_{\mathrm{under}}=3(t_1+\pi)=3t_1+3\pi\equiv 3t_1+\pi
\pmod{2\pi}$, a difference of exactly $\pi$, independent of $t_1$ and
hence identical at all three crossings (verified directly for
$t_1\in\{\pi/6,\,5\pi/6,\,3\pi/2\}$: each gives
$\theta_{\mathrm{over}}=\pi/2$, $\theta_{\mathrm{under}}=3\pi/2$, jump
$=\pi$ exactly). The jump survives in the continuum limit because it
is a discontinuity in the assignment $\mathbf{x}\mapsto t_{\mathrm{near}}(\mathbf{x})$
itself (a property of the nearest-point map, not of the grid), not an
artifact of finite spacing.
\qed
\end{proof}

\begin{remark}[The jump is not a removable gauge artifact]
\label{P16:rem:not_a_pole}
A discontinuity in an azimuthal-type phase is harmless where it
multiplies $\sin f_0=0$ (the familiar coordinate singularity at the
poles of $S^2$, as for ordinary spherical coordinates at the poles).
At the crossing midline, however, $\rho$ equals the half-separation
$r_0$ between the two strand centrelines (by the symmetry used in the
proof above), giving $\sin f_0(r_0\,C_3^*)=0.275$ at the values used
in Paper~XV ($r_0=0.874$, $C_3^*=2.5062$): not small. The resulting
jump in $n$ itself,
$|\Delta n| = 2\sin f_0(r_0 C_3^*) \approx 0.55$
(out of a maximum possible antipodal-flip value of $2$), is a
genuine, physically substantial discontinuity in the field, not a
coordinate artifact.
\end{remark}

\begin{proposition}[The discontinuity carries divergent gradient
  energy under refinement]
\label{P16:prop:divergent_energy}
The local Faddeev--Niemi gradient energy density
$g^2=\sum_a|\partial_a n|^2$, evaluated near a crossing midline on a
sequence of grids of spacing $h$ via central differences, has a peak
value that grows as $h\to0$, consistent with $g^2_{\max}\propto h^{-2}$
(the scaling expected for a true jump discontinuity sampled by a
centred finite difference), rather than saturating at a finite value.
\end{proposition}

\begin{proof}[Evidence]
A direct numerical comparison (script \texttt{crossing\_resolution\allowbreak\_scan.py},
available with this paper's companion files) evaluates $g^2$ on a
sequence of local grids of half-width $1.5\,r_0$ centred on crossing
region~1, for $N_{\mathrm{grid}}\in\{41,61,81,121,161,241\}$ points
per axis. The peak density grows from $64.6$ at $N=41$
($h=0.0656$) to $1964$ at $N=241$ ($h=0.0109$); a log-log fit gives
$g^2_{\max}\propto h^{-1.91}$, consistent with the $h^{-2}$ scaling
of a genuine jump discontinuity under a centred difference, and
inconsistent with the bounded scaling expected of a smooth but
sharply-varying feature. The total local gradient energy in the
sampled region grows correspondingly, from $98.3$ to $195.0$ over the
same resolution range, rather than converging to a finite value.
\end{proof}

\begin{remark}[A naive smoothing does not resolve the discontinuity:
  a topological obstruction]
\label{P16:rem:naive_smoothing_fails}
A natural candidate repair is to replace the single-nearest-point
phase with a continuous blend of the two nearby strands' phases,
weighted by inverse-square distance and averaged on the unit circle
(i.e.\ via $\arg\sum_i w_i\,e^{i\theta_i}$ rather than a direct
average of $t_{\mathrm{near}}$). This reduces the local gradient
energy substantially --- by a factor of $\approx3.1$ at the crossing
tested, stable across the resolution range above --- but does not
eliminate the discontinuity. At the exact crossing midpoint, the two
strands receive equal weight, and the two phases being averaged
differ by exactly $\pi$ (Proposition~\ref{P16:prop:phase_jump}): averaging
two equal-weight unit vectors at antipodal points of the circle gives
the zero vector, leaving the blended phase $\arg(0)$ undefined, and
the surrounding neighbourhood inherits a residual discontinuous snap
as the relative weighting crosses zero. This is not a defect of this
particular smoothing scheme: it is a generic obstruction for any
construction that assigns a single $U(1)$-valued phase as a function
of two nearby source phases differing by exactly $\pi$. Since the
$\pi$ separation is itself forced (Proposition~\ref{P16:prop:phase_jump}),
no smooth, single-phase repair of this kind can remove the
singularity; only an ansatz with genuinely independent phase
freedom for the two strands --- i.e.\ a field valued in the full
$S^3$ Hopf-fibre bundle rather than collapsed to a single real
$\theta(\mathbf{x})$ before the energy functional is evaluated ---
can avoid forcing the two strands to agree on a single phase value at
the point where they are equally relevant. We do not yet have such an
ansatz; identifying one, and determining whether the additional
phase freedom is genuinely needed (as opposed to consistently
fixable by a different gauge choice) is left as the natural next step,
and is recorded as a candidate refinement to the saddle-finder
remedies of Section~\ref{P16:sec:saddle_prerequisite}.6.
\end{remark}

\begin{remark}[Relationship to the inter-tube overlap mechanism]
\label{P16:rem:two_mechanisms}
This phase discontinuity is logically independent of the inter-tube
field-overlap mechanism Paper~XV attributes the saddle-finder's
non-convergence to (tube radius $2.6\times$ the minimum inter-tube
separation at the tested initialisation $C^*=1.5$): the overlap
argument concerns the radial profile $f_0(\rho)$ and does not
reference the internal phase $\theta$ at all, while the discontinuity
identified here is purely a property of $\theta$'s assignment and is
present regardless of $C^*$. The two effects are not shown to be
mutually exclusive and may both contribute additively to the observed
non-convergence; this paper does not establish that either fully
accounts for it, only that the phase discontinuity is a second, real,
independently verifiable candidate cause that should be controlled
for before attributing the full non-convergence to inter-tube overlap
alone.
\end{remark}

\subsection{A candidate repair: a genuine $S^3$ (complex-doublet)
  superposition resolves the discontinuity}
\label{P16:sec:s3_repair}

Remark~\ref{P16:rem:naive_smoothing_fails} shows that no construction
which first collapses each strand's local field data to a single
real $S^2$ phase $\theta$, and only then averages or blends those
phases, can avoid the degeneracy at the crossing midline: two
real phases differing by exactly $\pi$ average to an undefined
direction at equal weight. We now show that working instead in the
full $S^3$ pre-image of the Hopf map --- i.e.\ combining each
strand's own complex doublet \emph{before} projecting to $n\in S^2$
--- removes the obstruction entirely, at least locally at a single
crossing.

\begin{construction}[Complex-doublet superposition ansatz]
\label{P16:constr:s3lift}
Recall the Hopf doublet parametrisation $n=Z^\dagger\boldsymbol\sigma Z$
for $Z=(z_1,z_2)\in\mathbb{C}^2$, $|Z|=1$, with
$Z(f,\theta)=(\cos(f/2),\,\sin(f/2)\,e^{i\theta})$ reproducing the
familiar $n=(\sin f\cos\theta,\sin f\sin\theta,\cos f)$. Near a
crossing, build each strand's own doublet independently, using that
strand's own distance and a fixed strand phase:
\begin{equation}
  Z_{\mathrm{over}}(\mathbf{x}) = \bigl(\cos\tfrac{f_{\mathrm{over}}}2,\,
    \sin\tfrac{f_{\mathrm{over}}}2\,e^{i\theta_{\mathrm{over}}}\bigr),
  \qquad
  Z_{\mathrm{under}}(\mathbf{x}) = \bigl(\cos\tfrac{f_{\mathrm{under}}}2,\,
    \sin\tfrac{f_{\mathrm{under}}}2\,e^{i\theta_{\mathrm{under}}}\bigr),
  \label{P16:eq:s3_doublets}
\end{equation}
with $f_{\mathrm{over}}=f_0(d_{\mathrm{over}}(\mathbf{x})\,C_3^*)$,
$d_{\mathrm{over}}$ the Euclidean distance to the over-strand
centreline (and likewise for under), and
$\theta_{\mathrm{over}},\theta_{\mathrm{under}}$ the fixed phases
$3t_1$, $3(t_1+\pi)$ of Proposition~\ref{P16:prop:phase_jump}. Combine by
inverse-square-distance-weighted complex superposition,
\begin{equation}
  Z(\mathbf{x}) \;=\;
  \frac{w_{\mathrm{over}}Z_{\mathrm{over}} + w_{\mathrm{under}}Z_{\mathrm{under}}}
       {\bigl|w_{\mathrm{over}}Z_{\mathrm{over}} + w_{\mathrm{under}}Z_{\mathrm{under}}\bigr|},
  \qquad w_s = d_s(\mathbf{x})^{-2},
  \label{P16:eq:s3_blend}
\end{equation}
and recover $n(\mathbf{x}) = Z(\mathbf{x})^\dagger\boldsymbol\sigma Z(\mathbf{x})$.
\end{construction}

\begin{proposition}[The complex-doublet superposition is well-defined
  and smooth at the crossing midline]
\label{P16:prop:s3_smooth}
At the geometric midpoint of a crossing region (equal weight,
$w_{\mathrm{over}}=w_{\mathrm{under}}$, where $f_{\mathrm{over}}=f_{\mathrm{under}}\equiv f_*$
by the symmetry of Paper~XV's antipodal isometry), the unnormalised
sum in equation~\eqref{P16:eq:s3_blend} is
\begin{equation}
  Z_{\mathrm{over}}+Z_{\mathrm{under}}
  \;=\; \bigl(2\cos\tfrac{f_*}2,\;0\bigr),
  \label{P16:eq:s3_sum}
\end{equation}
a nonzero vector for any $f_*\neq\pi$, normalising to $Z=(1,0)$ and
hence $n=(0,0,1)$ (the Hopf north pole) at the midpoint --- a
well-defined value, in contrast to the undefined $\arg(0)$ of the
real-phase blend (Remark~\ref{P16:rem:naive_smoothing_fails}). The
resulting field is, moreover, differentiable (not merely continuous)
across the midline: a direct numerical check shows the derivative of
each component of $n$ along the line joining the two strand
centrelines agrees to three significant figures on either side of
the midpoint.
\end{proposition}

\begin{proof}
Equation~\eqref{P16:eq:s3_sum} follows directly from
$e^{i\theta_{\mathrm{over}}}+e^{i\theta_{\mathrm{under}}}
=e^{i\theta_{\mathrm{over}}}(1+e^{i\pi})=0$ (the second doublet
component cancels exactly, by the forced $\pi$ phase separation of
Proposition~\ref{P16:prop:phase_jump}), while the first component
$2\cos(f_*/2)$ does not depend on either strand phase and is generic.
The differentiability claim is verified numerically (script
\texttt{crossing\_S3\_lift\_test\_v3.py}): sampling $n_y$ along the
strand-to-strand line at $s=\pm0.025,\pm0.05$ relative to the
midpoint gives a one-sided numerical derivative agreeing to within
$2.5\%$ on either side, with no sign change or discontinuity.
\end{proof}

\begin{theorem}[The $S^3$-lift ansatz removes the divergent gradient
  energy of Proposition~\ref{P16:prop:divergent_energy}]
\label{P16:thm:s3_bounded_energy}
The local Faddeev--Niemi gradient energy density of
Construction~\ref{P16:constr:s3lift}, evaluated near crossing region~1 on
the same sequence of grids as Proposition~\ref{P16:prop:divergent_energy},
converges to a finite peak value as $h\to0$, in contrast to the
single-nearest-point ansatz's divergence.
\end{theorem}

\begin{proof}[Evidence]
The resolution scan of script \texttt{crossing\_S3\_lift\_test\_v3.py}
(extended to $N_{\mathrm{grid}}=321$, $h=0.00819$) gives the peak
density $g^2_{\max}$ for the single-nearest-point ansatz (A) and the
$S^3$-lift ansatz (C):
\begin{center}
\renewcommand{\arraystretch}{1.2}
\begin{tabular}{rrrrr}
\toprule
$N_{\mathrm{grid}}$ & $h$ & $g^2_{\max}$(A) & $g^2_{\max}$(C) & ratio \\
\midrule
$41$  & $0.0656$ & $64.6$   & $47.4$ & $1.36$ \\
$81$  & $0.0328$ & $233.6$  & $51.2$ & $4.56$ \\
$161$ & $0.0164$ & $888.0$  & $52.7$ & $16.8$ \\
$241$ & $0.0109$ & $1964.5$ & $52.9$ & $37.1$ \\
$321$ & $0.0082$ & $3462.9$ & $53.1$ & $65.3$ \\
\bottomrule
\end{tabular}
\end{center}
A log-log fit over this range gives
$g^2_{\max}(\mathrm{A})\propto h^{-1.92}$, consistent with
Proposition~\ref{P16:prop:divergent_energy}'s divergent
($h^{-2}$-type) scaling, while $g^2_{\max}(\mathrm{C})$'s successive
increments shrink geometrically\\
($2.6,\,1.2,\,1.1,\,0.4,\,0.2,\,0.1$
across the six refinement steps from $N=41$ to $N=321$),\\
consistent with convergence to a finite limit near $53.1$--$53.2$, not
divergence. The total local gradient energy shows the same pattern:
$E(\mathrm{A})$ grows slowly with residual box-edge effects
($98.3\to214.7$ over the same range, not a clean power law, but not
converging either) while $E(\mathrm{C})$ is essentially flat
($35.2\to45.1$, converged by $N=121$).
\end{proof}

\begin{remark}[Interpretation: why the complex lift succeeds where the
  real blend fails]
\label{P16:rem:why_s3_works}
The mechanism is specific to working in $\mathbb{C}^2$ rather than on
the circle of phases directly. A weighted real-number average of two
angles separated by exactly $\pi$ is ill-defined at equal weight
because the two unit vectors $e^{i\theta_{\mathrm{over}}}$,
$e^{i\theta_{\mathrm{under}}}$ are themselves antipodal and cancel.
The Hopf doublet, however, carries an \emph{additional} real degree
of freedom --- the relative weight between $z_1=\cos(f/2)$ (which
does not depend on $\theta$ at all) and $z_2=\sin(f/2)e^{i\theta}$
(which does) --- and it is precisely the $z_1$ component, invariant
under the antipodal phase flip, that survives the cancellation and
fixes the combined doublet to a well-defined, nonzero value. This is
a structural feature of the full $S^3$ Hopf bundle that is invisible
once the doublet has already been projected down to a single real
$S^2$ phase, and is consistent with the original motivating
hypothesis (tracking the trefoil's strands through their full $S^3$
lift rather than the base $S^2$ image alone) for this investigation.
This result is local --- it resolves the discontinuity at one
crossing's midline. Global consistency with the trefoil's unit Hopf
charge is checked in Section~\ref{P16:sec:s3lift_charge_zero}, with a
negative result ($\QH=0$, not $3$); testing inside the full energy
functional of the existing solver remains undone and is moot until a
charge-correct version of this construction is found.
\end{remark}

\subsection{Global non-degeneracy of the $S^3$-lift ansatz}
\label{P16:sec:global_nondegen}

Construction~\ref{P16:constr:s3lift} is stated for a single crossing
region. We now extend it to the whole trefoil and check the first of
the two prerequisites flagged above: that the resulting field is
well-defined (the underlying doublet never vanishes) everywhere, not
just at the one crossing tested in
Section~\ref{P16:sec:s3_repair}.

\begin{construction}[Global two-strand $S^3$-lift ansatz]
\label{P16:constr:global_s3lift}
At a point $\mathbf{x}$, let $t_1(\mathbf{x})$ be the curve parameter
of the nearest point on $\Gamma$, and $t_2(\mathbf{x})$ the curve
parameter of the nearest point at curve-parameter distance greater
than $\pi/2$ from $t_1$ (i.e.\ the nearest point on a genuinely
different strand passage, found in practice by querying the $k=80$
nearest curve samples and taking the first with
$|t-t_1|\bmod 2\pi>\pi/2$). Build $Z_1,Z_2$ as in
equation~\eqref{P16:eq:s3_doublets} using $d_1=|\mathbf{x}-\Gamma(t_1)|$,
$d_2=|\mathbf{x}-\Gamma(t_2)|$ and
$\theta_i=3t_i$, and combine via inverse-square weighting as in
equation~\eqref{P16:eq:s3_blend}. Where no second strand is found within
the search radius (the generic case, far from all three crossings),
the construction reduces to the single-doublet ansatz with $w_2=0$.
\end{construction}

\begin{proposition}[The global ansatz reduces correctly away from
  crossings]
\label{P16:prop:reduces_away}
At a generic point on the trefoil tube, away from all three crossing
regions, Construction~\ref{P16:constr:global_s3lift} agrees with the
original single-nearest-point ansatz to machine precision.
\end{proposition}

\begin{proof}[Evidence]
Direct evaluation at a representative point
($t=1.0$, offset $0.3r_0$ off the centreline, comparable to but
distinct from the three crossing parameters $\pi/6$, $5\pi/6$,
$3\pi/2$ and their antipodes) gives agreement between the two
ansätze to $6\times10^{-14}$ (script
\texttt{global\_s3\_ansatz\_charge\_test.py}), consistent with
floating-point round-off rather than a genuine discrepancy: the
second strand's weight $w_2\propto d_2^{-2}$ is negligible there
since $d_2=O(R_0)\gg d_1=O(r_0)$.
\end{proof}

\begin{proposition}[Global non-degeneracy]
\label{P16:prop:global_nondegen}
The unnormalised doublet $Z_1\!+\!Z_2$ (weighted as in
Construction~\ref{P16:constr:global_s3lift}) does not vanish at any
sampled point in a neighbourhood of the trefoil.
\end{proposition}

\begin{proof}[Evidence]
Two independent numerical scans (script
\texttt{global\_s3\_ansatz\_charge\_test.py}): (i) a broad scan of
$\sim6.4\times10^5$ grid points within $3.5\,r_0$ of the trefoil
centreline, covering the full tube and all three crossing regions,
finds $|Z_1+Z_2|_{\min}=0.107$; (ii) a dense, focused scan of
$60^3$ points in a box of half-width $1.2\,r_0$ centred on each of
the three crossing midpoints individually finds
$|Z_1+Z_2|_{\min}\in\{0.386,0.386,0.465\}$ respectively (the small
asymmetry at crossing~3 is attributable to the discrete curve
sampling, not a physical difference, the three crossings being
related by the exact $\mathbb{Z}_3$ symmetry of $T_{2,3}$ established
in Paper~XV). No zero of $Z_1+Z_2$ was found in either scan.
\end{proof}

\begin{remark}[What non-degeneracy establishes, and what it does not]
\label{P16:rem:nondegen_scope}
Propositions~\ref{P16:prop:reduces_away}--\ref{P16:prop:global_nondegen}
together show that Construction~\ref{P16:constr:global_s3lift} is a
well-defined, continuous map from (a neighbourhood of) the trefoil
into $S^2$, reducing to the proven $\QH=3$ ansatz away from the
crossings. This is the necessary prerequisite for asking what its
Hopf charge is, but it is not by itself a computation of that charge:
a continuous deformation of a map can in principle change the degree
if the deformation is not itself realised as a genuine homotopy of
non-singular maps throughout, and the comparison here is to the
\emph{original ansatz only away from the crossings} (where the two
constructions agree exactly), not via an explicit homotopy connecting
the two field configurations as wholes --- the original
single-nearest-point ansatz is itself discontinuous at the crossings
(Proposition~\ref{P16:prop:phase_jump}), so no continuous homotopy can
have it as an endpoint, and the charge of
Construction~\ref{P16:constr:global_s3lift} must therefore be computed
directly rather than inferred by continuity from the old ansatz's
known value $\QH=3$.
\end{remark}

\subsection{Direct computation: the $S^3$-lift ansatz has $\QH=0$}
\label{P16:sec:s3lift_charge_zero}

The direct computation flagged in Remark~\ref{P16:rem:nondegen_scope} has
been carried out, via the Whitehead (Berry--Chern--Simons) integral
$\QH=(1/4\pi^2)\int A\cdot(\nabla\times A)\,d^3x$ with
$A_i=\Im(Z^\dagger\partial_iZ)$, and gives an unambiguous negative
result: $\QH=0$ for Construction~\ref{P16:constr:global_s3lift}, not
$\QH=3$.

\begin{proposition}[The global $S^3$-lift ansatz is Hopf-trivial]
\label{P16:prop:s3lift_charge_zero}
The Whitehead integral evaluated on a $40^3$--$80^3$ grid converges
to $\QH\approx0$ for Construction~\ref{P16:constr:global_s3lift}, while
the same integral, evaluated identically on the standard
inverse-stereographic $\QH=1$ reference doublet
($Z_1=(1-r^2+2iz)/(1+r^2)$, $Z_2=2(x+iy)/(1+r^2)$), converges
monotonically toward $1$ over the same resolution range, confirming
the method and grid are adequate to resolve a genuine unit charge.
\end{proposition}

\begin{proof}[Evidence]
Script \texttt{whitehead\_hopf\_charge.py} (companion file) gives, at
three resolutions:
\begin{center}
\renewcommand{\arraystretch}{1.2}
\begin{tabular}{rrr}
\toprule
$N_{\mathrm{grid}}$ & $\QH$ ($\QH=1$ reference) & $\QH$ ($S^3$-lift ansatz) \\
\midrule
$40$ & $0.887$ & $-0.0003$ \\
$56$ & $0.939$ & $-0.0002$ \\
$80$ & $0.967$ & $-0.0001$ \\
\bottomrule
\end{tabular}
\end{center}
The reference column converges monotonically toward the known value
$1$ as the grid is refined (the residual deficit at each $N$ is the
expected finite-resolution underestimate of a smooth, compactly
localised integrand). The $S^3$-lift column converges monotonically
toward $0$, not toward $\pm3$ or any other nonzero integer: the
result strengthens, rather than weakens, with resolution, ruling out
a coarse-grid artifact as the explanation.
\end{proof}

\begin{remark}[Diagnosis: the missing poloidal winding]
\label{P16:rem:missing_poloidal_winding}
The vanishing follows from a genuine structural gap in
Construction~\ref{P16:constr:global_s3lift}, traceable back to the
single-nearest-point ansatz it is built from. The thin-torus profile
integral used throughout this series (Paper~I, Paper~XIV),
$J_{2a}\propto\int\sin^4\!f\,[(f')^2+\sin^2\!f/\rho^2]\,\rho\,d\rho$,
contains an angular term $\sin^2\!f/\rho^2$ that is the gradient cost
of a phase $\Phi$ winding once around the tube's poloidal angle
$\chi$ at fixed $\rho$ --- the standard axisymmetric construction is
$\Phi=\chi+N\phi$ for a charge-$N$ Hopfion, with the $\chi$-winding
supplying the base unit of Hopf linking and the toroidal winding $N$
multiplying it. The phase actually implemented by \texttt{build\_ic}
(and reproduced exactly in
Construction~\ref{P16:constr:global_s3lift})
is $\theta=3t$ \emph{only}: a numerical check confirms the assigned
field value is independent of $\chi$ at fixed $(\rho,t)$, with no
poloidal winding term at all. Tracing the level sets of this field
analytically confirms the consequence directly: the preimage of a
generic regular value consists of three small, mutually disjoint,
unlinked transverse circles (one at each of three isolated curve
parameters satisfying $3t\equiv\theta_{\mathrm{target}}\pmod{2\pi}$,
with the poloidal angle $\chi$ free and not constrained at all),
rather than a single connected curve threading the trefoil's
crossings with nonzero linking number. This is consistent with, and
gives an independent analytic confirmation of, the numerical
$\QH=0$ result above. The discrepancy is not a flaw specific to the
$S^3$-lift repair: it is inherited from the underlying
single-nearest-point ansatz, which was evidently never intended to
carry the full topology exactly at the level of the numerical
solver, only to approximate the energetics of the analytically
understood $\Phi=\chi+3t$ profile.
\end{remark}

\begin{remark}[Consequence for the saddle-finder programme]
\label{P16:rem:s3lift_not_sufficient}
The smoothness result of Section~\ref{P16:sec:s3_repair}
(Theorem~\ref{P16:thm:s3_bounded_energy}) stands: the complex-doublet
superposition genuinely removes the divergent gradient energy at a
crossing, and remains useful evidence that $\pi$-phase
discontinuities of this kind are removable in principle by working
in the full $S^3$ pre-image rather than a single real $S^2$ phase.
What Proposition~\ref{P16:prop:s3lift_charge_zero} shows is that
Construction~\ref{P16:constr:global_s3lift}, \emph{as stated}, is not a
valid replacement initial condition for the saddle-finder problem of
Section~\ref{P16:sec:saddle_prerequisite}, since it does not carry the
required $\QH=3$ topology: the construction blends two copies of a
phase that is itself missing the poloidal winding, and blending
cannot supply a winding that neither input possesses. A corrected
construction would need to build in the poloidal dependence
explicitly, e.g.\ a phase of the form $\Phi=\chi+3t$ (or an
equivalent two-strand generalisation with each strand contributing
its own poloidal winding) before the crossing-region superposition is
applied, and to verify directly --- by the same Whitehead-integral
test as above --- that the corrected construction restores $\QH=3$
while retaining the bounded-energy property of
Theorem~\ref{P16:thm:s3_bounded_energy}. This is left for future work,
and is recorded as the corrected next step in
Section~\ref{P16:sec:open}.
\end{remark}

\subsection{A geometric threshold in the profile parameter $C^*$, and
  why blending strands cannot repair the missing winding}
\label{P16:sec:cstar_threshold}

The diagnosis of Remark~\ref{P16:rem:missing_poloidal_winding} --- that
the poloidal winding is structurally absent from
Construction~\ref{P16:constr:global_s3lift} regardless of how the
crossing-region blending is performed --- can be sharpened and given
an independent, purely geometric confirmation by examining how the
construction behaves as the profile parameter $C^*$ is varied away
from its physical value $C_3^*=2.5062$.

\begin{proposition}[Critical profile parameter from the trefoil's
  injectivity radius]
\label{P16:prop:cstar_crit}
The trefoil curve $\Gamma=T_{2,3}$ has injectivity radius exactly
$r_0$: for any point within distance $r_0$ of $\Gamma$, the nearest
point on $\Gamma$ is unique except on a measure-zero set (the three
crossing midline surfaces of Proposition~\ref{P16:prop:phase_jump}).
Beyond distance $r_0$, the nearest-point map is multi-valued on a
substantially larger set. Since the profile $f_0(\rho)=2\arctan(\rho^{-C^*}\!)$
has characteristic tube radius $1/C^*$, this defines a critical
value
\begin{equation}
  C^*_{\mathrm{crit}} = \frac{1}{r_0} \approx 1.1442,
  \label{P16:eq:cstar_crit}
\end{equation}
below which the field's physically relevant support (where
$\sin^2\!f_0$ is non-negligible) extends past the injectivity radius
of the curve it is built from.
\end{proposition}

\begin{proof}[Evidence]
A direct numerical search (script \texttt{cstar\_sweep\_energy.py})
for the global minimum, over all non-adjacent pairs of points on
$\Gamma$ (excluding a curve-parameter neighbourhood of width $\pi/2$
to avoid trivially nearby points on the same local arc), of the
distance between them, gives exactly $2r_0=1.748$, attained at the
three crossing midpoints already characterised in
Proposition~\ref{P16:prop:phase_jump} --- confirming that the crossings
are the curve's own points of closest self-approach, and that the
injectivity radius is $r_0$ exactly. The physical value
$C_3^*=2.5062$ gives tube radius $1/C_3^*=0.399$, a factor of $2.2$
below $r_0$, so the three-crossing picture studied throughout this
section is safely within the regime where it applies.
\end{proof}

\begin{remark}[The energetic signature of crossing $C^*_{\mathrm{crit}}$]
\label{P16:rem:cstar_energy_signature}
At fixed grid resolution, the total gradient energy of the
single-nearest-point ansatz (Construction giving rise to
Proposition~\ref{P16:prop:phase_jump}), swept over $C^*$ from $2.5062$
down to $0.1$, shows no special feature exactly at
$C^*_{\mathrm{crit}}$ itself, but does show its steepest rate of
increase in the neighbourhood $C^*\in[0.7,1.0]$, immediately below
$C^*_{\mathrm{crit}}$, before saturating for $C^*\lesssim0.3$
(script \texttt{cstar\_sweep\_energy.py}). This is consistent with
the qualitative behaviour visible in renderings of the construction
at varying $C^*$: three locally distinct, sharply concentrated
crossing regions for $C^*>C^*_{\mathrm{crit}}$, merging into a single
diffuse, near-isotropic region for $C^*\ll C^*_{\mathrm{crit}}$.
\end{remark}

A natural question raised by this threshold is whether the global
$S^3$-lift ansatz's Hopf charge (Section~\ref{P16:sec:s3lift_charge_zero})
changes character as $C^*$ crosses $C^*_{\mathrm{crit}}$ --- in
particular, whether it is the natural setting for the transient
$\QH=3\to\QH=2$ mechanism considered elsewhere in this series. The
following shows that it is not, for an identifiable and instructive
reason.

\begin{proposition}[The $S^3$-lift's nonzero charge at small $C^*$ is
  a strand-counting artifact]
\label{P16:prop:cstar_artifact}
The Whitehead-integral charge of Construction~\ref{P16:constr:global_s3lift}
(extended to vary with $C^*$), evaluated across
$C^*\in[0.1,2.5062]$, is non-monotonic: it is negligible for
$C^*\gtrsim2$, grows in magnitude through $C^*_{\mathrm{crit}}$, peaks
near $C^*\approx0.3$--$0.4$ at $\QH\approx-0.13$ to $-0.15$, and
decays back toward zero as $C^*\to0$. This non-zero, non-integer
signal is attributable to the construction's two-strand truncation
rather than to any genuine topological effect: replacing the
two-strand superposition with an otherwise identical three-strand
superposition (each point's field built from the three nearest
\emph{distinct} trefoil lobes, weighted identically) reduces the peak
magnitude by roughly half, to $\QH\approx-0.05$ to $-0.06$ over the
same $C^*$ range, with a flatter and broader peak.
\end{proposition}

\begin{proof}[Evidence]
Script \texttt{cstar\_sweep\_charge.py} (two-strand) and
\texttt{cstar\_sweep\_3strand.py} (three-strand) give, at matched
$C^*$ values:
\begin{center}
\renewcommand{\arraystretch}{1.2}
\begin{tabular}{rrr}
\toprule
$C^*$ & $\QH$ (2-strand) & $\QH$ (3-strand) \\
\midrule
$1.5000$ & $-0.0225$ & $-0.0108$ \\
$1.1442$ & $-0.0552$ & $-0.0246$ \\
$0.8000$ & $-0.0957$ & $-0.0429$ \\
$0.6000$ & $-0.1182$ & $-0.0521$ \\
$0.4000$ & $-0.1325$ & $-0.0558$ \\
$0.3000$ & $-0.1285$\footnotemark & $-0.0542$ \\
$0.2000$ & $-0.1346$ & $-0.0495$ \\
\bottomrule
\end{tabular}
\end{center}
\footnotetext{Reproduced from the finer-resolution rerun; the coarser
  initial sweep gave $-0.1259$ at this $C^*$, consistent within the
  expected grid sensitivity.}
A direct check (script \texttt{cstar\_sweep\_charge.py}, querying the
field weight $\sin^2\!f_0$ of the curve's three lobes at the crossing
region) confirms the mechanism: at the physical $C_3^*=2.5062$, the
third-nearest lobe's field weight is negligible relative to the
second's (ratio $\approx0.13$), justifying the two-strand truncation
there, but this ratio grows monotonically as $C^*$ decreases, crossing
unity (the third lobe becoming as important as the second) between
$C^*=0.4$ and $C^*=0.3$ --- the same range where the two-strand
charge anomaly peaks.
\end{proof}

\begin{remark}[Consequence: blending cannot supply the missing
  winding, at any strand count]
\label{P16:rem:blending_insufficient}
Proposition~\ref{P16:prop:cstar_artifact} sharpens
Remark~\ref{P16:rem:missing_poloidal_winding}'s diagnosis into a stronger
claim. It is not merely that the two-strand blend of
Construction~\ref{P16:constr:global_s3lift} happens to omit a relevant
third strand at small $C^*$; it is that \emph{no amount of correctly
weighted strand-blending can repair the deficiency}, because every
strand contributes a phase $\theta_i=3t_i$ that is individually
poloidal-winding-free. Summing or superposing more such terms changes
the Berry curvature's integral quantitatively (as the two-strand
versus three-strand comparison shows) but supplies no new angular
degree of freedom: the construction remains built entirely from phases
that do not depend on $\chi$, and a sum of $\chi$-independent
functions is $\chi$-independent. The shrinking, not vanishing,
residual when more strands are added is consistent with this reading:
each additional correctly-included strand removes one further source
of the same truncation error, but the underlying defect --- absence
of $\chi$ in $\theta_i$ --- is untouched by strand count. A
corrected construction must instead supply $\chi$-dependence
\emph{within each strand's own phase}, e.g.\ $\Phi_i=\chi+3t_i$,
before any crossing-region superposition is applied; this is taken up
directly in Section~\ref{P16:sec:perstrand_winding}.
\end{remark}

\subsection{The corrected construction: per-strand poloidal winding
  restores $\QH=3$}
\label{P16:sec:perstrand_winding}

Remark~\ref{P16:rem:blending_insufficient} identifies the precise repair
needed: each strand's phase must carry its own poloidal dependence,
$\Phi_i=\chi_i+3t_i$, before any crossing-region superposition is
applied. Constructing $\chi_i$ consistently requires a smooth,
globally-defined poloidal angle for the trefoil curve, which in turn
requires confronting a geometric subtlety not previously encountered
in this series.

\subsubsection{A closed, $\mathbb{Z}_3$-symmetric framing}

\begin{proposition}[The trefoil's Bishop frame has nontrivial holonomy]
\label{P16:prop:bishop_holonomy}
Let $(N_1(t),N_2(t))$ be the frame obtained by parallel-transporting
an initial vector perpendicular to $\Gamma'(t)$ along $T_{2,3}$ (the
Bishop, or relatively parallel adapted, frame: at each step $N_1$ is
rotated by exactly the rotation taking $T(t)$ to $T(t+dt)$, with no
additional twist about the tangent). This frame does not close up
after one traversal: $N_1(2\pi)\neq N_1(0)$, with holonomy angle
\begin{equation}
  H = -63.0191^\circ = -1.09989~\mathrm{rad}.
  \label{P16:eq:bishop_holonomy}
\end{equation}
The holonomy is exactly $\mathbb{Z}_3$-symmetric: the rotation
relating $N_1(t_0+2\pi/3)$ to the curve's own rigid rotation of
$N_1(t_0)$ (the extra twist beyond the rigid motion already present
in $T_{2,3}$'s proven $\mathbb{Z}_3$ symmetry) is $H/3=-21.0064^\circ$
at every $t_0$ tested, summing consistently to $H$ over the three
segments.
\end{proposition}

\begin{proof}[Evidence]
Numerical parallel transport (script \texttt{bishop\_frame.py}),
converged to four decimal places by $N_T=5000$ curve samples and
unchanged through $N_T=80000$, gives $H=-63.0191^\circ$. The
per-segment decomposition (comparing $N_1(t_0+2\pi/3)$ to
$N_1(t_0)$ rotated by the curve's own $-120^\circ$ rigid rotation
about the $z$-axis, the same rotation relating $\Gamma(t_0+2\pi/3)$ to
$\Gamma(t_0)$) gives $-21.0064^\circ$, and $3\times(-21.0064^\circ)=
-63.0192^\circ$ agrees with $H$ to within $0.0001^\circ$.
\end{proof}

\begin{remark}[Distinction from the existing $\theta=3t$ winding]
\label{P16:rem:holonomy_distinct}
This holonomy is a property of the curve's normal bundle alone,
unrelated to the coordinate winding number $3$ already present in
$\theta=3t$ (an algebraic fact about the parametrisation, present even
for a hypothetical curve with trivial framing holonomy). The two
numbers are not expected to agree and do not need to: $H\approx
63^\circ$ measures genuine parallel-transport holonomy, while the
ansatz's separate $3\!\times\!2\pi$ total coordinate winding is
already accounted for in $\theta_i=3t_i$.
\end{remark}

\begin{construction}[Holonomy-compensated, arc-length-uniform frame]
\label{P16:constr:compensated_frame}
Define the corrected frame by adding a compensating rotation about
$T(t)$, distributed \emph{proportionally to arc length} rather than
to the parameter $t$:
\begin{equation}
  \varphi_{\mathrm{comp}}(t) = -H\,\frac{s(t)}{L_3},
  \qquad
  s(t)=\int_0^t|\Gamma'(u)|\,du,
  \label{P16:eq:arclength_comp}
\end{equation}
where $L_3$ is the trefoil's total arc length, and rotate
\begin{equation}
(N_1,N_2)\mapsto(\cos\varphi_{\mathrm{comp}}N_1+\sin\varphi_{\mathrm{comp}}N_2,\,
-\sin\varphi_{\mathrm{comp}}N_1+\cos\varphi_{\mathrm{comp}}N_2). 
\end{equation}
Since
$\varphi_{\mathrm{comp}}(0)=0$ and $\varphi_{\mathrm{comp}}(L_3)=-H$,
the corrected frame closes exactly. The arc-length (rather than
parameter) weighting is the natural physical choice, since it
distributes the correction at a constant rate per unit physical
distance rather than per unit of an arbitrary parametrisation.
\end{construction}

\begin{proposition}[Closure]
\label{P16:prop:frame_closure}
The frame of Construction~\ref{P16:constr:compensated_frame} closes to
within $3\times10^{-6}$ degrees (script \texttt{bishop\_frame\_v2.py}),
consistent with exact closure up to floating-point accumulation over
the numerical integration.
\end{proposition}

\subsubsection{The per-strand ansatz and its local smoothness}

\begin{construction}[Per-strand poloidal phase]
\label{P16:constr:perstrand}
For a point $\mathbf{x}$ near strand $i$ (curve parameter $t_i$,
nearest point $\Gamma(t_i)$), define the poloidal angle
$\chi_i(\mathbf{x})$ as the angle of $\mathbf{x}-\Gamma(t_i)$ in the
$(N_1(t_i),N_2(t_i))$ plane of the compensated frame
(Construction~\ref{P16:constr:compensated_frame}), and set
\begin{equation}
  \Phi_i(\mathbf{x}) = \chi_i(\mathbf{x}) + 3t_i,
  \qquad
  Z_i(\mathbf{x}) = \bigl(\cos\tfrac{f_i}2,\,\sin\tfrac{f_i}2\,e^{i\Phi_i(\mathbf{x})}\bigr).
  \label{P16:eq:perstrand_doublet}
\end{equation}
At a crossing, combine $Z_1,Z_2$ (the two nearest distinct strands,
found via the robust three-lobe partition search of
Section~\ref{P16:sec:global_nondegen}, which avoids a $k$-nearest-neighbour
search radius failure identified and corrected during this
investigation) by the same inverse-square-weighted complex
superposition as Construction~\ref{P16:constr:s3lift}.
\end{construction}

\begin{proposition}[$\chi$ is a genuine poloidal angle]
\label{P16:prop:chi_winds}
$\chi_i(\mathbf{x})$ winds exactly once around any small transverse
loop encircling the curve at fixed $t_i$.
\end{proposition}

\begin{proof}[Evidence]
Direct evaluation along a transverse circle of radius $0.3r_0$ at a
generic $t_i$ gives winding number $1.0000$ (script
\texttt{perstrand\_winding\_test.py}), confirming the construction
supplies the poloidal degree of freedom absent from $\theta_i=3t_i$
alone.
\end{proof}

\begin{proposition}[Construction~\ref{P16:constr:perstrand} is smooth at
  all three crossings]
\label{P16:prop:perstrand_smooth}
The field built from Construction~\ref{P16:constr:perstrand} is
continuous and differentiable across the midline of all three
crossing regions, with identical (to within discrete-sampling
precision) behaviour at each, consistent with the proven
$\mathbb{Z}_3$ symmetry of $T_{2,3}$.
\end{proposition}

\begin{proof}[Evidence]
A scan of $n(\mathbf{x})$ across each crossing's midline (script
\texttt{perstrand\_s3lift\_v2.py}) shows all three components passing
smoothly through the midpoint with matching left- and right-side
derivatives (e.g.\ $-0.3056$ on both sides for $n_x$ at crossing~1),
and the values at all three crossings agree to five decimal places at
matched offsets. A resolution scan of the local gradient energy
(margin-trimmed core region, $N_{\mathrm{grid}}\in\{41,61,81,121\}$)
gives $g^2_{\max}\in\{56.2,58.4,59.3,60.0\}$, converging cleanly to a
finite value comparable to the chi-free construction's $\approx53$
(Theorem~\ref{P16:thm:s3_bounded_energy}), confirming the additional
poloidal structure does not reintroduce a divergence.
\end{proof}

\subsubsection{Global Hopf charge: confirmed $\QH=3$}

\begin{theorem}[The corrected per-strand construction restores $\QH=3$]
\label{P16:thm:perstrand_charge_three}
The Whitehead-integral Hopf charge of the global per-strand
construction converges to $\QH=3$, ruling out $\QH=0,1,2$.
\end{theorem}

\begin{proof}[Evidence]
Script \texttt{whitehead\_perstrand\_charge.py}, run at increasing
resolution (validated throughout against the same inverse-stereographic
$\QH=1$ reference doublet used in Proposition~\ref{P16:prop:s3lift_charge_zero}):
\begin{center}
\renewcommand{\arraystretch}{1.2}
\begin{tabular}{rrr}
\toprule
$N_{\mathrm{grid}}$ & $\QH$ (reference, $\to1$) & $\QH$ (per-strand construction) \\
\midrule
$50$  & $0.925$ & $1.251$ \\
$56$  & $0.939$ & $1.268$ \\
$64$  & $0.952$ & $1.052$ \\
$70$  & $0.959$ & $1.176$ \\
$120$ & $0.983$ & $1.796$ \\
$240$ & $0.992$ & $2.499$ \\
\bottomrule
\end{tabular}
\end{center}
The reference column converges monotonically toward $1$, as expected.
The construction's own column is non-monotonic at the lowest
resolutions tested ($N=50,56$, where it briefly exceeds its
$N=64$ value, a finite-size transient not unusual for a more elaborate
integrand than the reference case) but is cleanly monotonically
increasing for $N\geq64$: $1.052\to1.176\to1.796\to2.499$. Since this
value already exceeds $2$ at $N=240$, and the same qualitative
convergence pattern as the validated reference case approaches its
limit from below, $\QH=1$ and $\QH=2$ are excluded as the limiting
value: a monotonically increasing sequence cannot converge to a value
it has already surpassed. $\QH=3$ is the unique integer consistent
with the data, and the deficit $3-\QH(N)$ decreases cleanly with $N$
($1.948,\,1.824,\,1.204,\,0.502$ at $N=64,70,120,240$ respectively),
fitting a power law $3-\QH(N)\propto N^{-1.02}$, consistent with the
expected $O(1/N)$ global accuracy of a central-difference curl
integral.
\end{proof}

\begin{remark}[Status of the construction]
\label{P16:rem:perstrand_summary}
Construction~\ref{P16:constr:perstrand}, combined with
Construction~\ref{P16:constr:compensated_frame}, is therefore the first
ansatz in this investigation to satisfy all three requirements
simultaneously: smoothness at every crossing
(Proposition~\ref{P16:prop:perstrand_smooth}), bounded gradient energy
under refinement (matching Theorem~\ref{P16:thm:s3_bounded_energy}'s
chi-free result in order of magnitude), and the correct topological
charge $\QH=3$ (Theorem~\ref{P16:thm:perstrand_charge_three}). This
resolves the gap left open by Remark~\ref{P16:rem:s3lift_not_sufficient}
and confirms the diagnosis of
Remark~\ref{P16:rem:missing_poloidal_winding}: the poloidal winding,
supplied per-strand rather than recovered by blending chi-blind
phases (Remark~\ref{P16:rem:blending_insufficient}), was indeed the
missing structural ingredient. Two items were identified as
remaining before this construction could be proposed as a replacement
initial condition for Paper~XV's saddle-finder: testing it inside the
actual solver's full energy functional, and confirming that the
corrected construction's local energy reduction
(Theorem~\ref{P16:thm:s3_bounded_energy}) is not offset by additional cost
elsewhere in the field. Both are addressed directly in
Section~\ref{P16:sec:solver_functional_test}.
\end{remark}

\subsection{Testing against the solver's actual energy functional}
\label{P16:sec:solver_functional_test}

The diagnostic checks of Section~\ref{P16:sec:perstrand_winding} use a
local gradient-energy proxy near a single crossing. The solver
referenced throughout this series (\texttt{qh3\_trefoil\_solver\_3d\_v11.py})
minimises a different, more specific functional,
$E_{\mathrm{geom}}=K\!\cdot\!J_4$ with $K=J_{2a}+\mu J_{2\mathrm{iso}}$
the (weighted) gradient energy and $J_4=\int\bigl(F_{xy}^2+F_{xz}^2+F_{yz}^2\bigr)d^3x$
the quartic Faddeev--Niemi term built from the field's own Jacobian
structure, $\mu=3-\varphi$. This is the functional the corrected
construction must actually be tested against.

\begin{proposition}[Global functional comparison at the solver's
  working resolution]
\label{P16:prop:global_functional_comparison}
At the solver's default grid ($N=64$, $h=0.26$, $C^*=2.4987$), the
three constructions compare as follows:
\begin{center}
\renewcommand{\arraystretch}{1.2}
\begin{tabular}{lrrrrr}
\toprule
Construction & $K$ & $J_4$ & $\bar r$ & $E_{\mathrm{geom}}=K\!\cdot\!J_4$ \\
\midrule
A (original, discontinuous)     & $2911.2$ & $126.6$  & $2.980$ & $368{,}438$ \\
B (chi-free $S^3$-lift, $\QH=0$) & $1510.4$ & $76.8$   & $2.956$ & $115{,}960$ \\
C (per-strand, $\QH=3$)          & $1736.4$ & $2181.3$ & $3.191$ & $3{,}787{,}624$ \\
\bottomrule
\end{tabular}
\end{center}
At this resolution, Construction~C's $J_4$ is roughly $17$ times
larger than Construction~B's, and its $E_{\mathrm{geom}}$ is the
largest of the three.
\end{proposition}

This headline comparison, taken alone, would read as a serious
problem for the corrected construction. A direct investigation traces
it to grid resolution rather than to the construction itself.

\begin{proposition}[The large $J_4$ is a resolution artifact, not a
  divergence]
\label{P16:prop:j4_resolution_artifact}
The $J_4$ excess of Construction~C is concentrated almost entirely
within one tube radius of the curve ($\rho<0.5$): bucketing the
solver-grid $\rho_{J_4}$ density by distance from the nearest curve
point, $98\%$ of Construction~C's total $J_4$ comes from $\rho<1.0$,
versus a comparable concentration but $\sim30\times$ smaller magnitude
for Construction~B in the same region. Under local grid refinement
near a single crossing, Construction~C's $J_4$ converges cleanly to a
finite plateau ($453\to632\to729\to749\to729$ across
$N_{\mathrm{grid}}=41,61,81,121,161$, with peak density similarly
flattening), in sharp contrast to Construction~A, whose $J_4$ under
the same refinement grows without bound (peak density
$87\to202\to352\to843\to1569$, consistent with the proven gradient
energy divergence of Theorem~\ref{P16:thm:s3_bounded_energy}). A direct
check at a generic, non-crossing point on the tube reproduces the same
per-unit-length $J_4$ excess found at crossings ($135.3$ versus $124.3$
at $h=0.26$ in matched local boxes), and tiling this generic
contribution over the trefoil's full arc length
($L_3\approx41.26$) predicts a global total of $2128.8$, matching the
directly computed global value of $2181.3$ to within $2.4\%$.
\end{proposition}

\begin{proof}[Evidence]
Scripts\\
\texttt{compute\_global\_comparison.py},\\
\texttt{j4\_core\_refinement\_check.py}, and\\
\texttt{j4\_construction\_A\_refinement.py}.
\end{proof}

\begin{remark}[Interpretation: a generic, well-understood resolution
  requirement]
\label{P16:rem:j4_resolution_interpretation}
The poloidal angle $\chi_i$ satisfies $|\nabla\chi_i|\sim1/\rho$ near
the curve, the same coordinate behaviour as any angular coordinate
near its own axis. This is present at every point along the tube, not
only at crossings, so the excess is uniform per unit arc length rather
than crossing-localised --- consistent with Proposition~\ref{P16:prop:j4_resolution_artifact}'s
generic-point check. The solver's default grid spacing ($h=0.26$)
gives fewer than two points across the tube diameter
($2/C^*\approx0.8$), under-resolving this winding everywhere along the
curve; this is the standard resolution requirement for any field with
genuine angular winding around a thin core (vortex lines, cosmic
strings, monopole cores), not a defect specific to this construction.
By contrast, Construction~A's divergence under refinement
(Proposition~\ref{P16:prop:j4_resolution_artifact}) is a genuine
discontinuity, not a resolvable coarse-grid effect: refining the grid
makes Construction~A's local energy worse, while it makes
Construction~C's local energy converge. This addresses the second item
of Remark~\ref{P16:rem:perstrand_summary}: the corrected construction's
local energy reduction near crossings is not offset by a hidden
divergence elsewhere in the field; the larger total seen at coarse
resolution is a finite, resolvable, and well-understood discretisation
effect, fully accounted for by the generic tube-winding mechanism
above.
\end{remark}

\begin{remark}[Practical recommendation for solver use]
\label{P16:rem:solver_practical_recommendation}
Using Construction~C as an actual solver initial
condition is not yet demonstrated, since adequately resolving the tube
core at the local length scale established in
Proposition~\ref{P16:prop:perstrand_smooth} ($h\lesssim0.03$) is not
feasible on a uniform global grid at the box sizes used in this series
($N^3\sim5\times10^8$ points required, versus $N^3=2.6\times10^5$ at
the solver's current default). A non-uniform or adaptively refined
mesh, concentrated within one to two tube radii of $T_{2,3}$, is the
natural next step; accepting the large initial $J_4$ at coarse
resolution and relying on the solver's own annealing to relax it is
not recommended, given this series' documented history of the
relaxation dynamics itself producing unwanted degeneracies (dilution
and spiking, Sections referenced in the solver's own development
notes) when given a poorly-conditioned starting point. This addresses
the first item of Remark~\ref{P16:rem:perstrand_summary}: the construction
has now been tested against the solver's actual functional, with the
result fully diagnosed, but a genuine solver run remains future work
pending the grid-resolution question above.
\end{remark}

%══════════════════════════════════════════════════════════════════════════════
\section{Gradient-flow minimisation: landscape of $E_{\mathrm{geom}}$ in
  the $\QH=3$ sector}
\label{P16:sec:gradient_flow}
%══════════════════════════════════════════════════════════════════════════════

With Construction~\ref{P16:constr:perstrand} providing a verified
$\QH=3$ initial field, it is natural to ask what happens when that
field is driven toward a stationary point of $E_{\mathrm{geom}}=K\!\cdot\!J_4$
by gradient flow.  This section records the results of direct numerical
experiments on a $144^3$ grid ($h=0.075$,
$10.6$ grid points across the tube diameter at $C^*=2.5062$), and draws
two conclusions: the landscape has a clean Phase~1 / Phase~2 structure
that reveals where the field's excess energy lives; and the
$\QH=3$ topological floor, real in the continuum, is not yet
numerically enforced at this grid resolution.

\subsection{Phase 1 --- gradient-energy cleanup}
\label{P16:sec:gf_phase1}

The initial $E_{\mathrm{geom}}$ of Construction~\ref{P16:constr:perstrand} at
$N=144$ is $1.647\times10^7$ ($K=2756.9$, $J_4=5972.6$,
$\bar r=3.198$).  Running Adam gradient flow with learning rate
$\mathrm{lr}=3\times10^{-4}$ (without any topology constraint), the
field undergoes two visually distinct phases.

\emph{Phase~1} (steps 1--710): $K$ and $J_4$ both decrease together,
$\bar r$ decreases slightly from $3.198$ toward $3.180$, and
clamped\% $=0$.  The field is shedding the excess gradient energy that
the per-strand construction distributes broadly around the tube through
the under-resolved poloidal winding (Proposition~\ref{P16:prop:j4_resolution_artifact}).
$K$ reaches its minimum of $K_{\min}=2313.9$ at step $\approx711$.
The $J_4$-weighted mean radius $\bar r$ reaches its minimum $3.180$
simultaneously and the gradient norm $|\nabla E|$ approaches a
local plateau at $\approx6.7\times10^4$.

\begin{proposition}[Phase~1 energy removal]
\label{P16:prop:gf_phase1}
Starting from Construction~\ref{P16:constr:perstrand} at $N=144$,
$h=0.075$, gradient flow with $\mathrm{lr}=3\times10^{-4}$ reduces
$E_{\mathrm{geom}}$ by a factor of $\approx2$ in $710$ steps, with
$\bar r$ remaining within $0.018$ of $R_0=3$ throughout.  The
$K$-minimum at step $\approx711$ marks the point at which the initial
excess gradient energy has been removed and the field enters the
nontrivial $K$--$J_4$ competition region.
\end{proposition}

\subsection{Phase 2 --- the $K$--$J_4$ competition regime and topology loss}
\label{P16:sec:gf_phase2}

After step $\approx711$, $K$ begins to rise while $J_4$ continues
to fall.  $E_{\mathrm{geom}}=K\!\cdot\!J_4$ keeps decreasing, but now
through a trade-off: reducing the quartic structure ($J_4$) costs
gradient energy ($K$).  Continued running with the tighter
Phase~2 schedule ($\mathrm{lr}=10^{-4}$, per-point rotation clamp at
$9^\circ$) shows:

\begin{itemize}
\item The default $9^\circ$ angle clamp fires \emph{zero times} in $3220$
Phase~2 steps.  At $\mathrm{lr}=10^{-4}$, Adam's adaptive per-point
effective rotation is approximately $10^{-4}$~rad ($0.006^\circ$) per
step---more than $1500\times$ smaller than the clamp threshold.
Topology loss is therefore \emph{not} caused by large single-step
jumps; it is a cumulative drift over many individually-safe steps.

\item $J_4$ decays exponentially from $3237$ to $293$ over $3220$
steps, with no sign of approaching a finite floor.  $J_4/K$ falls
from $1.40$ to $0.099$.

\item $\bar r$ drifts from $3.189$ to $3.215$ ($+0.026$ over 3220
steps)---dramatically more controlled than the unconstrained
predecessor ($\bar r\to4.4$ in $\approx1000$ steps), but still
monotonically outward.

\item The gradient norm $|\nabla E|$ halves from $6.8\times10^4$ to
$2.9\times10^4$, converging steadily but without any sign of
approaching zero.
\end{itemize}

\begin{proposition}[Topology loss via cumulative drift]
\label{P16:prop:gf_topology_loss}
On the $N=144$, $h=0.075$ grid, gradient flow from
Construction~\ref{P16:constr:perstrand} does not converge to a
$\QH=3$ energy minimum.  The field instead drifts continuously toward
the vacuum state ($J_4\to0$, $\QH\to0$, $E_{\mathrm{geom}}\to0$) through
a process of slow, cumulative small-angle rotations, each individually
within the discrete-topology-safe regime.  The per-point rotation clamp
(Option~A, Section~\ref{P16:sec:gradient_flow}) prevents catastrophic
single-step topology changes but cannot prevent this cumulative
mechanism.
\end{proposition}

\subsection{Interpretation: discrete topological floor}
\label{P16:sec:gf_floor}

In the continuum, $E_{\mathrm{geom}}$ in the $\QH=3$ sector is bounded
below by a positive constant proportional to $|\QH|^{3/4}$
(Faddeev--Niemi bound).  This floor prevents gradient flow from
reaching the vacuum.  On a \emph{discrete} grid, however, the
numerical approximation of $E_{\mathrm{geom}}$ can approach zero even
when the \emph{continuum} field carries $\QH=3$ winding, provided the
grid is too coarse to resolve that winding faithfully.  The tube radius
$1/C^*\approx0.399$ at the target $C^*=2.5062$ requires approximately
$15$ grid points across the tube diameter for the local refinement scans
of Proposition~\ref{P16:prop:j4_resolution_artifact} to reach their
converged values; the $N=144$ grid provides only $10.6$.

\begin{remark}[Two-phase landscape, $\QH$ transition, $K_{\min}$ convergence, and preliminary cascade observations]
\label{P16:rem:gf_twophase}

\textbf{Resolution-convergence of $K_{\min}$.}
The $N=144$, $h=0.075$ run gives $K_{\min}=2313.9$ and the $N=192$, $h=0.05$
run gives $K_{\min}\approx2294$. The shift of $-19.9$ ($1.0\%$) is consistent
with an $O(h^2)$ grid error of $\approx(0.05)^2/(0.075)^2\approx0.44$ times the
$N=144$ error, as expected for a second-order discretisation.

Phase~1 (steps $1$--$710$ for the $N=144$ run, steps $1$--$840$ for $N=192$)
genuinely operates within the $\QH=3$ sector: the field sharpens without
topology change, $\bar r$ drifts \emph{toward} $R_0=3$, and the gradient norm
approaches a natural plateau. The $N=192$, $h=0.05$ gradient-flow run
(Phase~1 exit -at step $\approx840$--$850$, $K \approx2306$, and
$K_{\min}\approx2294$ at step $\sim1290$) extends the $N=144$ and $N=192$
results in two directions: it sharpens the saddle estimate, and it reveals
a qualitative change in Phase~2 behaviour that provides direct evidence for
the topological picture of Section~\ref{P16:sec:gf_floor}.

\textbf{Phase~2 behaviour and $\mathbb{Z}_3$-symmetry breaking ($N=192$ run).}
Visual inspection of the $\rho_{J_4}$ field at step $n=4000$ at the end
of the $N=192$ run ($J_4\approx322$), (Figure~\ref{P16:fig:tre_foil_density}) shows
the final solver state of the $\mathbb{Z}_3$-symmetry that was spontaneously
broken during Phase~2 at ($J_4\approx702$). In the image, one outer lobe
together with its adjacent midpoint are visibly depleted while the two
remaining lobes show energy accumulation on either side of the depleted
region.  This $(2{+}1)$ lobe-splitting pattern arises because the gradient
flow on the flat-bottomed $\QH=3$ energy landscape (Remark~7.6)%\ref{P18:rem:saddle_difficulty}
of Paper~XVIII~\cite{Paper18}) has no restoring force once floating-point
accumulation in Adam's moment estimates ($\sim10^{-7}$ per step) kicks the
field off the symmetric locus. The particular direction selected varies across
seeds; all three equivalent $(2{+}1)$ directions are equally accessible.

\textbf{Physical interpretation.}
The $(2+1)$ splitting is the lightest single string-breaking mode of
the baryon: one colour flux tube is released while the remaining
two-colour diquark system stays bound.  This is consistent with the
fraying-site geometry (Section~\ref{P16:sec:arc_segment}), in which a
single midpoint site can deplete before its two $\mathbb{Z}_3$ images,
giving an asymmetric $(2+1)$ configuration that decays further rather
than relaxing back to the symmetric vacuum.

A higher-resolution run ($N=256$, $h=0.0375$) is the natural next
step to check whether $K_{\min}$ has converged further and whether
all three $(2+1)$ breaking directions are energetically equivalent
under deliberate seeding.

\textbf{Preliminary observations from the $N=256$, $h=0.0375$ run.}
A higher-resolution run ($N=256$, $h=0.0375$, $\delta_{\max}=2^\circ$,
Phase~2 $\mathrm{lr}=2.5\times10^{-5}$, at phase transition: K rises
$> 1.0$ above its minimum, reported fully in Paper~XVIII~\cite{Paper18})
gives $K_{\min}\approx2268$ and reveals two qualitative improvements
over the $N=192$ results.

First, the tighter angle clamp ($2^\circ$ versus $9^\circ$) suppresses
the accumulated numerical asymmetry: the fraying remains approximately
$\mathbb{Z}_3$-symmetric throughout Phase~2, confirming that the $(2{+}1)$
breaking in the $N=192$ run was a numerical artefact of the larger
per-step rotations rather than a physical instability.

Second, the fraying site ordering is directly observable.  At step~50
(immediately after Phase~1 exit), the three outer distal lobes
($R_{xy}=R_0+r_0=3.874$, $z=0$) show clear energy depletion at their
outward faces, while the crossing sites show depletion on their outer
(vacuum-facing) side of the tube.  Midpoint disturbance is minor and
appears later.  By step~1000, crossing depletion has propagated from
the outer to the inner side of the tube, and by step~2000 the energy
density at the lobe breaks appears to reflect inward rather than exit
outward --- consistent with the topological barrier preventing vacuum
escape within the gradient-flow dynamics.  The midpoints remain the
last and least disturbed category throughout.

The empirical fraying ordering --- distal lobes first, crossing depletion
second, midpoints last --- is discussed in the context of the revised
cascade picture in Paper~XVIII~\cite{Paper18},
Remark~7.5.%\ref{P18:rem:cascade_ordering}
\end{remark}

%══════════════════════════════════════════════════════════════════════════════
\section{Arc-segment energy partition of the trefoil}
\label{P16:sec:arc_segment}
%══════════════════════════════════════════════════════════════════════════════

Visual inspection of the energy-density field $\rho_{J_4}$ from the
gradient-flow runs reveals a structurally non-uniform distribution
along the trefoil tube with a 12-site fraying pattern.  Three geometrically
distinct categories of field disturbance are visible, each internally
related by the $\mathbb{Z}_3$ symmetry of the knot:
\begin{figure}[h]
    \centering
        \begin{subfigure}{0.45\linewidth}
        \centering
        \includegraphics[width=\linewidth]{trefoil_energy1.png}
        \caption{Top-down ($z$-axis) view. The $\mathbb{Z}_3$ symmetry
        of the trefoil is clearly visible; the three inner loops and
        three outer lobes are resolved, and the crossing structure fraying
        is apparent at this projection angle.}
        \label{P16:fig:trefoil_1}
    \end{subfigure}
    \hfill
    \begin{subfigure}{0.45\linewidth}
        \centering
        \includegraphics[width=\linewidth]{trefoil_energy2.png}
        \caption{Oblique view. The three-lobed knotted tube is visible
        with its characteristic writhe; multiple site fraying is apparent
        at this projection angle.\\}
        \label{P16:fig:trefoil_2}
    \end{subfigure}
    \caption{$\QH=3$ condensate field (\texttt{np\_final.py}, $n=4000$)
    visualised by \texttt{hopfion\_viewer.html} using $1.5\times10^5$
    particles. Grid: N=192, h=0.05, C*=2.5062, Phase 1: lr=0.0003, no
    angle clamp (cleanup), Phase 2: lr=0.0001, delta\_max=$9.0^\circ$ per
    point, Phase transition: K rises $> 2.0$ above its minimum.
    \emph{Hue} encodes the direction of the unit field $\mathbf{n}$
    at each point (colour = $\mathbf{n}$ direction per the viewer
    legend); \emph{brightness} encodes the local field density $\rho$,
    with white/yellow indicating high $\rho$ (topological flux
    concentration in the A-segment tube bodies) and dark blue
    indicating low $\rho$ (flux depletion at the crossing vertices,
    consistent with the $\rho_{J_4}$--$S_{\mathrm{eff}}$
    anti-correlation of Remark~\ref{P16:rem:rhoJ4_anticorrelation}).
    The three crossing vertices are visibly depleted at $z=+r_0$,
    elevated above the equatorial plane.}
    \label{P16:fig:tre_foil_density}
\end{figure}
\begin{enumerate}[label=(\roman*)]
\item \textbf{Crossing approach/departure} (6 sites, 2 per colour crossing,
  120° apart within each sub-category): fraying begins on approach to
  each crossing, reaches maximum intensity at the crossing vertex where
  $S_{\mathrm{eff}}$ is highest and $\rho_{J_4}$ is lowest, then
  diminishes symmetrically on departure.  The two sub-categories
  (approach and departure sides) are offset by 60° azimuthally.
\item \textbf{Midpoint torsion twist} (3 sites, 120° apart): a minor
  fraying event at each colour-mixing midpoint, where the poloidal
  Frenet frame undergoes its maximum rate of rotation.
\item \textbf{Outer-lobe distal} (3 sites, 120° apart): a moderate
  fraying event at the most radially distant point of each lobe petal,
  equidistant in 3D from the two adjacent crossing vertices.  This site
  corresponds to no extremum of $S_{\mathrm{eff}}$ and does not coincide
  with any arc-segment boundary.
\end{enumerate}

The three boundary types carry distinct geometric data, which can be
quantified with the integrated $J_4$ asymmetry by partitioning the tube
into six arc-segment integrals and identifying which geometric features
control the energy distribution.

\subsection{Geometry of the six arc-segments}
\label{P16:sec:arcseg_geometry}

The trefoil $T_{2,3}$ has six structurally distinct arcs, alternating
between three \emph{crossings} and three \emph{midpoints}.

\begin{definition}[Arc-segment boundaries]
\label{P16:def:arc_boundaries}
Define the six distinguished parameter values on $T_{2,3}$ in
monotone order:
\begin{align}
  t_{\mathrm{C}_0} &= \tfrac{\pi}{6},\quad
  t_{\mathrm{M}_0} = \tfrac{\pi}{3},\quad
  t_{\mathrm{C}_1} = \tfrac{\pi}{6}+\tfrac{2\pi}{3},\quad
  t_{\mathrm{M}_1} = \tfrac{\pi}{3}+\tfrac{2\pi}{3},\nonumber\\
  t_{\mathrm{C}_2} &= \tfrac{\pi}{6}+\tfrac{4\pi}{3},\quad
  t_{\mathrm{M}_2} = \tfrac{\pi}{3}+\tfrac{4\pi}{3}.
\end{align}
The \textbf{B-segments} (crossing to midpoint) are the three arcs
$[\mathrm{C}_n, \mathrm{M}_n]$, each of arc-parameter width
$\tfrac{\pi}{6}$ (30°).  The \textbf{A-segments} (midpoint to
crossing) are the three arcs $[\mathrm{M}_n, \mathrm{C}_{n+1}]$,
each of width $\tfrac{\pi}{2}$ (90°).
\end{definition}

\subsection{Numerical results}
\label{P16:sec:arcseg_numerics}

We compute $J_4^{(s)} = \int_{\text{seg }s} \rho_{J_4}\,\mathrm{d}^3x$
by evaluating Construction~\ref{P16:constr:perstrand} analytically on uniform grids, assigning
each grid point to its nearest arc-segment via a KDTree query on the
sampled curve $\Gamma(t)$, and summing $\rho_{J_4}\,h^3$ within each
segment.  Two resolutions are used to assess grid artifacts.

\begin{table}[h]
\centering
\begin{tabular}{lcccc}
\toprule
 & \multicolumn{2}{c}{$N=64$, $h=0.175$} & \multicolumn{2}{c}{$N=96$, $h=0.113$}\\
\cmidrule(lr){2-3}\cmidrule(lr){4-5}
Segment & $J_4^{(s)}$ & \% & $J_4^{(s)}$ & \% \\
\midrule
$B_0$ (C$_0\to$M$_0$, 30°) & 263.6 & 7.03 & 362.5 & 7.01 \\
$A_0$ (M$_0\to$C$_1$, 90°) & 989.4 & 26.40 & 1362.5 & 26.35 \\
$B_1$ (C$_1\to$M$_1$, 30°) & 262.3 & 7.00 & 363.4 & 7.03 \\
$A_1$ (M$_1\to$C$_2$, 90°) & 974.8 & 26.01 & 1351.7 & 26.15 \\
$B_2$ (C$_2\to$M$_2$, 30°) & 276.9 & 7.39 & 374.0 & 7.23 \\
$A_2$ (M$_2\to$C$_0$, 90°) & 980.7 & 26.17 & 1356.0 & 26.23 \\
\midrule
$J_4^{\text{total}}$ & 3747.6 & 100 & 5170.1 & 100 \\
\bottomrule
\end{tabular}
\caption{Arc-segment $J_4$ integrals at two grid resolutions.
  B-segments (crossing to midpoint, 30°) and A-segments (midpoint to
  crossing, 90°) each form a $\mathbb{Z}_3$-symmetric triplet to within
  2.5\% at $N=64$ and 1.4\% at $N=96$.  The A-segment dominance
  (79\% of total $J_4$ in 75\% of arc length) reflects the
  $\rho_{J_4}$--$S_{\mathrm{eff}}$ anti-correlation: flux concentrates
  in the bodies of the tube where curvature is intermediate, not at
  the high-curvature crossing vertices.}
\label{P16:tab:arc_segments}
\end{table}

\begin{proposition}[$\mathbb{Z}_3$ symmetry of arc-segment energies]
\label{P16:prop:z3_symmetry}
For Construction~\ref{P16:constr:perstrand} on a cubic grid at both resolutions tested, the
three B-segment integrals agree to within $2.5\%$ ($N=64$) and $1.4\%$
($N=96$), and the three A-segment integrals agree to within $0.6\%$
and $0.3\%$ respectively.  The residual asymmetry decreases with
resolution and is consistent with a grid discretisation artifact.
The field therefore partitions its $J_4$ charge in a $\mathbb{Z}_3$-symmetric
pattern, with each arc carrying the same energy as its two $\mathbb{Z}_3$
images to numerical precision.
\end{proposition}

\subsection{Energy density asymmetry}
\label{P16:sec:arcseg_asymmetry}

The raw segment integrals $J_4^{A} \approx 3 \times J_4^{B}$ are
partially explained by the 3:1 arc-length ratio (A-segments span 90°,
B-segments span 30°).  The physically meaningful quantity is the
energy \emph{density} per unit arc-parameter:

\begin{equation}
  \rho_A = \frac{J_4^A}{\pi/2} \approx 864,\qquad
  \rho_B = \frac{J_4^B}{\pi/6} \approx 700 \quad (N=96).
\end{equation}

The density ratio $\rho_B/\rho_A \approx 0.81$ is consistent across
both resolutions (0.818 at $N=64$, 0.811 at $N=96$), indicating that
energy density is genuinely higher on the A-segments (midpoint-to-crossing,
approaching each crossing) than on the B-segments (crossing-to-midpoint,
departing each crossing).  The crossing vertex itself sits at the A/B
boundary and carries the local minimum of $\rho_{J_4}$ 
Remark~\ref{P16:rem:rhoJ4_anticorrelation}), so the integrated A-segment
dominance reflects flux accumulated across the entire 90° approach arc,
not at the vertex.

With the arc-segment data above, and (script \texttt{fraying\_site\allowbreak\_geometry.py})
we record the exact or numerically confirmed results below.

\begin{proposition}[Geometric data of the three distinguished point types and
arc-segments]
\label{P16:prop:geom_data}
For the trefoil $T_{2,3}$ with $R_0=3$, $r_0=0.874$, $n=0,1,2$:
\begin{enumerate}[label=(\roman*)]
\item \textbf{Crossings} 
($t=t_{\mathrm{C}_n} = \tfrac{\pi}{6}+\tfrac{2\pi n}{3}$):
  curvature $\kappa=0.399$ (global maximum along the tube), torsion
  $\tau\approx-0.17$, $S_{\mathrm{eff}}\propto\kappa^2$ at its maximum,
  colour charge $n$ maximally pure (one weight dominates), $\rho_{J_4}$ at
  a local minimum.  Position: $R_{xy}=R_0=3$, $z=+r_0=+0.874$ for all
  three, 120° apart in azimuth; the tangent vector lies \emph{exactly} in the
  equatorial plane ($\dot\Gamma_z=3r_0\cos(\pi/2)=0$), so any energy
  redirected at the crossing exits at the adjacent B-segment departure
  zone and propagates radially outward in the horizontal plane at
  $z=+r_0$.
\item \textbf{Midpoints} ($t=t_{\mathrm{M}_n}=\tfrac{\pi}{3}+\tfrac{2\pi n}{3}$):
  curvature $\kappa=0.026$ (global minimum, factor~15 below crossings),
  torsion $|\tau|\approx11.3$ (factor~67 above crossings), $S_{\mathrm{eff}}$
  at its minimum, colour weights exactly equal ($w_1=w_2=w_3=\tfrac{1}{3}$).
  Position: $R_{xy}=R_0-r_0=2.126$, $z=0$ for all three (equatorial plane).
  The tangent vector has a nonzero $z$-component
  ($\dot\Gamma_z\approx -0.525\,|\dot\Gamma|$ at $t_{\mathrm{M}_k}$),
  so any energy ejected at a midpoint propagates out of the equatorial plane.
  Field disturbance is minor: a small fraying at the point of maximum
  Frenet-frame rotation.
\item \textbf{Outer-lobe distal points} ($t_{\mathrm{D}_n}=\tfrac{2\pi n}{3}$):
  curvature $\kappa=0.349$ (intermediate), torsion is $+0.014$; $\tau\approx0$
  (the lobe sits at a torsion zero by the local $\mathbb{Z}_2$
  symmetry of $T_{2,3}$ at $t=2\pi n/3$),
  $S_{\mathrm{eff}}$ intermediate (neither maximum nor minimum).
  Position: $R_{xy}=R_0+r_0=3.874$ (maximum radial extent of the knot),
  $z=0$ for all three (equatorial plane, 120°apart in azimuth).
  Each distal point is equidistant in three-dimensional space from its
  two nearest crossing vertices:
  $|\Gamma(t_{\mathrm{D}_0})-\Gamma(t_{\mathrm{C}_n})|=3.626$.
  (exact by $\mathbb{Z}_3$ symmetry). The tangent vector has a nonzero
  $z$-component ($\dot\Gamma_z\approx +0.321\,|\dot\Gamma|$), so energy
  ejected here also propagates out of the equatorial plane.  The field
  shows moderate fraying here, not correlated with $S_{\mathrm{eff}}$,
  $\rho_{J_4}$ extrema, or any arc-segment boundary.
\item \textbf{Arc interiors} (A- and B-segment bodies): the bulk of
  $\rho_{J_4}$ resides in the A-segment bodies (midpoint-to-crossing,
  90°), which are smooth, non-fraying, and carry $\approx79\%$ of the
  total $J_4$ charge at energy density $\rho_A\approx864$ per unit
  arc-parameter.  The B-segment bodies (crossing-to-midpoint, 30°)
  carry the remaining $\approx21\%$ at density $\rho_B\approx700$.
  Neither segment body shows visible fraying; all fraying is localised
  to the site categories above.
\end{enumerate}
The three crossings all sit at $z=+r_0$ (above the equatorial plane),
while the three midpoints and the three outer-lobe distal points all
sit at $z=0$.  This is an exact $\mathbb{Z}_2$ asymmetry of the writhe:
the knot's over-strands and under-strands are not mirror-symmetric
under $z\to -z$.  Only the crossing category has a tangent lying exactly
in the equatorial plane; midpoint and distal-lobe tangents have nonzero
$z$-components of opposite sign ($-0.525$ and $+0.321$ respectively).
\end{proposition}

\begin{remark}[Torsion contrast and field blindness at the midpoints]
\label{P16:rem:torsion_contrast}
The torsion ratio $|\tau_{\mathrm{M}}|/|\tau_{\mathrm{C}}|\approx67$
is dramatically large.  The midpoints are the sites of the most violent
Frenet-frame rotation on the entire curve --- yet Paper~XV proves that
torsion makes \emph{zero} contribution to $S_{\mathrm{eff}}$ at all orders
(Theorem~6.1 of Paper~XV~\cite{Paper15}).  The field
is maximally blind to geometry exactly where the geometry is most active. At
the midpoints, the poloidal frame must reverse its writhe direction (the ribbon
twist changes sign), while the action integral $S_{\mathrm{eff}}$ is at a
\emph{local minimum} there.  The minor midpoint fraying visible in 
$\rho_{J_4}$ is therefore caused by the frame reversal, not by any restoring
force from torsion. This combination --- geometric maximum of torsion, field
minimum of $S_{\mathrm{eff}}$, colour mixing complete --- identifies the
midpoints as structurally weak loci: the tube is locally under-compressed
($\kappa_{\min}$), the colour charge is maximally diffuse, and the field
exerts zero restoring force from the torsion direction. The disturbance appears
to propagate from each midpoint \emph{inward} along the A-segment
toward the adjacent crossing, where it is redirected into the
B-segment departure channel; the distal lobe then receives secondary
energy from the B-segment arc.
\end{remark}

\begin{remark}[$\rho_{J_4}$ anti-correlates with $S_{\mathrm{eff}}$ along the tube]
\label{P16:rem:rhoJ4_anticorrelation}
The visual field and the arc-segment integrals together establish an 
anti-correlation: $\rho_{J_4}$ is lowest precisely where
$S_{\mathrm{eff}}\propto\kappa^2$ is highest (at the crossings), and
the bulk of the $J_4$ flux is concentrated in the A-segment bodies
where $\kappa$ is intermediate and declining toward the midpoints.
This is consistent with the Hopf flux being repelled from regions of
high tube curvature: the field is maximally compressed geometrically
at the crossings but carries the least topological flux there.  The
fraying at each crossing is therefore a \emph{deficit} signature ---
a structural weakness where the field cannot maintain its profile under
the geometric constraint --- rather than a concentration.
\end{remark}

\begin{remark}[Physical interpretation: 12-site fraying and decay geometry]
\label{P16:rem:arcseg_physics}
The arc-segment partition, together with the geometric data of
Proposition~\ref{P16:prop:geom_data} and
Remarks~\ref{P16:rem:torsion_contrast}--\ref{P16:rem:rhoJ4_anticorrelation},
supports the following picture, recorded here as a precise
phenomenological description of the numerically observed field.

\textbf{Twelve fraying sites in three categories.}
The $\rho_{J_4}$ field shows 12 distinct sites of field disturbance,
each category internally related by $\mathbb{Z}_3$:
(i)~six crossing-associated sites (approach and departure sides of each
vertex, 2~per crossing, the two sub-categories offset by 60°
azimuthally): fraying intensity maximal at the vertex where $S_{\mathrm{eff}}$
peaks and $\rho_{J_4}$ is depleted;
(ii)~three midpoint torsion-twist sites (minor, at $t_{\mathrm{M}_k}$):
frame reversal with negligible $S_{\mathrm{eff}}$ response;
(iii)~three outer-lobe distal sites (moderate, at $t_{\mathrm{D}_k}$):
located at maximum radial extent $R_{xy}=R_0+r_0$, $z=0$, equidistant
from both adjacent crossing vertices in 3D ($d=3.626$), at an intermediate
$S_{\mathrm{eff}}$ ($\kappa=0.349$), between the crossing maximum $0.399$
and midpoint minimum $0.026$).

The distal lobe sites are geometrically notable: they are not associated
with any extremum of $S_{\mathrm{eff}}$, any arc-segment boundary, or any
torsion event.  Their equidistance from the two adjacent crossings
($d=3.626$ in 3D, exact by $\mathbb{Z}_3$ symmetry) identifies them as
the points on each outer lobe that are maximally remote from the nearest
confinement vertices in physical space, positioned to accept the disturbance
from the torsion-twist sites as it propagates outward, fraying the distal lobe.

\textbf{Energy anti-correlation with $S_{\mathrm{eff}}$.}
The dominant pattern --- $\rho_{J_4}$ depleted where $S_{\mathrm{eff}}$
is highest, concentrated where $S_{\mathrm{eff}}$ is intermediate ---
is consistent with the Hopf flux being repelled from high-curvature
regions.  The A-segment bodies (midpoint-to-crossing, intermediate
$\kappa$, smooth) carry 79\% of the total $J_4$.  The crossing vertices
(maximum $\kappa$) carry the least.

\textbf{Quark--flux-tube identification.}
If the three crossings ($z=+r_0$, $\kappa_{\max}$, pure colour~$k$,
local $\rho_{J_4}$ minimum) are identified as quark colour vertices
and the A-segment bodies as the flux-tube arms connecting them, then
the flux-tube energy is concentrated in the arm interiors and depleted
at the vertices --- consistent with the standard QCD picture in which
string energy dips near the quark endpoints in the centre-of-mass
frame.  The crossing vertex acts as a topological lens: incoming
A-segment flux is redirected at the crossing and exits through the
adjacent B-segment departure zone, which is the primary jet-ejection
site.  The fraying at each crossing is a structural deficit rather
than a concentration, reflecting the field's inability to maintain
its profile at the confinement vertex.

\textbf{Degeneracy and its breaking.}
The $\mathbb{Z}_3$ symmetry of Proposition~\ref{P16:prop:z3_symmetry}
renders all members of each category energetically equivalent.
The only geometric asymmetry available at this level of description is
the $\mathbb{Z}_2$ asymmetry of the writhe: all three crossings sit at
$z=+r_0$, all other special points at $z=0$. Under an external field
distinguishing $z>0$ from $z=0$, this asymmetry could in principle
favour one string-breaking mode over another, selecting where energy
emerges first.  Whether this geometric $\mathbb{Z}_2$ connects to the
physical breaking of generation degeneracy
(Remark~\ref{P16:rem:gf_twophase}) remains an open question.

\textbf{String breaking at the midpoints.}
The torsion contrast (Remark~\ref{P16:rem:torsion_contrast}) identifies
the midpoints as the natural primary string-breaking loci.  When the
$\QH=3$ baryon loses topological integrity at these B-segments, the
disturbance propagates inward along the A-segment toward the crossing,
is redirected at the crossing vertex (topological lens), and exits at
the B-segment departure zone.  The distal lobe receives secondary
energy from the B-segment arc.
\end{remark}

%══════════════════════════════════════════════════════════════════════════════
\section{Discussion: What the Exclusions Tell Us}
\label{P16:sec:open}
%══════════════════════════════════════════════════════════════════════════════

Each of Sections~\ref{P16:sec:vertices}--\ref{P16:sec:colourshift} tests a
distinct hypothesis and finds a distinct, specific reason for its
failure: a category mismatch between combinatorial and dynamical
quantities (Section~\ref{P16:sec:vertices}); a single automorphism
masquerading as two independent structures
(Section~\ref{P16:sec:colourauto}); a conformal embedding that exists only
at one level, foreclosing the obvious lepton-sector analogy
(Section~\ref{P16:sec:level}); a genuinely independent algebraic structure
whose internal shape nonetheless does not match an ordered ladder
(Section~\ref{P16:sec:q8}); and a generation-universal piece of physics
that is real but provably insufficient on its own
(Section~\ref{P16:sec:colourshift}).

Three patterns emerge across all the failures, worth stating
explicitly as guidance for future attempts.

\begin{enumerate}[label=(\Roman*)]
\item \textbf{The lepton sector's dominant hierarchy is a dynamical,
  not a combinatorial, fact.} The integer $n_e=20$ is fixed by a
  transcendental thermodynamic matching condition
  (Paper~XII~\cite{Paper12}), not by any property of $\twoI$'s finite
  group structure. Every hypothesis in
  Sections~\ref{P16:sec:vertices}--\ref{P16:sec:q8} attempted to extract $n_q$
  from $\twoT/\twoI$ group theory directly, and every one failed. This
  suggests the quark tower level, if it exists in analogous form,
  will likewise come from a dynamical matching condition specific to
  the $\QH=3$ sector's own condensate coupling, not from the
  representation theory of $\twoT$ or its embedding in $\twoI$.
\item \textbf{Generation-universal mechanisms are necessary but not
  sufficient.} Both the T-matrix phases of Paper~XV (open
  problem~(iv)(b)) and the colour level shift of Paper~V
  (Section~\ref{P16:sec:colourshift}) are generation-universal,
  proved results that nonetheless leave the bulk of the hierarchy
  unexplained. A successful mechanism must supply
  \emph{generation-dependent} input, and Proposition~\ref{P16:prop:colourshift_fail}
  shows this input is not a smooth function of $g$ in the most natural
  variable ($\vp$-exponent shift) tried here.
\item \textbf{Colour is an output of the $\twoI\to\twoT$ breaking, not
  an input.} Theorem~\ref{P16:thm:colour_exclusion} proves that the colour
  $\mathbb{Z}_3$ automorphism $\sigma$ has no realisation at the
  $\twoI$ level: every $\QH=2$ configuration is a colour singlet, and
  colour charge first becomes a well-defined quantum number only at the
  $\QH=3$ level where $\sigma$ acts on $\twoT$.  In retrospect, all
  six mechanisms of Sections~\ref{P16:sec:vertices}--\ref{P16:sec:e8coxeter}
  assumed colour was already a given and searched for an additional
  structure layered on top of it; the theorem shows this search was
  necessarily circular.  The correct architecture for any quark mass
  formula must therefore begin from $\QH=2$ quantities (specifically,
  lepton masses $m_\ell^{(g)}$) and add $\QH=3$-specific contributions
  as first-order corrections --- exactly the form of
  equation~\eqref{P16:eq:quark_mass_recap}.  What the formula is missing,
  by Proposition~\ref{P16:prop:colourshift_fail}, is a
  \emph{generation-dependent} $\QH=3$ contribution, and by
  Pattern~(I), that contribution is most likely dynamical rather than
  combinatorial.  The string-tension energy $\sigma\cdot3R_0$
  (Section~\ref{P16:sec:stringtension}) is qualitatively the right kind of
  object: it is strictly $\QH=3$-specific (absent for the $\QH=2$
  torus), it is generation-dependent through the Y-junction geometry,
  and it would enter the mass formula as an additive correction to the
  profile energy rather than as a multiplicative factor on the lepton
  mass.
\end{enumerate}

The open problem stated in Paper~XV is accordingly sharpened, not
resolved: the quark tower level $n_q^{(g)}$, for each generation $g$
and each of the up-type/down-type sectors, requires a dynamical
derivation analogous to Paper~XII's thermodynamic matching condition
for $n_e$, evaluated for the $\QH=3$ sector's own condensate coupling
--- a calculation not attempted in this paper, and which the
exclusions above suggest is unlikely to reduce to a representation-
theoretic or combinatorial shortcut.

The $E_8$ Coxeter formula of Section~\ref{P16:sec:e8coxeter}, while not
itself that derivation, encodes two algebraic constraints
(Remark~\ref{P16:rem:nq_structure}) that any successful derivation must
reproduce: the cross-sector identity $2k\,h(E_8)(T_1-T_3) = |\twoT|$
linking the lepton T-phase span to the order of the quark-sector
group, and the specific distribution of $E_8$ Coxeter exponents
between isospin sectors that produces the non-monotone generation
structure of the down-type masses.

%══════════════════════════════════════════════════════════════════════════════
\section{Summary}
\label{P16:sec:summary}
%══════════════════════════════════════════════════════════════════════════════

This paper reports six ruled-out mechanisms for the quark generation
mass hierarchy left open by Paper~XV, each falsified for an
identifiable structural reason rather than a failure to fit data:

\begin{enumerate}[noitemsep]
\item The $96$ $\vp$-valued vertices broken under $\twoI\to\twoT$ do
  not determine the quark tower level $n_q$, under any of three
  precise readings tested (Proposition~\ref{P16:prop:vertices_fail}); the
  underlying geometric decomposition and the uniqueness of the
  $\vp$-free conjugate among five copies of $\twoT$ in $\twoI$ remain
  valid (Proposition~\ref{P16:prop:five_copies}) but do not extend to a
  mass formula.
\item The three $1$-dimensional irreducible representations of
  $\twoT$ do not supply a generation label independent of colour,
  since the same outer automorphism permutes both triples
  simultaneously (Proposition~\ref{P16:prop:bosonic_fail}).
\item Varying the $\SU(3)$ WZW level has no consistent McKay-type
  embedding beyond $k=1$, since the central-charge match underlying
  the entire $\QH=3$ sector is exact only there
  (Proposition~\ref{P16:prop:level_fail}).
\item The quotient $\twoT/Q_8\cong\mathbb{Z}_3$ is a genuine,
  independent algebraic structure (Proposition~\ref{P16:prop:q8_independent}),
  but its three cosets split as one distinguished coset plus a
  conjugate pair, not as three mutually exchangeable slots
  (Proposition~\ref{P16:prop:q8_fail}).
\item Paper~V's proved, generation-universal colour level shift
  $\delta n^{(C)}=\tfrac12$ leaves a residual gap to measured quark
  masses that is not a smooth function of generation index, ruling
  out a simple per-generation correction layered on top of it
  (Proposition~\ref{P16:prop:colourshift_fail}); a separate scale ambiguity
  between the condensate-scale formula and Paper~V's own
  low-energy CKM-angle usage remains unresolved
  (Remark~\ref{P16:rem:scale_ambiguity}).
\item Paper~V's alternative integer-exact ratio formula
  $m_q/m_{\ell(g)}=\exp(n_q\pi/9)$, built on Paper~VI's proved
  incommensurability of $\pi$ and $\ln\vp$, rests on exactly one
  independently derived input --- the strange-quark $E_8$ Coxeter
  coincidence $m_s=13$ --- with the remaining five quark-to-exponent
  assignments fixed only by that input together with the already-known
  mass ordering, via a monotonicity argument with no further
  independent content (Propositions~\ref{P16:prop:e8_unique}--\ref{P16:prop:e8_mixed});
  the resulting predictions miss measured masses by $26\%$ to $70\%$
  for four of the five fitted quarks, and are undefined at tree level
  for the sixth (Proposition~\ref{P16:prop:e8_accuracy}).
\end{enumerate}

A seventh investigation, complementary to mechanisms~(1)--(6) and not
a falsification, establishes the first positive group-theoretic result
of this paper: the colour $\mathbb{Z}_3$ automorphism $\sigma$ does
not extend to any automorphism of $\twoI$, and consequently every
$\QH=2$ configuration is a colour singlet
(Theorem~\ref{P16:thm:colour_exclusion}, Section~\ref{P16:sec:colour_exclusion}).
The Dynkin-diagram formulation is transparent: $\mathrm{Aut}(\hat{E}_8)=\{1\}$,
while the colour $\mathbb{Z}_3$ is the non-trivial graph automorphism
of $\hat{E}_6$.  This reframes mechanisms~(1)--(6) retrospectively:
they all searched for a generation label within colour, but colour
itself only comes into being at the $\twoI\to\twoT$ transition
(Remark~\ref{P16:rem:symmetry_breaking_colour}).  The only logically
consistent formula for $m_q^{(g)}$ must therefore start from the
$\QH=2$ lepton masses and add $\QH=3$-specific corrections ---
exactly the structure of equation~\eqref{P16:eq:quark_mass_recap}.

An eighth investigation, motivated by colour confinement rather than
by group theory or WZW data, examines the trefoil's Y-junction
string-tension energy $\sigma\cdot3R_0$ (Section~\ref{P16:sec:stringtension}).
Unlike mechanisms~(1)--(6), this is not falsified: $\sigma$ is simply
unvalued in Paper~XV (Section~\ref{P16:sec:sigma_unvalued}), and we show
that the chameleon-sector field and the Hopf profile function are the
same field (Section~\ref{P16:sec:field_identification}), removing one
possible obstruction, while a second obstruction --- the Y-junction
arms lie off the trefoil's own tube
(Proposition~\ref{P16:prop:arms_off_curve}) --- means $\sigma$ cannot be
extracted from existing tube integrals. Both obstructions are resolved:
$\sigma\cdot3R_0 = 0.04405\,E_{\mathrm{profile}}^*$ analytically
(Proposition~\ref{P16:prop:sigma_ext_formula}), and
$E_\cup/E_{\mathrm{baryon}}=6.9\times10^{-8}$ numerically
(Proposition~\ref{P16:prop:ejunc_zero}), giving the exact baryon energy
formula $E_{\mathrm{baryon}}=1.04405\,E_{\mathrm{profile}}^*$
(Remark~\ref{P16:rem:ejunc_consequence}).

A ninth investigation revisits the scale ambiguity of
Remark~\ref{P16:rem:scale_ambiguity} directly, by running the colour-shift
predictions of mechanism~(5) from $\LUV$ to the PDG comparison scale
with a proper one-loop, multi-threshold calculation
(Section~\ref{P16:sec:rgrunning}). The strange-quark prediction matches
measurement to $0.09\%$ (Proposition~\ref{P16:prop:strange_match}) ---
independently of, and not reducible to, the $E_8$ Coxeter route's own
exact strange-quark coincidence
(Remark~\ref{P16:rem:two_strange_results}). The other five quarks remain
off by factors of $1.7$ to $49$, but with a systematic structure: every
prediction undershoots, and the deviation grows with generation,
the qualitative signature of a missing generation-dependent,
always-positive energy contribution. This is recorded as quantitative
motivation for Section~\ref{P16:sec:stringtension}'s string-tension
candidate, not as evidence that the candidate is correct.

A tenth investigation asks whether the pole mass's renormalon
ambiguity corresponds to an identifiable geometric feature of the
$\QH=3$ condensate (Section~\ref{P16:sec:saddle_prerequisite}). The
$\QH=2$ torus's own Derrick instability is fully resolved by a
scale-invariant geometric-mean functional, but Paper~XV's open
saddle-finder problem for $\QH=3$ reports a qualitatively more severe
failure: every numerical method tried fails to converge to any
saddle, with the instability traced to self-interaction between the
trefoil's three crossing regions specifically --- a structural,
non-speculative analogue of the renormalon obstruction
(Remark~\ref{P16:rem:saddle_renormalon_parallel}). A second, independent
candidate cause is identified directly in the field ansatz itself: the
existing numerical solver's initial condition has a forced, exact
$\pi$ phase discontinuity at every crossing region
(Proposition~\ref{P16:prop:phase_jump}), carrying gradient energy that
diverges under mesh refinement (Theorem~\ref{P16:thm:s3_bounded_energy}'s
ansatz~A column) rather than the bounded value expected of a genuine
field configuration. A candidate local repair --- combining each
strand's own Hopf doublet by complex superposition before projecting
to $S^2$, rather than blending their projected phases (Construction~B,
chi-free, of Proposition \ref{P16:prop:global_functional_comparison}) ---
removes the divergence and converges to a finite gradient
energy (Theorem~\ref{P16:thm:s3_bounded_energy}, ansatz~C column), and
extends to a globally non-degenerate construction
(Proposition~\ref{P16:prop:global_nondegen}).
However, a direct Whitehead-integral computation of this construction's
Hopf charge gives the negative result $\QH=0$, not $\QH=3$
(Proposition~\ref{P16:prop:s3lift_charge_zero}), traced to a genuine
structural gap shared with the original ansatz: the numerical solver's
phase $\theta=3t$ carries no poloidal ($\chi$) winding, unlike the
$\Phi=\chi+3t$ implicit in the profile integrals used analytically
throughout the series (Remark~\ref{P16:rem:missing_poloidal_winding}).
A geometric threshold in the profile parameter, $C^*_{\mathrm{crit}}=1/r_0$
(Proposition~\ref{P16:prop:cstar_crit}), was investigated as a possible
site of a genuine topological transition, but the resulting
Whitehead-integral signal there was traced to the same two-strand
truncation rather than to new physics (Proposition~\ref{P16:prop:cstar_artifact}),
sharpening the diagnosis into a fully structural statement: blending
any number of chi-blind strands cannot supply the missing winding
(Remark~\ref{P16:rem:blending_insufficient}). Acting on that diagnosis
directly, a corrected construction supplying the poloidal dependence
\emph{within} each strand's own phase, $\Phi_i=\chi_i+3t_i$, built on
a holonomy-closed, arc-length-uniform-compensated parallel frame
(Construction~\ref{P16:constr:compensated_frame}), is smooth at all three
crossings (Proposition~\ref{P16:prop:perstrand_smooth}) and confirmed by
direct Whitehead-integral computation to carry the correct charge,
$\QH=3$ (Theorem~\ref{P16:thm:perstrand_charge_three}). This is the first
construction in the investigation to satisfy smoothness, bounded
energy, and correct topology simultaneously
(Remark~\ref{P16:rem:perstrand_summary}). Tested directly against the
solver's own minimised functional, $E_{\mathrm{geom}}=K\!\cdot\!J_4$
(Proposition~\ref{P16:prop:global_functional_comparison}), the corrected
construction shows a substantially larger $J_4$ at the solver's
working resolution; this is traced precisely to an under-resolved
poloidal winding present uniformly along the tube (not concentrated
at, or specific to, the crossings), confirmed by a refinement scan
showing clean convergence in sharp contrast to the original ansatz's
confirmed divergence under the same refinement
(Proposition~\ref{P16:prop:j4_resolution_artifact}). The conclusion
is that the motivating renormalon question still cannot yet be posed
precisely without a confirmed $\QH=3$ saddle, but the prerequisite
obstruction has narrowed considerably: a construction satisfying every
structural requirement identified so far now exists and has been
tested against the actual functional it would need to minimise, with
the remaining barrier identified concretely as a grid-resolution
requirement (Remark~\ref{P16:rem:solver_practical_recommendation}) rather
than an open conceptual question. A further
complication is noted independently: if the
$\QH=3$ sector is populated only transiently, the relevant object may
be a family of energy-perturbation-dependent configurations rather
than one static saddle (Remark~\ref{P16:rem:transient_complication}).

The recurring lesson (Section~\ref{P16:sec:open}) is that the lepton
sector's analogous hierarchy is fixed dynamically, not
combinatorially, and that the quark sector's tower level likely
requires an equally dynamical derivation specific to the $\QH=3$
condensate coupling --- a genuinely open calculation, not a
shortcut through the finite group theory of $\twoT\subset\twoI$,
nor through fitting an integer label to an already-known mass
ordering. The string-tension energy of Section~\ref{P16:sec:stringtension}
is the most promising candidate for that dynamical derivation
identified in this paper, precisely because it is qualitatively
absent from the lepton sector rather than merely quantitatively
different from it, and Section~\ref{P16:sec:rgrunning}'s residual pattern
is the first quantitative hint that such a term is doing real
explanatory work rather than serving as a post-hoc patch. Both that
candidate and the oscillation question of
Section~\ref{P16:sec:saddle_prerequisite} converge on the same
computational bottleneck: a working $\QH=3$ saddle-finder, currently
the single most consequential open numerical problem identified in
this paper.

%══════════════════════════════════════════════════════════════════════════════
\begin{thebibliography}{99}

\bibitem{Paper1}
F.~Manfredi,
\textit{The Density-Feedback Faddeev--Niemi Hopfion:
  Exact Algebraic Predictions for Standard-Model Parameters},
Zenodo (2026).
\href{https://doi.org/10.5281/zenodo.19342027}{doi:10.5281/zenodo.19342027}

\bibitem{Paper4}
F.~Manfredi,
\textit{The Hopf Spoke, WZW Fermion Mass Renormalization,
  and Three- and Four-Term Formulas for the Fine Structure Constant},
Zenodo (2026).
\href{https://doi.org/10.5281/zenodo.19504729}{doi:10.5281/zenodo.19504729}

\bibitem{Paper5}
F.~Manfredi,
\textit{Colour Confinement, the QCD Scale, and Hadronic Structure
  from the Density-Feedback Hopfion},
Zenodo (2026).
\href{https://doi.org/10.5281/zenodo.19638857}{doi:10.5281/zenodo.19638857}

\bibitem{Paper6}
F.~Manfredi,
\textit{The Golden-Ratio Tower: Fermion Scale Hierarchy from the Hopfion RG Fixed Point},
Zenodo (2026).
\href{https://doi.org/10.5281/zenodo.19646945}{doi:10.5281/zenodo.19646945}

\bibitem{Paper12}
F.~Manfredi,
\textit{The Profile Normalisation Conjecture ---
  Proof via CMB Energy Identity and Two-Route Equivalence},
Zenodo (2026).
\href{https://doi.org/10.5281/zenodo.20210460}{doi:10.5281/zenodo.20210460}

\bibitem{Paper15}
F.~Manfredi,
\textit{The $Q_H=3$ Sector of the Density-Feedback Hopfion},
Zenodo (2026).
\href{https://doi.org/10.5281/zenodo.20691001}{doi:10.5281/zenodo.20691001}

\bibitem{Paper17} F.~Manfredi,
\textit{The $Q_H=1$ Sector of the
Density-Feedback Hopfion: Topology, Group Structure, Colour Exclusion},
Zenodo (2026).
\href{https://doi.org/10.5281/zenodo.21013412}{doi:10.5281/zenodo.21013412}

\bibitem{Paper18} F.~Manfredi,
\textit{Preimage Topology and Confinement: The $Q_H=3$ Sector of the Density-Feedback Hopfion},
Zenodo (2026).
\href{https://doi.org/10.5281/zenodo.21047750}{doi:10.5281/zenodo.21047750}

\bibitem{Zenodo_P16}
F.~Manfredi,
\textit{Code supplement for Paper~XVI: Python scripts and html viewer},
Zenodo (2026).
\href{https://doi.org/10.5281/zenodo.20805458}{doi:10.5281/zenodo.20805458}

\bibitem{CIZ1987}
A.~Cappelli, C.~Itzykson, J.-B.~Zuber,
\textit{The ADE classification of minimal and $A_1^{(1)}$ conformal
invariant theories},
Commun.\ Math.\ Phys.\ \textbf{113}, 1 (1987).
\href{https://doi.org/10.1007/BF01221394}{doi:10.1007/BF01221394}

\bibitem{Nesterenko1996}
Y.~V.~Nesterenko,
\textit{Modular functions and transcendence questions},
Sb.\ Math.\ \textbf{187}, 1319 (1996).
\href{https://doi.org/10.1070/SM1996v187n09ABEH000158}{doi:10.1070/SM1996v187n09ABEH000158}

\bibitem{Beneke1998}
M.~Beneke,
\textit{A quark mass definition adequate for threshold problems},
Phys.\ Lett.\ B \textbf{434}, 115 (1998).
\href{https://doi.org/10.1016/S0370-2693(98)00741-2}{doi:10.1016/S0370-2693(98)00741-2}

\bibitem{Witten1984}
E.~Witten,
\textit{Non-abelian bosonization in two dimensions},
Commun.\ Math.\ Phys.\ \textbf{92}, 455 (1984).
\href{https://doi.org/10.1007/BF01215276}{doi:10.1007/BF01215276}

\bibitem{PDG2024}
R.~L. Workman et al.\ (Particle Data Group),
\textit{Review of Particle Physics},
Prog.\ Theor.\ Exp.\ Phys.\ \textbf{2022}, 083C01 (2022);
updated 2024.
\href{https://doi.org/10.1093/ptep/ptac097}{doi:10.1093/ptep/ptac097}

\bibitem{Hoang1999}
A.~H.~Hoang,
\textit{Top Quark Mass Definition and Top Quark Pair
Production Near Threshold at the NLC},
arXiv.hep-ph/9904468(1999).
\href{https://doi.org/10.48550/arXiv.hep-ph/9904468}{doi:10.48550/arXiv.hep-ph/9904468}

\bibitem{Hoang2002}
A.~H.~Hoang,
\textit{Heavy Quarkonium Dynamics,
in \textit{At the Frontier of Particle Physics}, World Scientific (2001)},
arXiv.hep-ph/0204299
\href{https://doi.org/10.48550/arXiv.hep-ph/0204299}{doi:10.48550/arXiv.hep-ph/0204299}

\bibitem{Rodrigo1997}
G.~Rodrigo,
\textit{Quark Mass Effects in QCD Jets},
Nuclear\ Physics\ B\ -\ Proceedings\ Supplements \textbf{54}, 60-65 (1997).
\href{https://doi.org/10.1016/s0920-5632(97)00015-7}{doi:10.1016/s0920-5632(97)00015-7}

\end{thebibliography}
\end{document}